\documentclass{situationist-preprint}
\usepackage{longtable}
\usepackage{array}
\usepackage{needspace}

\title{A Classification of Gettier Cases}
\msid{EPI-01-b}
\msfamily{EPI-01}
\msdiscipline{Epistemology}
\mstype{Companion manuscript}
\mskeywords{Gettier problem; safety; epistemic luck; veritic luck; false lemmas; fake barns; lotteries; testimony; case models}

\newcommand{\stnsources}{44}
\newcommand{\stnlit}{38}
\newcommand{\stnschem}{6}
\newcommand{\stnvariants}{267}
\newcommand{\stnmodels}{311}
\newcommand{\stnG}{109}
\newcommand{\stnK}{105}
\newcommand{\stnGrows}{10}
\newcommand{\stnagree}{32}
\newcommand{\stndepart}{6}
\newcommand{\stnundisputed}{29}
\newcommand{\stlocc}{58}
\newcommand{\stlocd}{51}
\newcommand{\stevsafe}{10}
\newcommand{\stevunsafe}{9}
\newcommand{\stnjtsrc}{41}
\newcommand{\stnnonrobust}{6}
\newcommand{\stnmulti}{21}
\newcommand{\stnGdeltaone}{10}
\newcommand{\stcnmodels}{42656}
\newcommand{\stcnchecks}{426560}
\newcommand{\stcnconfigs}{1333}
% operations: {'ded': 6, 'fal': 44, 'fl': 6, 'hos': 21, 'lie': 44, 'nfl': 9, 'nor': 16, 'rel': 44, 'rem': 20, 'sec': 16, 'unj': 41}
% classes (all models): {'F': 41, 'G-env': 26, 'G-fl': 7, 'G-int': 13, 'G-lot': 3, 'G-mrg': 4, 'G-nec': 2, 'G-nfl': 7, 'G-tmp': 5, 'G-tst': 42, 'K': 105, 'U': 56}
% departures: ['HAR', 'KAL', 'TMP', 'ZBM', 'HLC', 'ATC']
% nonrobust: ['RED', 'TMP', 'MOR', 'LTN', 'BKL', 'BKH']


\begin{document}
\maketitle

\begin{abstract}
EPI-01-a analyses knowledge as safe, justified belief and proves that a justified true belief fails to be knowledge exactly when it is unsafe, that is, exactly when it is a Gettier case. The cases from which the post-Gettier literature was built are, however, of visibly different kinds: some turn on a false lemma and some do not, some on a deceptive environment and some on a fact that intervenes between the evidence and the truth, some on time, chance, testimony or the limits of discrimination. We give a typology of these cases by the source of their luck, prove for each class a theorem connecting it to safety, and show that the classes partition the unsafe justified true beliefs. The typology is then applied to a catalogue of \stnlit{} cases from the literature and \stnschem{} schematic cases, each encoded as a finite case model, and to \stnvariants{} systematic variants obtained by changing one feature of a case. A program computes, for each of the \stnmodels{} models, whether the belief is justified, true, safe and known, and its class; it checks every theorem on every model. Every case the literature presents as a Gettier case is computed to be a justified true unsafe belief. The six cases in which the computed verdict departs from the source's are listed and discussed; each is one on which the literature is itself divided or whose verdict depends on how its basis is individuated.
\end{abstract}

\section{Introduction}

\subsection{The problem}

\citet{Gettier1963} gave two cases in which a subject has a justified true belief that is not knowledge. In the six decades since, the cases have multiplied, and so have the diagnoses. The first diagnosis was that Gettier's subjects reason from a false belief, and that knowledge requires that one's justification not depend on a falsehood \citep{Clark1963,Harman1973}. Fake barns \citep{Goldman1976} involve no inference at all; the stopped clock of \citet{Russell1948}, which predates Gettier, involves a false background belief but no premise; lotteries involve neither; and \citet{Feldman1974} showed that the inference in Gettier's cases can be rerouted through true intermediate beliefs without restoring knowledge. A second family of diagnoses appealed to defeat \citep{LehrerPaxson1969}, a third to causation \citep{Goldman1967}, a fourth to sensitivity \citep{Nozick1981}. Each fitted a subset of the cases. The impression this history leaves, and which \citet{Shope1983} documented, is that \enquote{Gettier case} names a heterogeneous family whose members share nothing beyond the fact that they defeat the tripartite analysis.

EPI-01 holds the opposite: knowledge is safe, justified true belief, and Gettier cases are exactly those in which a justified true belief is unsafe. EPI-01-a states the analysis and proves the characterization (its Theorem~3.3 and Corollary~3.4), and EPI-01-c reaches the same biconditional from knowledge-first premises. EPI-01-d shows, for Gettier's own two cases, that the false lemma is incidental: what makes them failures of knowledge is that the subject's basis could easily have produced the same belief while it was false. The principal manuscript checks seven standard cases in prose. A thesis of the form \enquote{exactly those} invites a more systematic test. If the Gettier cases are a single phenomenon, their visible differences must be differences in how a belief comes to be unsafe, and a classification by source of luck should be a classification of the ways in which safety fails. If they are not, there should be a class of cases, recognized in the literature, which the classification cannot place.

This manuscript carries out the test. We give a typology of Gettier cases by the source of the luck that makes the belief true; we prove, for each class, a theorem stating how its members fail to be safe and what feature of the case would restore safety; and we show that the classes partition the unsafe justified true beliefs. We then encode the cases of the literature, together with systematically generated variants, as finite case models, and compute every verdict.

\subsection{The typologies in the literature}

Classifications of epistemic luck are older than the safety analysis. \citet{Unger1968} proposed that knowledge excludes its being an accident that the subject is right. \citet{Pritchard2005} distinguished veritic luck, the luck that a belief is true given its basis, from evidential luck, the luck that one has the basis at all, and from content and capacity luck, and argued that only veritic luck is incompatible with knowledge; he later distinguished, within veritic luck, intervening luck, in which something comes between the subject's basis and the fact (Chisholm's sheep), from environmental luck, in which the basis connects with the fact in the normal way but the environment makes error easy (fake barns) \citep{Pritchard2012}. Reflective luck, the luck that, from the subject's own point of view, her belief is true, was separated from veritic luck by \citet{Pritchard2005} and corresponds to the distinction \citet{Sosa2007} draws between animal and reflective knowledge. \citet{Zagzebski1994} showed how to manufacture Gettier cases: take a justified false belief and add a stroke of good luck that makes it true without touching the basis.

These distinctions are sound, and we use them. What they leave open is the connection that EPI-01 asserts: that every kind of luck which defeats knowledge is a way in which safety fails, and that every way in which safety fails is a kind of luck which defeats knowledge. We supply that connection, kind by kind.

\subsection{Plan}

Section~\ref{s:frame} defines case models and recalls the analysis. Section~\ref{s:typ} gives the typology. Section~\ref{s:thm} proves the theorems connecting each class to safety and the partition theorem. Section~\ref{s:cat} describes the catalogue and the mechanical check and discusses particular entries. Section~\ref{s:obj} answers objections, and Section~\ref{s:not} records what the classification does not establish. Appendices~A to~E contain the computed material: the catalogue of cases with their models and verdicts, the systematic variants, the safety profiles and rival verdicts of every model, the incidence tables of the typology, and an exhaustive check of the theorems on all case models up to a bound.

\section{Case models}\label{s:frame}

\subsection{Definitions}

EPI-01-a works with possible worlds, a closeness relation subject only to centring (its Assumption~2.1), bases, and a primitive notion of justification. To compute verdicts we need finite representations of cases. We use the following.

\begin{definition}[Case model]\label{d:model}
A \emph{case model} consists of a subject $S$, a proposition $p$, a basis $m$, a justification value $J \in \{0,1\}$, and a finite set of worlds containing an actual world $w_0$. Each world $v$ carries a distance $d(v) \in \{0,1,2,3,4\}$, with $d(v) = 0$ exactly when $v = w_0$; a belief value $B(v)$, which is $1$ when $S$ believes $p$ on basis $m$ at $v$; a truth value $p(v)$; and a tag recording what differs between $v$ and $w_0$. Optionally a world carries the truth value of a weaker proposition $q$ entailed by $p$, and a value $\mathrm{alt}(v) = 0$ when, at $v$, $S$ believes on basis $m$ some other proposition, which is false there. The model also carries a threshold $\theta \geq 0$, fixed by the context of the case, and descriptive features: the route of the basis, the false lemma it relies on if any, whether the fact that makes $p$ true is causally connected to the basis in the normal way, whether there is a defeater, and whether the subject has reflective access to the reliability of her basis. In every model $B(w_0) = 1$.
\end{definition}

\begin{definition}[Closeness in a model]\label{d:close}
In a model with threshold $\theta$, $v$ is close to $w_0$ when $d(v) \leq \theta$.
\end{definition}

Since $d(w_0) = 0$, Definition~\ref{d:close} satisfies centring for every $\theta$. Distances are ordinal: $1$ marks what could very easily have happened given the case as described, $2$ what could easily have happened, $3$ what could have happened only with some change in the circumstances, and $4$ what is remote. The default threshold is $2$. Raising $\theta$ makes more worlds close, which is the model counterpart of the contextualist thesis of EPI-13 that attending to a possibility makes it close; the bank cases (BKL, BKH) are encoded this way.

\begin{definition}[Safety, knowledge and Gettier cases in a model]\label{d:safe}
The belief is \emph{locally safe} at $\theta$ when $p(v) = 1$ for every close $v$ with $B(v) = 1$; this is Definition~2.2 of EPI-01-a. It is \emph{basis-wide safe} when, in addition, $\mathrm{alt}(v) \neq 0$ for every close $v$; this is the amendment for necessary truths which EPI-01-a accepts in reply to its third objection. The subject \emph{knows} when $J = 1$ and the belief is basis-wide safe (Definition~2.3 of EPI-01-a with the amendment). The model is a \emph{Gettier case} when $J = 1$, $p(w_0) = 1$, and the belief is not basis-wide safe (Definition~2.4 of EPI-01-a with the same amendment).
\end{definition}

The amendment changes nothing for contingent propositions in the catalogue: the value $\mathrm{alt}$ is encoded only in the two schematic cases about arithmetic, SNC and SNS, and in all other models local and basis-wide safety coincide.

\begin{definition}[Error distance]\label{d:delta}
The \emph{error distance} $\delta$ of a model is the least $d(v)$ over the worlds $v$ at which $B(v) = 1$ and $p(v) = 0$, or $\mathrm{alt}(v) = 0$; it is $\infty$ if there is no such world.
\end{definition}

\begin{lemma}\label{l:delta}
A belief in a case model is basis-wide safe at $\theta$ if and only if $\delta > \theta$. In particular, a model with a justified true belief is a Gettier case if and only if $\delta \leq \theta$.
\end{lemma}

\begin{proof}
The belief fails to be basis-wide safe exactly when some close world is an error world, that is, when some error world has distance at most $\theta$, which is to say that the least such distance is at most $\theta$. The second claim follows by Definition~\ref{d:safe}.
\end{proof}

Lemma~\ref{l:delta} reduces every verdict to a single number. Its corollary, which the computations make vivid, is that knowledge comes in a profile: as $\theta$ ranges over $0, \ldots, 4$, a belief is safe at every threshold below $\delta$ and unsafe at every threshold from $\delta$ on. Appendix~C prints this profile for every model as a string of five letters, \texttt{S} for safe and \texttt{u} for unsafe.

\begin{assumption}[Faithful encoding]\label{a:faithful}
For each case of the catalogue, the encoded model contains, at the stated distances, the variations of the actual situation which the case description presents as easy (or not easy), with the belief and truth values that the description fixes for them; and the justification value is the one the source assigns.
\end{assumption}

Assumption~\ref{a:faithful} is where the philosophical judgment enters, and we return to it in Section~\ref{s:obj}. It does not assume the analysis: the encodings record what could easily have happened and whether the subject would then have believed truly, which is what the case descriptions say, and the source's verdict on knowledge is recorded separately and compared with the computed one.

\subsection{Kinds of luck}

\begin{definition}[Luck]\label{d:luck}
In a model with $J = 1$ and $p(w_0) = 1$:
\begin{enumerate}[label=(\alph*)]
\item the belief is \emph{veritically lucky} if there is a close error world;
\item it is \emph{evidentially lucky} if there is a close world at which the subject lacks the basis (tag \texttt{basis});
\item it is \emph{reflectively lucky} if the subject has no reflective access to the reliability of her basis;
\item it is \emph{content-lucky} if the truth of $p$ is itself a matter of chance, the basis bearing on the outcome after the fact.
\end{enumerate}
\end{definition}

These are the distinctions of \citet{Pritchard2005}, restated for finite models. Clause~(a) is the luck the Gettier literature is about; clauses (b) to (d) are the kinds which, on Pritchard's view and on the analysis of EPI-01-a (its section on the taxonomy of luck), are compatible with knowledge.

\section{The typology}\label{s:typ}

A Gettier case is unsafe because of some close error world. The typology classifies Gettier cases by what differs in the nearest error worlds, that is, by the tag of the worlds at distance $\delta$, and, among cases in which the truth-maker varies, by the route of the basis. The tags are as follows. \texttt{tm}: the fact that makes $p$ true is absent or different, while the basis is held fixed. \texttt{env}: a look-alike, a trick or a faulty device presents the subject with the same appearance. \texttt{time}: the moment at which the subject consults her source differs. \texttt{draw}: the outcome of a chance process differs. \texttt{marg}: a magnitude differs by less than the subject can discriminate. \texttt{src}: the testimonial source differs or is insincere. Two further tags mark worlds that are not error worlds: \texttt{basis}, where the subject lacks the basis, and \texttt{def}, where she meets a defeater and suspends judgment; \texttt{ok} marks a variation in which the belief remains true.

\begin{definition}[Gettier classes]\label{d:classes}
Let a Gettier case have error distance $\delta$ and let $\Sigma$ be the set of tags of its error worlds at distance $\delta$. Its class is fixed by the first applicable clause:
\begin{enumerate}[label=(\roman*)]
\item \textbf{G-nec} (necessary content): the belief is locally safe, so that only the basis-wide condition fails;
\item \textbf{G-tmp} (temporal): $\texttt{time} \in \Sigma$;
\item \textbf{G-lot} (lottery): $\texttt{draw} \in \Sigma$;
\item \textbf{G-mrg} (margin): $\texttt{marg} \in \Sigma$;
\item \textbf{G-tst} (testimonial): $\texttt{src} \in \Sigma$;
\item \textbf{G-env} (environmental): $\texttt{env} \in \Sigma$;
\item \textbf{G-fl} (false lemma): $\texttt{tm} \in \Sigma$, the route is inferential (inference, induction or deduction), and the basis relies on a false lemma;
\item \textbf{G-nfl} (no false lemma): $\texttt{tm} \in \Sigma$, the route is inferential, and there is no false lemma;
\item \textbf{G-int} (intervening): $\texttt{tm} \in \Sigma$ and the route is not inferential.
\end{enumerate}
\end{definition}

\begin{definition}[Locus]\label{d:locus}
A Gettier case is \emph{connected} if the fact that makes $p$ true is causally connected to the basis in the normal way, and \emph{disconnected} otherwise. We write the locus after the class, as in G-env/c.
\end{definition}

Connected Gettier cases are those that \citet{Pritchard2012} calls cases of environmental luck: the subject's faculties make contact with the fact, and the danger lies in what surrounds the contact. Disconnected cases are cases of intervening luck: something other than what the basis bears on makes the belief true. The locus is a separate dimension from the class: a testimonial case may be connected (HLC, where the director knows where the party is) or disconnected (STL, where the liar is right by mistake).

\begin{definition}[Non-Gettier outcomes]\label{d:other}
A model that is not a Gettier case is classed \textbf{U} if $J = 0$, \textbf{F} if $J = 1$ and $p(w_0) = 0$, and \textbf{K} otherwise. A K model is annotated with the kinds of luck it exhibits that are compatible with knowledge: \texttt{ev} (evidential), \texttt{rf} (reflective), \texttt{ct} (content), \texttt{md} (a misleading defeater exists), \texttt{fx} (the truth is secured by another agent's intervention).
\end{definition}

\begin{remark}\label{r:priority}
The order of the clauses in Definition~\ref{d:classes} settles the label when the nearest error worlds differ in more than one way. It has no effect on any verdict: whether a model is a Gettier case depends only on $\delta$ and $\theta$ (Lemma~\ref{l:delta}). The order puts the more specific descriptions first: a stopped clock is also a case in which the truth-maker varies, but the temporal description says how. Among the \stnmodels{} models, \stnmulti{} have nearest error worlds with more than one tag.
\end{remark}

The typology answers to the distinctions listed in the brief history of the problem. The inferential classes G-fl and G-nfl separate Gettier's cases and Lehrer's Nogot case from Feldman's version. G-env and G-int are Pritchard's environmental and intervening luck, refined by route. G-tmp, G-lot, G-tst and G-mrg are the classes that the literature treats under separate headings: the stopped clock, the lottery, testimony, and Williamson's margins for error. G-nec is the class of necessary truths reached on a defective basis, which a safety theorist must treat by the amendment. Reflective luck does not appear among the Gettier classes, because, as Proposition~\ref{p:reflective} shows, it is not a way of being unsafe.

\section{Theorems}\label{s:thm}

Each of the following is proved in prose from the definitions and checked by the program on every model to which it applies; the checks are assertions, and the program halts if one fails.

\begin{proposition}[Veritic luck]\label{p:veritic}
In a model with a justified true belief, the belief is veritically lucky if and only if it is not basis-wide safe. Hence the Gettier cases are exactly the models with a justified true belief that is veritically lucky.
\end{proposition}

\begin{proof}
Veritic luck is the existence of a close error world (Definition~\ref{d:luck}(a)); failure of basis-wide safety is the same condition (Definition~\ref{d:safe}). The second sentence follows from Definition~\ref{d:safe}.
\end{proof}

Proposition~\ref{p:veritic} is the model form of the reading of EPI-01-a's characterization theorem on which knowledge excludes veritic luck and no other kind. The remaining propositions show how each class realizes veritic luck.

\begin{proposition}[Evidential luck]\label{p:evidential}
Evidential luck neither entails nor excludes safety. A justified true belief that is evidentially lucky and not veritically lucky is knowledge.
\end{proposition}

\begin{proof}
At a world tagged \texttt{basis}, $B(v) = 0$, so such a world is not an error world and does not bear on safety. Whether a model has close \texttt{basis} worlds is therefore independent of whether it has close error worlds. The second sentence follows from Proposition~\ref{p:veritic} and Definition~\ref{d:safe}.
\end{proof}

Among the justified true models of the catalogue and variants, \stevsafe{} are evidentially lucky and safe, and \stevunsafe{} evidentially lucky and unsafe; SEV (the witness at the window) and GRN (Nozick's grandmother) are examples of the first kind, the mistaken liar (STL), who hears from a liar who could easily have lied about another train and not to him, is an example of the second.

\begin{proposition}[False lemmas]\label{p:lemma}
Adding or removing a false lemma while leaving the worlds of a model unchanged leaves its safety, and so its status as a Gettier case, unchanged. Consequently, relying on a false lemma is neither necessary nor sufficient for a justified true belief to be a Gettier case.
\end{proposition}

\begin{proof}
Safety is computed from the distances, belief values and truth values of the worlds alone (Definition~\ref{d:safe}); the lemma is not among them. For the second sentence, Feldman's version of the Nogot case (FEL) is a Gettier case without a false lemma, and Warfield's handout case (WAR) is a justified true belief relying on a false lemma whose model is safe.
\end{proof}

The proposition presupposes that removing the lemma does not change the worlds, and the presupposition is the substance of the matter. In Gettier's cases it holds because the falsity of the lemma and the truth of $p$ are independent: the error worlds of GE1 are worlds in which Smith's own pocket holds nine coins, and nothing about the route from the evidence about Jones to the conclusion bears on them. That is the observation of EPI-01-d, and the variants GE1.nfl, GE2.nfl, NOG.nfl and their kind confirm it model by model. In WAR the false belief that there are 53 people present is within the margin by which the conclusion, that 100 handouts suffice, is robust; the worlds in which the count is off by one or two are worlds in which the conclusion is still true. The no-false-lemmas condition fails in both directions for the same reason: it inspects the route of the inference, while knowledge depends on the modal profile of its conclusion.

\begin{proposition}[Environment and truth-maker]\label{p:env}
\begin{enumerate}[label=(\alph*)]
\item If every nearest error world of a Gettier case is tagged \texttt{tm}, \texttt{draw} or \texttt{time}, and no other close error world exists, then securing the truth-maker, so that at those worlds the subject would lack the basis, makes the belief safe.
\item If every close error world is tagged \texttt{env} or \texttt{src}, then normalizing the environment, by removing the look-alikes and the bad sources, makes the belief safe.
\item Placing a look-alike at distance $1$ makes any belief formed by perception, instrument or testimony unsafe; if the belief was safe, justified and true, the result is a case of class G-env.
\end{enumerate}
\end{proposition}

\begin{proof}
(a) Securing the truth-maker sets $B(v) = 0$ at the worlds in which the truth-maker varies, so they are no longer error worlds; if no other close error world exists, the belief is safe by Lemma~\ref{l:delta}. (b) Similarly, removing the worlds tagged \texttt{env} or \texttt{src} removes all close error worlds. (c) The added world is close at every $\theta \geq 1$ and is an error world. If the belief was safe, it had no close error world, so the new nearest error world is the added one and its tag is \texttt{env}.
\end{proof}

Proposition~\ref{p:env} explains why the two loci call for different remedies. Intervening luck is removed by connecting the basis to the truth-maker: Smith would know that the man who will get the job has ten coins if his evidence concerned his own pocket. Environmental luck is removed only by changing the environment: nothing Henry could do with his eyes from the road would make his belief safe in a district of facades. The variants marked \texttt{sec}, \texttt{nor} and \texttt{hos} in Appendix~B carry out these operations on every case to which they apply.

\begin{proposition}[Causal connection]\label{p:causal}
Among justified true beliefs, causal connection between the truth-maker and the basis is neither necessary nor sufficient for safety.
\end{proposition}

\begin{proof}
All four combinations are realized among the models (Table~\ref{tab:combos} in Appendix~D): BRN is causally connected and unsafe; TMP is not causally connected and safe; RED and GE1 realize the remaining two combinations.
\end{proof}

The connected and unsafe models are exactly the connected Gettier cases (G-$x$/c), \stlocc{} of the \stnG{} Gettier cases among all models; the remaining \stlocd{} are disconnected. A causal theory of knowledge in the manner of \citet{Goldman1967} therefore classifies the disconnected cases correctly and the connected ones incorrectly, which is the pattern Goldman himself recognized when he introduced the fake-barn case.

\begin{proposition}[Temporal cases]\label{p:temporal}
Let a belief about the time be formed by reading an instrument that has stopped. The belief is safe at $\theta$ if and only if no consultation at a different moment lies within distance $\theta$. If the instrument is going and accurate, the belief is safe whatever moment of consultation is close.
\end{proposition}

\begin{proof}
At a consultation at a different moment, a stopped instrument shows the same reading, the subject believes on the same basis that it is the time shown, and the belief is false; such a world is an error world, and safety requires that none be close. A going and accurate instrument shows the time at each moment, so at each consultation world the belief formed is true.
\end{proof}

\begin{proposition}[Necessary truths]\label{p:nec}
If $p$ is true at every world of a model, the belief that $p$ is locally safe at every threshold, and no such model is a Gettier case in the sense of Definition~2.4 of EPI-01-a. It is a Gettier case in the basis-wide sense exactly when, at some close world, the same basis yields a false belief in another proposition.
\end{proposition}

\begin{proof}
Local safety requires $p(v) = 1$ at close worlds where $B(v) = 1$, which holds when $p$ is true everywhere. The second sentence is Definition~\ref{d:safe} with $p$ true everywhere.
\end{proof}

The pupil of SNC, whose procedure gives the right product for $12 \times 13$ and wrong products for its neighbors, is the schematic case; SNS, with a sound procedure, is its control. The proposition records why the amendment of EPI-01-a is needed and what it costs: nothing, since in contingent cases the amended and unamended conditions agree in every model of the catalogue.

\begin{proposition}[Testimony]\label{p:testimony}
Suppose a hearer believes $p$ on the testimony of a speaker, at exactly the worlds at which the speaker believes $p$ on her own basis and sincerely asserts it, and suppose the speaker's sincerity is modally robust (no close world at which she asserts $p$ without believing it). Then the hearer's belief is safe if and only if the speaker's is. If instead there is a close world at which the speaker asserts $p$ while $p$ is false, the hearer's belief is unsafe, whatever the safety of the speaker's.
\end{proposition}

\begin{proof}
Under the first supposition the hearer's belief worlds and truth values coincide with the speaker's, so the two beliefs have the same close error worlds. Under the second, the world described is a close world at which the hearer believes $p$ on the same testimonial basis and $p$ is false.
\end{proof}

The first half is the model form of the thesis of EPI-08 that testimony transmits knowledge but does not generate it: a hearer's safety is borrowed from the speaker's. (In the relay variants the hearer's justification is likewise taken over from the speaker, as transmission requires.) The second half is the structure of the Halloween case on its coarse individuation (HLC) and of the mistaken liar (STL), the testimonial Gettier cases of class G-tst. Every case of the catalogue whose belief is true has two testimonial variants, \texttt{rel} and \texttt{lie}, and the program checks the proposition on all of them.

\begin{proposition}[Lotteries]\label{p:lottery}
In a fair lottery in which the draw of each ticket is as close as the draw of any other, the belief, formed on the odds, that one's ticket will lose is unsafe whatever the number of tickets. The belief, formed on a reliable report of the result, that one's ticket has lost is safe unless a misreport is close.
\end{proposition}

\begin{proof}
The draw of one's own ticket is a world at which one believes, on the odds, that it will lose, and in which it wins. Since it is as close as the actual draw, which is close by centring, it is close at every $\theta \geq d$, where $d$ is the common distance of the draws; the number of tickets does not enter. On the report basis, at the world in which one's ticket is drawn the report says so, and one does not believe that it lost; only a misreport would be an error world.
\end{proof}

This is the asymmetry between the odds and the newspaper (LOT and LTN). On the analysis it is the asymmetry between a Gettier case of class G-lot and knowledge, in agreement with EPI-09's conclusion that one does not know that one's ticket will lose, though one may rationally act as if it will.

\begin{proposition}[Margins and thresholds]\label{p:margin}
\begin{enumerate}[label=(\alph*)]
\item If a belief is safe at $\theta$, it is safe at every $\theta' \leq \theta$.
\item Moving every error world of a model to distance $4$ makes the belief safe at every $\theta \leq 3$ and unsafe at $\theta = 4$.
\item A belief that a magnitude is not $x$, formed by an estimate that cannot discriminate $x$ from the actual magnitude, is unsafe at every $\theta \geq 1$ when the world in which the magnitude is $x$ is at distance $1$.
\end{enumerate}
\end{proposition}

\begin{proof}
(a) The close worlds at $\theta'$ are among those at $\theta$. (b) By Lemma~\ref{l:delta}, since $\delta$ becomes $4$. (c) The world in which the magnitude is $x$ looks the same, so the subject believes on the same basis that it is not $x$, falsely; it is an error world at distance $1$.
\end{proof}

Clause~(c) is Williamson's margin-for-error principle in model form, and Proposition~8.1 of EPI-01-a is its general statement; MAG and MGW are the cases. Clause~(a) is what makes the contextualism of EPI-13 coherent with the analysis: raising the standards can turn knowledge into its absence, never the reverse, and a belief known in a demanding context is known in every less demanding one. The bank cases BKL and BKH are a single model at two thresholds.

\begin{proposition}[Reflective luck]\label{p:reflective}
Whether the subject has reflective access to the reliability of her basis does not enter into safety. A justified, safe belief that is reflectively lucky is knowledge.
\end{proposition}

\begin{proof}
Reflective access is a descriptive feature of the model, not a component of any world (Definition~\ref{d:model}); safety and knowledge are defined without it (Definition~\ref{d:safe}).
\end{proof}

The unreflective birder (SRF) is the schematic case. Norman (CLV) and Truetemp (TTP) are reflectively lucky too, but they are not justified by their sources' lights, and fail to know for that reason: their class is U, not a Gettier class. Whether a subject who lacks reflective access can be justified is a question for the theory of justification, which, by EPI-06, is internalist; the analysis of knowledge takes its verdict as given.

\begin{proposition}[Deduction]\label{p:deduction}
If the belief that $p$ is safe and $q$ is believed by competent deduction from it, the belief that $q$ is safe. The belief that $q$ may be safe when the belief that $p$ is not.
\end{proposition}

\begin{proof}
The first sentence is Proposition~6.1 of EPI-01-a: at every close world where the subject believes $q$ by deduction she believes $p$, which is true there, and $q$ is entailed by $p$. For the second, the variant BRN.ded is safe: Henry's belief that there is a barn in the district, deduced from his unsafe belief that what he is looking at is a barn, is true in every world of the model.
\end{proof}

The second sentence is less familiar than the first and worth noting. A Gettier case does not contaminate everything deduced from it: a weaker conclusion whose truth does not depend on the luck can be known. Smith knows, by deduction from his Gettiered belief, that someone will get the job (GE1.ded). Kripke's observation about the red barn is an instance of the first sentence: Henry's safe belief that he sees a red barn yields, by deduction, a safe belief that he sees a barn (RED.ded), while his perceptual belief that he sees a barn, formed on shape alone, is a Gettier case of class G-env (RDP).

We can now state the partition result.

\begin{theorem}[Partition]\label{t:partition}
Every case model falls under exactly one of the classes G-nec, G-tmp, G-lot, G-mrg, G-tst, G-env, G-fl, G-nfl, G-int, K, U and F. A model falls under one of the nine Gettier classes if and only if it has a justified true belief that is not basis-wide safe.
\end{theorem}

\begin{proof}
If $J = 0$ the model is U, and if $J = 1$ and $p(w_0) = 0$ it is F. Suppose $J = 1$ and $p(w_0) = 1$. If the belief is basis-wide safe, the model is K. Otherwise it is a Gettier case, and either it is locally safe, and then G-nec, or there is a close error world for $p$, so that $\delta \leq \theta$ and the set $\Sigma$ of tags at distance $\delta$ is non-empty. Error worlds carry one of the six tags \texttt{time}, \texttt{draw}, \texttt{marg}, \texttt{src}, \texttt{env} and \texttt{tm} (the remaining tags mark worlds at which the belief is true or absent), so one of the clauses (ii) to (ix) of Definition~\ref{d:classes} applies, and the first applicable one fixes the class. The classes are exclusive because the clauses are applied in order and U, F and K are separated by $J$, truth and safety. The second sentence is the case distinction just made.
\end{proof}

\begin{corollary}\label{c:exactly}
In case models, the Gettier cases of every class are exactly the justified true beliefs that are unsafe; and each class is a way of being unsafe, specified by what varies in the nearest error worlds.
\end{corollary}

Corollary~\ref{c:exactly} is the conclusion of EPI-01 stated for the classification. It says that the heterogeneity of the Gettier cases is real, and that it is a heterogeneity of mechanisms, not of kinds: what differs from class to class is what could easily have been otherwise, while what makes each member a failure of knowledge is the same.

\section{The catalogue and the check}\label{s:cat}

\subsection{What is encoded}

Appendix~A gives \stnsources{} case models. \stnlit{} encode cases from the literature whose source we can name with confidence: Gettier's two cases; Lehrer's Nogot case and Feldman's version without a false lemma; Russell's stopped clock; the fake barns of \citet{Goldman1976}, which Goldman credits to Carl Ginet, and Kripke's red barn; Chisholm's sheep; Zagzebski's two cases; Skyrms's pyromaniac; Tom Grabit; Harman's assassination case; BonJour's Norman and Lehrer's Truetemp; Sosa's rubbish chute and kaleidoscope; Pritchard's Temp; Lackey's Morris; Hawthorne's lottery and newspaper; Williamson's Mr Magoo; Dretske's zebras; Nozick's grandmother; Comesa\~na's Halloween party and Bogardus's atomic clock, each on two individuations of the basis; Warfield's handouts; DeRose's bank cases; Sosa's belief that one is not a brain in a vat, Nozick's Jesse James, and Vogel's ice cubes \citep{Vogel1987}; and, as the oldest entry, the jury of Plato's \emph{Theaetetus} (201a--c), brought to a true judgment by persuasion \citep{PlatoTht}. The remaining \stnschem{} are schematic cases constructed to fill classes that the literature discusses in general terms (necessary truths, evidential luck, content luck, reflective luck, the mistaken liar). Each entry gives the source, a description in our words, the proposition, the basis and its route, the descriptive features, the worlds of the model in a table, and the computed verdicts.

Appendix~B applies to every case each of eleven operations that changes one feature: removing or adding a false lemma (\texttt{nfl}, \texttt{fl}); securing the truth-maker (\texttt{sec}); normalizing the environment (\texttt{nor}) or making it hostile (\texttt{hos}); moving the error worlds to distance $4$ (\texttt{rem}); making the belief false (\texttt{fal}); removing justification (\texttt{unj}); relaying the belief by sincere testimony (\texttt{rel}) or by a speaker who could easily have lied (\texttt{lie}); and competent deduction of a weaker proposition (\texttt{ded}). An operation is applied where it makes sense, which yields \stnvariants{} variants. For each, the program checks the prediction of the proposition that governs it. Appendix~C gives, for all \stnmodels{} models, the error distance, the safety profile, and the verdicts of four rival conditions added to justified true belief: no false lemmas (NFL), causal connection (CAU), sensitivity (SEN) and the absence of defeaters (UND), next to safety (SAF). Appendix~D gives the incidence of the classes and the summary of the check.

\subsection{The check}

The program verifies three things about the catalogue, besides the propositions of Section~\ref{s:thm}.

First, every case that its source presents as a Gettier case, a counterexample to the sufficiency of justified true belief, is computed to be a justified true belief that is unsafe. There are \stnGrows{} such cases, and the check holds in all of them; each has error distance $1$, so that the verdict holds at every threshold from $1$ upward.

Second, among the \stnundisputed{} literature cases that carry no dispute record of kind \texttt{dep} or \texttt{bas}, the source's verdict on knowledge coincides with the computed one in every case: the source says \enquote{not knowledge} of a justified true belief exactly when the belief is computed to be unsafe.

Third, the cases in which the computed verdict departs from the source's are listed with a reason (Table~\ref{tab:departures} in Appendix~D). There are \stndepart{}: HAR, KAL, TMP, ZBM, HLC and ATC. We discuss them below. In each, either the literature is divided on the verdict, or the verdict turns on the individuation of the basis, so that another row of the catalogue, with the basis individuated differently, agrees with the source.

Among the rival conditions, safety is the only one that classifies every undisputed literature case correctly (Table~\ref{tab:rivals}); the others err on between three and eight cases each. On the variants, where no literature verdict exists, the table records how often each rival agrees with safety. Sensitivity agrees most often, which is what one would expect of a condition that, in ordinary cases, consults the same worlds.

\subsection{Particular entries}

\paragraph{GE1, FEL and the false lemma.} Gettier's first case (GE1) is of class G-fl/d: its nearest error world is the one in which Smith's pocket holds nine coins, tagged \texttt{tm}, and the truth-maker is disconnected from the evidence about Jones. Its variant GE1.nfl, in which the same evidence reaches the conclusion through true intermediates, has the same worlds and the same error distance, and is of class G-nfl/d; Feldman's case (FEL) is the literature's instance of that class. The variant GE1.sec, in which Smith's evidence concerns his own pocket, is knowledge, although it keeps the false lemma; so is GE1.ded, Smith's belief that someone will get the job. The four rows together show the false lemma to be neither necessary nor sufficient, as Proposition~\ref{p:lemma} says, and locate the defect where EPI-01-d locates it.

\paragraph{BRN, RDP and RED.ded.} The fake-barn case is G-env/c, the paradigm of connected luck: Henry's belief is caused by the barn and is unsafe because facades are a glance away. Its variant BRN.nor, with the facades removed, is knowledge, and BRN.sec does not exist, since there is no truth-maker to secure. Kripke's red barn (RED) is safe, because a glance at a blue facade does not produce the belief that one sees a red barn; RDP, the belief formed on shape alone in the same district, is G-env/c; and RED.ded, the deduced belief that one sees a barn, is safe. The three rows are the formal content of the reply that EPI-01-a gives to Kripke.

\paragraph{HLF and HLC.} Comesa\~na's Halloween case appears twice. On the fine individuation of the basis (HLF), the guest believes on the testimony of a director disposed to tell him the truth; at the world in which the circumstances obtain, the director's disposition and so the basis differ, the guest does not believe on that basis, and the belief is safe: knowledge, as Comesa\~na holds. On the coarse individuation (HLC), the same world is an error world, and the case is G-tst/c. Bogardus's atomic clock (ATF, ATC) has the same structure with an instrument in place of a speaker. The analysis does not choose between the individuations; Assumption~\ref{a:faithful} is silent on it. What the pair shows is that the dispute about these cases is a dispute about bases, not about safety, which is the reply of EPI-01-a.

\paragraph{TMP and MOR.} Pritchard's Temp is the one case of the catalogue in which the computed verdict is knowledge and the source's verdict is not, by reason of how the truth is secured: by another agent who adjusts the world to the belief. The belief is justified on internalist standards, being formed by reading an instrument Temp has no reason to distrust, and it is safe. We accept the verdict that Temp knows, for a reason the catalogue makes visible. Lackey's Morris (MOR), who believes the directions of a passer-by, is another case of a safe belief whose safety is owed to another person rather than to the subject's own competence, and the source's verdict is that Morris knows. The difference between the cases lies in the direction of fit of the helper's contribution: the passer-by makes Morris's belief fit the world, while Temp's helper makes the world fit Temp's belief. That difference bears on whether the subject deserves credit for the truth of his belief. It does not bear on whether the belief could easily have been false, and on the analysis it is the latter that knowledge requires. TMP carries the annotation \texttt{fx} for this reason.

\paragraph{LOT and LTN.} The lottery belief (LOT) is G-lot/d, with error distance $1$ in any fair lottery (Proposition~\ref{p:lottery}); the newspaper belief (LTN) is knowledge, with error distance $3$, the distance of a misprint. Its profile \texttt{SSSuu} records that in a context in which misprints are attended to, it ceases to be knowledge, which is the contextualist's account of the sceptical pressure that lottery propositions exert (EPI-13).

\paragraph{The departures.} Harman's assassination case (HAR) is computed to be knowledge: Jill's belief is true at every close world at which she believes on the basis of the report, since at the worlds in which she hears the denials she does not believe at all. Harman's verdict is that she does not know. The structure of the case is that of Tom Grabit (GRB), whose source's verdict is knowledge: in both, a misleading defeater exists which the subject does not possess, and the existence of the defeater makes no difference to whether her belief could easily have been false. A verdict that is to be consistent across the two cases must treat them alike, and the analysis does so. Sosa's kaleidoscope (KAL) is computed to be a Gettier case of class G-env/c, where Sosa holds that the perceiver has animal but not reflective knowledge; on the analysis, the jokester's easy alternative makes the belief unsafe, and the case is a fake-barn case with a person in place of the facades. Dretske's disguised mules (ZBM) are computed to be known, against Dretske, because the belief is insensitive but safe; this is the conclusion of EPI-20, and the catalogue shows why the two verdicts part: SEN and SAF differ on ZBM as they do on the rubbish chute (CHU), the ice cubes (ICE), the brain in a vat (BIV), the zebras themselves (ZBP) and Tom Grabit (GRB): in each, the nearest world in which the belief is false is one in which the subject would still hold it, but that world is remote. TMP was discussed above, and HLC and ATC are the coarse individuations of cases whose fine individuations agree with their sources.

\section{Objections and replies}\label{s:obj}

\begin{objection}[The encodings decide the verdicts]\label{o:encode}
The program computes safety from distances and truth values which we chose. If one wanted a different verdict, one could choose different distances. The agreement between the computed verdicts and the literature is manufactured.
\end{objection}

\begin{reply}
The distances encode what the case descriptions say could easily have happened: that Smith could easily have had nine coins, that Henry's glance could easily have fallen on a facade, that the clock could easily have been consulted five minutes later. These are not verdicts about knowledge, and Assumption~\ref{a:faithful} makes explicit that the encoding is answerable to the descriptions, not to the analysis. Two features of the computed material limit the room for manipulation. First, every case presented as a Gettier case has error distance $1$: the verdict holds at every positive threshold, and no choice of threshold short of $0$ would reverse it. Of the \stnjtsrc{} justified true source beliefs, only \stnnonrobust{} have verdicts that change as $\theta$ ranges over $1$ to $3$, and they are RED, TMP, MOR, LTN and the bank cases BKL and BKH, which are the cases whose descriptions mention a remote or context-dependent possibility. Second, the variants are generated by fixed operations from the source encodings, and their verdicts are predicted by the propositions of Section~\ref{s:thm} before being computed; an encoding tuned to one verdict would not, without luck, produce correct predictions for the eleven operations.
\end{reply}

\begin{objection}[The classes are arbitrary]\label{o:arbitrary}
The order of the clauses in Definition~\ref{d:classes} is stipulated. A different order would yield a different classification, and nothing in the analysis privileges this one.
\end{objection}

\begin{reply}
The order matters only to the label of a model whose nearest error worlds differ in more than one way (Remark~\ref{r:priority}), which happens in \stnmulti{} of \stnmodels{} models, and never to its verdict. The order puts specific descriptions before general ones: every temporal, lottery, margin, testimonial and environmental error world is also, in a broad sense, one in which what makes the belief true is absent, and the clause for \texttt{tm} is the residual case. A typology must make some such choice; what Theorem~\ref{t:partition} establishes is that any order yields a partition of the unsafe justified true beliefs, because the verdict depends only on $\delta$ and $\theta$.
\end{reply}

\begin{objection}[The basis is gerrymandered]\label{o:basis}
The catalogue lists the Halloween and atomic-clock cases twice, once with each individuation of the basis, and counts the fine individuation as agreeing with the source. This is the generality problem in another dress: by choosing the basis one can make any belief safe or unsafe.
\end{objection}

\begin{reply}
We agree that safety is relative to a basis and that the analysis does not reduce bases to anything else; EPI-01-a accepts this in its reply to the objection from generality. The catalogue does not hide the relativity; it displays it, and it displays its limits. Of the \stnlit{} literature rows only five carry a dispute about individuation: the two cases encoded twice, and RDP, whose deductive counterpart RED.ded agrees with Kripke; and in none of the cases presented as Gettier cases is there an individuation, respecting the description, on which the belief is safe: Smith's basis is his evidence about Jones under any description, and Henry's basis is the look of a barn from the road. The generality problem is a problem about where the boundaries of a basis lie in hard cases, not a licence to redraw them in easy ones.
\end{reply}

\begin{objection}[Safety is not sufficient]\label{o:sufficient}
Temp shows that a justified and safe belief can fail to be knowledge. The catalogue concedes that the analysis departs from the source's verdict, which is to concede a counterexample.
\end{objection}

\begin{reply}
Temp is not a Gettier case on any account, since the belief is safe; the characterization of the Gettier cases is untouched. The objection concerns the sufficiency of justified safe belief, and our reply was given with the entry: the feature that distinguishes Temp from Morris is the direction of fit of the helper's contribution, which bears on credit and not on the modal relation between belief and truth. \citet{Pritchard2012} adds an ability condition to safety in order to deny Temp knowledge; EPI-01-a holds that the justification condition does the work of such a condition where it needs doing, and Temp, unlike a subject whose beliefs are made true by a demon without grounds, has grounds. The cost of the verdict is that Temp knows something for which he deserves little credit; Morris shows that this cost is one any view must pay.
\end{reply}

\begin{objection}[Intuitions vary]\label{o:xphi}
Experimental work suggests that intuitions about Gettier cases are not uniform. \citet{WeinbergNicholsStich2001} reported differences between cultural groups; \citet{StarmansFriedman2012} found that lay participants attribute knowledge in some Gettier cases; and \citet{Colaco2014} found that many attribute knowledge in fake-barn cases. A classification calibrated to the verdicts of philosophers may be calibrated to a parochial concept.
\end{objection}

\begin{reply}
The findings are mixed in a way that the typology helps to sort. A large cross-cultural study found the core Gettier intuition to be robust across languages and cultures \citep{Machery2017}, while the instability reported for fake barns \citep{GendlerHawthorne2005,Colaco2014} concerns a case of the connected class G-env/c, in which the subject's faculties make normal contact with the fact. It is to be expected that intuitions should be least firm where the actual causal route is normal and the danger lies in alternatives the subject never confronts. The analysis does not predict that every speaker will apply \enquote{knows} in accordance with it; it predicts that the cases in which competent judges withhold knowledge from justified true belief are the unsafe ones, and that disagreement will concentrate where unsafety is least visible. The catalogue records the fake-barn dispute and retains the source's verdict.
\end{reply}

\begin{objection}[Not all of these are Gettier cases]\label{o:extension}
Lotteries, margins for error and the bank cases are not usually called Gettier cases. Counting them as such inflates the thesis.
\end{objection}

\begin{reply}
The catalogue records the role each case plays in its source, and the check that every case presented as a Gettier case is a justified true unsafe belief is restricted to those so presented. That lottery beliefs and margin beliefs are justified, true and not knowledge is the reason they appear in the literature at all; the thesis of EPI-01 is that they fail for the same reason as Gettier's cases, and Definition~2.4 of EPI-01-a accordingly counts them as Gettier cases in its technical sense. The typology gives them their own classes, G-lot and G-mrg, so that nothing is lost by the extension, and Proposition~\ref{p:lottery} and Proposition~\ref{p:margin} say how each class is unsafe.
\end{reply}

\section{What the classification does not establish}\label{s:not}

The classification shows that the cases of the catalogue, as encoded, are classified correctly by safety and that the classes partition the unsafe justified true beliefs. It does not establish four things.

It does not establish that the encodings are the only faithful ones. Assumption~\ref{a:faithful} is a substantive premise, and a reader who thinks that, say, Russell's clock could not easily have been consulted at another moment will compute a different verdict from our tables. The tables make such disagreements precise, which is their use.

It does not settle the individuation of bases. The paired rows show where the verdict depends on it, and the paired rows are few, but the analysis inherits the generality problem in the form EPI-01-a states.

It does not show that the catalogue exhausts the literature. We have encoded only cases whose source we can identify with confidence, and the schematic cases are ours. The theorems of Section~\ref{s:thm} hold for all case models, not only for the catalogued ones; what the catalogue adds is that the literature's cases are case models of the kinds the theorems describe.

It does not adjudicate the experimental findings. Objection~\ref{o:xphi} explains how the typology bears on them; it does not show that the concept the analysis captures is shared by every speaker.

\section{Conclusion}

The Gettier cases are of many kinds, and the typology of Section~\ref{s:typ} sorts them by what varies in the nearest worlds in which the subject, on the same basis, believes falsely: the truth-maker, behind an inference with or without a false lemma or in perception; the environment; the time of consultation; the outcome of a chance; the source of testimony; a magnitude within the margin of discrimination; or, for necessary truths, the outputs of the basis on neighboring questions. For each kind a theorem states how safety fails and what would restore it, and the kinds together partition the unsafe justified true beliefs. Applied to the literature, the classification places every case presented as a Gettier case among them, and the cases in which its verdict departs from a source's are cases on which the literature is divided or whose basis admits two individuations. Gettier cases are exactly those in which a justified true belief is unsafe.

\appendix

\section{Catalogue of cases}\label{app:a}

\paragraph{What is encoded.} Each entry is one case model (Definition~\ref{d:model}). The header gives the identifier, the source, the route and the basis; then a description of the case in our words; then the proposition $p$, the weaker proposition $q$ where one is used for the deduction variant, and the descriptive features. The table lists the worlds of the model: the actual world first, then each other world with a short statement of how it differs, its distance $d$, whether it is close at the case's threshold $\theta$, whether the subject believes $p$ on the same basis there ($B$), whether $p$ is true there, its tag, and where used the truth of $q$ and the value \enquote{alt} of Definition~\ref{d:model}.

\paragraph{What is computed.} For each model the program computes justification and truth (read from the encoding), the error distance $\delta$, local and basis-wide safety at the case's threshold, whether the model is a Gettier case, whether the subject knows, the safety profile over $\theta = 0, \ldots, 4$, and the class (Definitions~\ref{d:classes} to \ref{d:other}). Where the source gives a verdict on knowledge, the entry records it, with the role the case plays in the source and whether the computed verdict agrees. Disputes are recorded with their kind: \texttt{int} (intuitions disputed, verdict retained), \texttt{dep} (computed verdict departs from the source's) or \texttt{bas} (verdict depends on the individuation of the basis).

\paragraph{Summary.} Of the \stnsources{} models, \stnlit{} are from the literature and \stnschem{} schematic. In \stnagree{} of the \stnlit{} literature cases the computed verdict agrees with the source's; the \stndepart{} departures are listed in Table~\ref{tab:departures}.

\needspace{12\baselineskip}
\subsection*{GE1 \enspace Ten coins}
\addcontentsline{toc}{subsection}{GE1 Ten coins}
\noindent\textit{Source.} \citet{Gettier1963}, case I.\quad \textit{Route:} inference.\quad \textit{Basis:} inference from evidence about Jones.\par
\noindent\textit{Case.} Smith has strong evidence that Jones will be offered the job for which both have applied, and that Jones has ten coins in his pocket. He infers that the man who will get the job has ten coins in his pocket. In fact Smith gets the job, and, without knowing it, has ten coins in his own pocket.\par
\noindent\textit{Belief} $p$: the man who will get the job has ten coins in his pocket.
 \textit{Weaker} $q$: someone will get the job.
 \textit{False lemma:} Jones will get the job. \textit{Causal connection:} no. \textit{Defeater:} genuine. \textit{Reflective access:} yes. \textit{Threshold} $\theta = 2$.\par
\begin{center}\begin{tabular}{@{}p{0.43\linewidth}rccclc@{}}
world & $d$ & close & $B$ & $p$ & tag & $q$ \\ \midrule
actual world & 0 & yes & yes & yes & act & yes \\
Smith has nine coins in his pocket & 1 & yes & yes & no & tm & yes \\
a third applicant is hired & 2 & yes & yes & no & tm & yes \\
Jones is hired, as the evidence indicated & 2 & yes & yes & yes & ok & yes \\
Smith never hears the assurance about Jones & 3 & no & no & yes & basis & yes \\
\end{tabular}\end{center}
\noindent\textit{Computed.} J yes, T yes, error distance $\delta = 1$, locally safe no, safe basis-wide no, Gettier case yes, knowledge no. Profile over $\theta = 0..4$: \texttt{Suuuu}. Class: \textbf{G-fl/d}.
 \textit{Source verdict:} not knowledge (presented as a Gettier case). Agreement: yes.
\par\medskip
\needspace{12\baselineskip}
\subsection*{GE2 \enspace Ford or Barcelona}
\addcontentsline{toc}{subsection}{GE2 Ford or Barcelona}
\noindent\textit{Source.} \citet{Gettier1963}, case II.\quad \textit{Route:} inference.\quad \textit{Basis:} disjunction introduced on evidence about Jones.\par
\noindent\textit{Case.} Smith has strong evidence that Jones owns a Ford. Choosing place names at random, he infers the disjunction that either Jones owns a Ford or his friend Brown is in Barcelona. Jones does not own a Ford, but Brown, by coincidence, is in Barcelona.\par
\noindent\textit{Belief} $p$: either Jones owns a Ford or Brown is in Barcelona.
 \textit{False lemma:} Jones owns a Ford. \textit{Causal connection:} no. \textit{Defeater:} genuine. \textit{Reflective access:} yes. \textit{Threshold} $\theta = 2$.\par
\begin{center}\begin{tabular}{@{}p{0.43\linewidth}rcccl@{}}
world & $d$ & close & $B$ & $p$ & tag \\ \midrule
actual world & 0 & yes & yes & yes & act \\
Brown is in Boston & 1 & yes & yes & no & tm \\
Brown is in Brest-Litovsk & 1 & yes & yes & no & tm \\
Jones owns the car he has been driving & 2 & yes & yes & yes & ok \\
\end{tabular}\end{center}
\noindent\textit{Computed.} J yes, T yes, error distance $\delta = 1$, locally safe no, safe basis-wide no, Gettier case yes, knowledge no. Profile over $\theta = 0..4$: \texttt{Suuuu}. Class: \textbf{G-fl/d}.
 \textit{Source verdict:} not knowledge (presented as a Gettier case). Agreement: yes.
\par\medskip
\needspace{12\baselineskip}
\subsection*{NOG \enspace Nogot and Havit}
\addcontentsline{toc}{subsection}{NOG Nogot and Havit}
\noindent\textit{Source.} \citet{Lehrer1965}.\quad \textit{Route:} inference.\quad \textit{Basis:} inference via a belief about Nogot.\par
\noindent\textit{Case.} A colleague, Nogot, gives the subject ample evidence that he owns a Ford: he says so, is seen driving one and shows papers. The subject infers that Nogot owns a Ford, and from that, that someone in the office owns a Ford. Nogot is deceiving; but another colleague, Havit, of whose car the subject knows nothing, owns a Ford.\par
\noindent\textit{Belief} $p$: someone in the office owns a Ford.
 \textit{Weaker} $q$: someone in the office has given me evidence that he owns a Ford.
 \textit{False lemma:} Nogot owns a Ford. \textit{Causal connection:} no. \textit{Defeater:} genuine. \textit{Reflective access:} yes. \textit{Threshold} $\theta = 2$.\par
\begin{center}\begin{tabular}{@{}p{0.43\linewidth}rccclc@{}}
world & $d$ & close & $B$ & $p$ & tag & $q$ \\ \midrule
actual world & 0 & yes & yes & yes & act & yes \\
Havit has sold his Ford & 1 & yes & yes & no & tm & yes \\
Havit works in another office & 1 & yes & yes & no & tm & yes \\
Nogot is sincere & 2 & yes & yes & yes & ok & yes \\
\end{tabular}\end{center}
\noindent\textit{Computed.} J yes, T yes, error distance $\delta = 1$, locally safe no, safe basis-wide no, Gettier case yes, knowledge no. Profile over $\theta = 0..4$: \texttt{Suuuu}. Class: \textbf{G-fl/d}.
 \textit{Source verdict:} not knowledge (presented as a Gettier case). Agreement: yes.
\par\medskip
\needspace{12\baselineskip}
\subsection*{FEL \enspace Nogot without the lemma}
\addcontentsline{toc}{subsection}{FEL Nogot without the lemma}
\noindent\textit{Source.} \citet{Feldman1974}.\quad \textit{Route:} inference.\quad \textit{Basis:} existential inference from true premises about testimony.\par
\noindent\textit{Case.} The evidence of the Nogot case, but the subject reasons only through true intermediate beliefs: from the fact that a colleague in the office has told him, and shown him, that he owns a Ford, he infers directly that someone in the office owns a Ford, without forming the belief that Nogot in particular owns one.\par
\noindent\textit{Belief} $p$: someone in the office owns a Ford.
 \textit{Weaker} $q$: someone in the office has given me evidence that he owns a Ford.
 \textit{False lemma:} none. \textit{Causal connection:} no. \textit{Defeater:} genuine. \textit{Reflective access:} yes. \textit{Threshold} $\theta = 2$.\par
\begin{center}\begin{tabular}{@{}p{0.43\linewidth}rccclc@{}}
world & $d$ & close & $B$ & $p$ & tag & $q$ \\ \midrule
actual world & 0 & yes & yes & yes & act & yes \\
Havit has sold his Ford & 1 & yes & yes & no & tm & yes \\
Havit works in another office & 1 & yes & yes & no & tm & yes \\
Nogot is sincere & 2 & yes & yes & yes & ok & yes \\
\end{tabular}\end{center}
\noindent\textit{Computed.} J yes, T yes, error distance $\delta = 1$, locally safe no, safe basis-wide no, Gettier case yes, knowledge no. Profile over $\theta = 0..4$: \texttt{Suuuu}. Class: \textbf{G-nfl/d}.
 \textit{Source verdict:} not knowledge (presented as a Gettier case). Agreement: yes.
\par\medskip
\needspace{12\baselineskip}
\subsection*{CLK \enspace Stopped clock}
\addcontentsline{toc}{subsection}{CLK Stopped clock}
\noindent\textit{Source.} \citet{Russell1948}.\quad \textit{Route:} instrument.\quad \textit{Basis:} reading the clock.\par
\noindent\textit{Case.} A person looks at a clock that has stopped, though she takes it to be going, and forms a belief about the time from its reading. She happens to look at one of the two moments of the day at which the reading is correct.\par
\noindent\textit{Belief} $p$: it is now the time the clock shows.
 \textit{False lemma:} the clock is going. \textit{Causal connection:} no. \textit{Defeater:} genuine. \textit{Reflective access:} yes. \textit{Threshold} $\theta = 2$.\par
\begin{center}\begin{tabular}{@{}p{0.43\linewidth}rcccl@{}}
world & $d$ & close & $B$ & $p$ & tag \\ \midrule
actual world & 0 & yes & yes & yes & act \\
she looks five minutes later & 1 & yes & yes & no & time \\
she looks five minutes earlier & 1 & yes & yes & no & time \\
the clock is going & 3 & no & yes & yes & ok \\
\end{tabular}\end{center}
\noindent\textit{Computed.} J yes, T yes, error distance $\delta = 1$, locally safe no, safe basis-wide no, Gettier case yes, knowledge no. Profile over $\theta = 0..4$: \texttt{Suuuu}. Class: \textbf{G-tmp/d}.
 \textit{Source verdict:} not knowledge (presented as a Gettier case). Agreement: yes.
\par\medskip
\needspace{12\baselineskip}
\subsection*{BRN \enspace Fake barns}
\addcontentsline{toc}{subsection}{BRN Fake barns}
\noindent\textit{Source.} \citet{Goldman1976}, credited to Ginet.\quad \textit{Route:} perception.\quad \textit{Basis:} visual appearance from the road.\par
\noindent\textit{Case.} Henry drives through a district in which, unknown to him, many facades have been put up which look from the road exactly like barns. He looks at the one real barn in view and believes that it is a barn.\par
\noindent\textit{Belief} $p$: what I am looking at is a barn.
 \textit{Weaker} $q$: there is a barn in this district.
 \textit{False lemma:} none. \textit{Causal connection:} yes. \textit{Defeater:} genuine. \textit{Reflective access:} yes. \textit{Threshold} $\theta = 2$.\par
\begin{center}\begin{tabular}{@{}p{0.43\linewidth}rccclc@{}}
world & $d$ & close & $B$ & $p$ & tag & $q$ \\ \midrule
actual world & 0 & yes & yes & yes & act & yes \\
his gaze falls on a facade & 1 & yes & yes & no & env & yes \\
he looks a moment earlier, at another facade & 1 & yes & yes & no & env & yes \\
no facades have been built & 3 & no & yes & yes & ok & yes \\
\end{tabular}\end{center}
\noindent\textit{Computed.} J yes, T yes, error distance $\delta = 1$, locally safe no, safe basis-wide no, Gettier case yes, knowledge no. Profile over $\theta = 0..4$: \texttt{Suuuu}. Class: \textbf{G-env/c}.
 \textit{Source verdict:} not knowledge (presented as a Gettier case). Agreement: yes.
 \textit{Dispute (int):} Intuitions are reported to be unstable across variants \citep{GendlerHawthorne2005} and many lay participants attribute knowledge \citep{Colaco2014}.
\par\medskip
\needspace{12\baselineskip}
\subsection*{RED \enspace Red barn}
\addcontentsline{toc}{subsection}{RED Red barn}
\noindent\textit{Source.} \citet{Kripke2011}.\quad \textit{Route:} perception.\quad \textit{Basis:} visual appearance, including color.\par
\noindent\textit{Case.} As in the fake-barn district, except that every facade is blue and the one real barn is red. Henry looks at the red barn and believes that he sees a red barn.\par
\noindent\textit{Belief} $p$: what I am looking at is a red barn.
 \textit{Weaker} $q$: what I am looking at is a barn.
 \textit{False lemma:} none. \textit{Causal connection:} yes. \textit{Defeater:} none. \textit{Reflective access:} yes. \textit{Threshold} $\theta = 2$.\par
\begin{center}\begin{tabular}{@{}p{0.43\linewidth}rccclc@{}}
world & $d$ & close & $B$ & $p$ & tag & $q$ \\ \midrule
actual world & 0 & yes & yes & yes & act & yes \\
his gaze falls on a blue facade & 1 & yes & no & no & env & no \\
he looks at another blue facade & 1 & yes & no & no & env & no \\
a red facade has been built & 3 & no & yes & no & env & no \\
\end{tabular}\end{center}
\noindent\textit{Computed.} J yes, T yes, error distance $\delta = 3$, locally safe yes, safe basis-wide yes, Gettier case no, knowledge yes. Profile over $\theta = 0..4$: \texttt{SSSuu}. Class: \textbf{K}.
 \textit{Source verdict:} knowledge (presented against a fourth condition). Agreement: yes.
\par\medskip
\needspace{12\baselineskip}
\subsection*{RDP \enspace Red barn, shape only}
\addcontentsline{toc}{subsection}{RDP Red barn, shape only}
\noindent\textit{Source.} \citet{Kripke2011}.\quad \textit{Route:} perception.\quad \textit{Basis:} visual appearance of shape.\par
\noindent\textit{Case.} In the district of the red-barn case, Henry believes that he sees a barn on the basis of the shape he sees, not of its color.\par
\noindent\textit{Belief} $p$: what I am looking at is a barn.
 \textit{False lemma:} none. \textit{Causal connection:} yes. \textit{Defeater:} genuine. \textit{Reflective access:} yes. \textit{Threshold} $\theta = 2$.\par
\begin{center}\begin{tabular}{@{}p{0.43\linewidth}rcccl@{}}
world & $d$ & close & $B$ & $p$ & tag \\ \midrule
actual world & 0 & yes & yes & yes & act \\
his gaze falls on a blue facade & 1 & yes & yes & no & env \\
he looks at another blue facade & 1 & yes & yes & no & env \\
\end{tabular}\end{center}
\noindent\textit{Computed.} J yes, T yes, error distance $\delta = 1$, locally safe no, safe basis-wide no, Gettier case yes, knowledge no. Profile over $\theta = 0..4$: \texttt{Suuuu}. Class: \textbf{G-env/c}.
 \textit{Source verdict:} not knowledge (presented against a fourth condition). Agreement: yes.
 \textit{Dispute (bas):} Kripke's point is that Henry knows he sees a barn; on the analysis he does so on the deductive basis (variant RED.ded), not on this one.
\par\medskip
\needspace{12\baselineskip}
\subsection*{SHP \enspace Sheep in the field}
\addcontentsline{toc}{subsection}{SHP Sheep in the field}
\noindent\textit{Source.} \citet{Chisholm1966}.\quad \textit{Route:} perception.\quad \textit{Basis:} seeing a sheep-like animal.\par
\noindent\textit{Case.} A person looks into a field and sees an animal which looks like a sheep, and believes that there is a sheep in the field. The animal is a dog; but there is a sheep in the field, hidden from view.\par
\noindent\textit{Belief} $p$: there is a sheep in the field.
 \textit{False lemma:} that animal is a sheep. \textit{Causal connection:} no. \textit{Defeater:} genuine. \textit{Reflective access:} yes. \textit{Threshold} $\theta = 2$.\par
\begin{center}\begin{tabular}{@{}p{0.43\linewidth}rcccl@{}}
world & $d$ & close & $B$ & $p$ & tag \\ \midrule
actual world & 0 & yes & yes & yes & act \\
the sheep is in the neighboring field & 1 & yes & yes & no & tm \\
the dog is not in the field & 3 & no & no & yes & basis \\
\end{tabular}\end{center}
\noindent\textit{Computed.} J yes, T yes, error distance $\delta = 1$, locally safe no, safe basis-wide no, Gettier case yes, knowledge no. Profile over $\theta = 0..4$: \texttt{Suuuu}. Class: \textbf{G-int/d}.
 \textit{Source verdict:} not knowledge (presented as a Gettier case). Agreement: yes.
\par\medskip
\needspace{12\baselineskip}
\subsection*{ZDR \enspace The diagnosis}
\addcontentsline{toc}{subsection}{ZDR The diagnosis}
\noindent\textit{Source.} \citet{Zagzebski1994}.\quad \textit{Route:} inference.\quad \textit{Basis:} inference from symptoms and tests.\par
\noindent\textit{Case.} A physician, on excellent evidence from symptoms and tests, judges that a patient has virus X. The symptoms are in fact produced by another virus; but the patient has also, very recently, contracted X, which has not yet produced symptoms.\par
\noindent\textit{Belief} $p$: the patient has virus X.
 \textit{False lemma:} the symptoms are produced by X. \textit{Causal connection:} no. \textit{Defeater:} genuine. \textit{Reflective access:} yes. \textit{Threshold} $\theta = 2$.\par
\begin{center}\begin{tabular}{@{}p{0.43\linewidth}rcccl@{}}
world & $d$ & close & $B$ & $p$ & tag \\ \midrule
actual world & 0 & yes & yes & yes & act \\
the patient has not contracted X & 1 & yes & yes & no & tm \\
the symptoms are produced by X & 2 & yes & yes & yes & ok \\
\end{tabular}\end{center}
\noindent\textit{Computed.} J yes, T yes, error distance $\delta = 1$, locally safe no, safe basis-wide no, Gettier case yes, knowledge no. Profile over $\theta = 0..4$: \texttt{Suuuu}. Class: \textbf{G-fl/d}.
 \textit{Source verdict:} not knowledge (presented as a Gettier case). Agreement: yes.
\par\medskip
\needspace{12\baselineskip}
\subsection*{ZHB \enspace The husband in the room}
\addcontentsline{toc}{subsection}{ZHB The husband in the room}
\noindent\textit{Source.} \citet{Zagzebski1994}.\quad \textit{Route:} perception.\quad \textit{Basis:} seeing a man she takes to be him.\par
\noindent\textit{Case.} Mary, whose eyesight is good, glances into the living room and takes the man she sees in a chair to be her husband. It is his brother, who looks very like him. Her husband is in fact in the room, out of her line of sight.\par
\noindent\textit{Belief} $p$: my husband is in the living room.
 \textit{False lemma:} the man in the chair is my husband. \textit{Causal connection:} no. \textit{Defeater:} genuine. \textit{Reflective access:} yes. \textit{Threshold} $\theta = 2$.\par
\begin{center}\begin{tabular}{@{}p{0.43\linewidth}rcccl@{}}
world & $d$ & close & $B$ & $p$ & tag \\ \midrule
actual world & 0 & yes & yes & yes & act \\
her husband is in the garden & 1 & yes & yes & no & tm \\
the man in the chair is her husband & 2 & yes & yes & yes & ok \\
\end{tabular}\end{center}
\noindent\textit{Computed.} J yes, T yes, error distance $\delta = 1$, locally safe no, safe basis-wide no, Gettier case yes, knowledge no. Profile over $\theta = 0..4$: \texttt{Suuuu}. Class: \textbf{G-int/d}.
 \textit{Source verdict:} not knowledge (presented as a Gettier case). Agreement: yes.
\par\medskip
\needspace{12\baselineskip}
\subsection*{PYR \enspace The pyromaniac}
\addcontentsline{toc}{subsection}{PYR The pyromaniac}
\noindent\textit{Source.} \citet{Skyrms1967}.\quad \textit{Route:} induction.\quad \textit{Basis:} induction from past strikings.\par
\noindent\textit{Case.} A man who has always found that matches of a certain brand light when struck believes, on that experience, that the match he is striking will light. This match has an impurity that would ordinarily prevent it from lighting; it lights because of an unusual physical occurrence at the moment of striking.\par
\noindent\textit{Belief} $p$: this match will light.
 \textit{False lemma:} this match is like the others. \textit{Causal connection:} no. \textit{Defeater:} genuine. \textit{Reflective access:} yes. \textit{Threshold} $\theta = 2$.\par
\begin{center}\begin{tabular}{@{}p{0.43\linewidth}rcccl@{}}
world & $d$ & close & $B$ & $p$ & tag \\ \midrule
actual world & 0 & yes & yes & yes & act \\
the unusual occurrence does not take place & 1 & yes & yes & no & tm \\
the match is free of the impurity & 3 & no & yes & yes & ok \\
\end{tabular}\end{center}
\noindent\textit{Computed.} J yes, T yes, error distance $\delta = 1$, locally safe no, safe basis-wide no, Gettier case yes, knowledge no. Profile over $\theta = 0..4$: \texttt{Suuuu}. Class: \textbf{G-fl/d}.
 \textit{Source verdict:} not knowledge (presented as a Gettier case). Agreement: yes.
\par\medskip
\needspace{12\baselineskip}
\subsection*{GRB \enspace Tom Grabit}
\addcontentsline{toc}{subsection}{GRB Tom Grabit}
\noindent\textit{Source.} \citet{LehrerPaxson1969}.\quad \textit{Route:} perception.\quad \textit{Basis:} seeing Tom take the book.\par
\noindent\textit{Case.} The subject sees Tom Grabit, whom she knows well, take a book from the library and conceal it, and believes that Tom stole the book. Unknown to her, Tom's mother has said that Tom has an identical twin who was in the library; the mother is deluded, and there is no twin.\par
\noindent\textit{Belief} $p$: Tom stole the book.
 \textit{False lemma:} none. \textit{Causal connection:} yes. \textit{Defeater:} misleading. \textit{Reflective access:} yes. \textit{Threshold} $\theta = 2$.\par
\begin{center}\begin{tabular}{@{}p{0.43\linewidth}rcccl@{}}
world & $d$ & close & $B$ & $p$ & tag \\ \midrule
actual world & 0 & yes & yes & yes & act \\
the mother's claim reaches her first & 1 & yes & no & yes & def \\
she sees Tom from another angle & 2 & yes & yes & yes & ok \\
Tom has a twin who was in the library & 4 & no & yes & no & env \\
\end{tabular}\end{center}
\noindent\textit{Computed.} J yes, T yes, error distance $\delta = 4$, locally safe yes, safe basis-wide yes, Gettier case no, knowledge yes. Profile over $\theta = 0..4$: \texttt{SSSSu}. Class: \textbf{K[md]}.
 \textit{Source verdict:} knowledge (presented against a fourth condition). Agreement: yes.
\par\medskip
\needspace{12\baselineskip}
\subsection*{HAR \enspace Assassination}
\addcontentsline{toc}{subsection}{HAR Assassination}
\noindent\textit{Source.} \citet{Harman1973}.\quad \textit{Route:} testimony.\quad \textit{Basis:} a report in a reliable newspaper.\par
\noindent\textit{Case.} Jill reads in a reliable newspaper that a political leader has been assassinated, and believes it. The report is true. The authorities then issue widespread false denials, which almost everyone hears and which Jill happens to miss.\par
\noindent\textit{Belief} $p$: the leader has been assassinated.
 \textit{False lemma:} none. \textit{Causal connection:} yes. \textit{Defeater:} misleading. \textit{Reflective access:} yes. \textit{Threshold} $\theta = 2$.\par
\begin{center}\begin{tabular}{@{}p{0.43\linewidth}rcccl@{}}
world & $d$ & close & $B$ & $p$ & tag \\ \midrule
actual world & 0 & yes & yes & yes & act \\
Jill hears the denials & 1 & yes & no & yes & def \\
she reads the report in another paper & 1 & yes & yes & yes & ok \\
\end{tabular}\end{center}
\noindent\textit{Computed.} J yes, T yes, error distance $\delta = \infty$, locally safe yes, safe basis-wide yes, Gettier case no, knowledge yes. Profile over $\theta = 0..4$: \texttt{SSSSS}. Class: \textbf{K[md]}.
 \textit{Source verdict:} not knowledge (presented against a fourth condition). Agreement: no.
 \textit{Dispute (dep):} Harman's verdict is that Jill does not know. The case has the structure of GRB, where the source's verdict is knowledge: a misleading defeater the subject does not possess.
\par\medskip
\needspace{12\baselineskip}
\subsection*{CLV \enspace Norman the clairvoyant}
\addcontentsline{toc}{subsection}{CLV Norman the clairvoyant}
\noindent\textit{Source.} \citet{BonJour1980}.\quad \textit{Route:} faculty.\quad \textit{Basis:} an unexplained clairvoyant faculty.\par
\noindent\textit{Case.} Norman is, in normal conditions, a completely reliable clairvoyant about a certain subject matter, but has no evidence for or against his having such a power. One day he comes to believe, for no reason he can give, that the President is in New York City. The belief is true and is produced by his power.\par
\noindent\textit{Belief} $p$: the President is in New York City.
 \textit{False lemma:} none. \textit{Causal connection:} yes. \textit{Defeater:} none. \textit{Reflective access:} no. \textit{Threshold} $\theta = 2$.\par
\begin{center}\begin{tabular}{@{}p{0.43\linewidth}rcccl@{}}
world & $d$ & close & $B$ & $p$ & tag \\ \midrule
actual world & 0 & yes & yes & yes & act \\
the President is in Washington & 1 & yes & no & no & tm \\
the President has just arrived & 2 & yes & yes & yes & ok \\
\end{tabular}\end{center}
\noindent\textit{Computed.} J no, T yes, error distance $\delta = \infty$, locally safe yes, safe basis-wide yes, Gettier case no, knowledge no. Profile over $\theta = 0..4$: \texttt{SSSSS}. Class: \textbf{U}.
 \textit{Source verdict:} not knowledge (presented against a fourth condition). Agreement: yes.
\par\medskip
\needspace{12\baselineskip}
\subsection*{TTP \enspace Truetemp}
\addcontentsline{toc}{subsection}{TTP Truetemp}
\noindent\textit{Source.} \citet{Lehrer1990}.\quad \textit{Route:} faculty.\quad \textit{Basis:} the output of an implanted device.\par
\noindent\textit{Case.} Mr Truetemp has had, without his knowledge, a device implanted in his head that measures the temperature accurately and causes him to believe what it measures. He believes that it is 104 degrees, and it is. He has no idea why he has such beliefs and has never checked them.\par
\noindent\textit{Belief} $p$: it is 104 degrees.
 \textit{False lemma:} none. \textit{Causal connection:} yes. \textit{Defeater:} none. \textit{Reflective access:} no. \textit{Threshold} $\theta = 2$.\par
\begin{center}\begin{tabular}{@{}p{0.43\linewidth}rcccl@{}}
world & $d$ & close & $B$ & $p$ & tag \\ \midrule
actual world & 0 & yes & yes & yes & act \\
it is 103 degrees & 1 & yes & no & no & tm \\
it is 104 degrees an hour later & 2 & yes & yes & yes & ok \\
\end{tabular}\end{center}
\noindent\textit{Computed.} J no, T yes, error distance $\delta = \infty$, locally safe yes, safe basis-wide yes, Gettier case no, knowledge no. Profile over $\theta = 0..4$: \texttt{SSSSS}. Class: \textbf{U}.
 \textit{Source verdict:} not knowledge (presented against a fourth condition). Agreement: yes.
\par\medskip
\needspace{12\baselineskip}
\subsection*{CHU \enspace Rubbish chute}
\addcontentsline{toc}{subsection}{CHU Rubbish chute}
\noindent\textit{Source.} \citet{Sosa1999}.\quad \textit{Route:} induction.\quad \textit{Basis:} induction from experience of the chute.\par
\noindent\textit{Case.} A resident of a high-rise building drops a bag of rubbish into the chute and believes, from long experience of the chute, that the bag is now in the basement.\par
\noindent\textit{Belief} $p$: the bag is in the basement.
 \textit{False lemma:} none. \textit{Causal connection:} yes. \textit{Defeater:} none. \textit{Reflective access:} yes. \textit{Threshold} $\theta = 2$.\par
\begin{center}\begin{tabular}{@{}p{0.43\linewidth}rcccl@{}}
world & $d$ & close & $B$ & $p$ & tag \\ \midrule
actual world & 0 & yes & yes & yes & act \\
she lets the bag go a moment later & 2 & yes & yes & yes & ok \\
the bag snags on the way down & 4 & no & yes & no & tm \\
\end{tabular}\end{center}
\noindent\textit{Computed.} J yes, T yes, error distance $\delta = 4$, locally safe yes, safe basis-wide yes, Gettier case no, knowledge yes. Profile over $\theta = 0..4$: \texttt{SSSSu}. Class: \textbf{K}.
 \textit{Source verdict:} knowledge (presented against a fourth condition). Agreement: yes.
\par\medskip
\needspace{12\baselineskip}
\subsection*{KAL \enspace Kaleidoscope}
\addcontentsline{toc}{subsection}{KAL Kaleidoscope}
\noindent\textit{Source.} \citet{Sosa2007}.\quad \textit{Route:} perception.\quad \textit{Basis:} visual appearance of color.\par
\noindent\textit{Case.} A perceiver looks at a surface that is red and lit by normal light, and believes that it is red. The scene is controlled by a jokester who could as easily have shown her a white surface lit by red light, which would have looked the same.\par
\noindent\textit{Belief} $p$: the surface is red.
 \textit{False lemma:} none. \textit{Causal connection:} yes. \textit{Defeater:} genuine. \textit{Reflective access:} no. \textit{Threshold} $\theta = 2$.\par
\begin{center}\begin{tabular}{@{}p{0.43\linewidth}rcccl@{}}
world & $d$ & close & $B$ & $p$ & tag \\ \midrule
actual world & 0 & yes & yes & yes & act \\
the jokester shows a white surface under red light & 1 & yes & yes & no & env \\
a different red surface is shown & 1 & yes & yes & yes & ok \\
\end{tabular}\end{center}
\noindent\textit{Computed.} J yes, T yes, error distance $\delta = 1$, locally safe no, safe basis-wide no, Gettier case yes, knowledge no. Profile over $\theta = 0..4$: \texttt{Suuuu}. Class: \textbf{G-env/c}.
 \textit{Source verdict:} knowledge (presented against a fourth condition). Agreement: no.
 \textit{Dispute (dep):} Sosa grants animal knowledge (apt belief) and denies reflective knowledge; the belief is unsafe, so on the analysis it is not knowledge.
\par\medskip
\needspace{12\baselineskip}
\subsection*{TMP \enspace Temp}
\addcontentsline{toc}{subsection}{TMP Temp}
\noindent\textit{Source.} \citet{Pritchard2012}.\quad \textit{Route:} instrument.\quad \textit{Basis:} reading the thermometer.\par
\noindent\textit{Case.} Temp forms beliefs about the temperature of a room by reading a thermometer which, unknown to him, is broken and shows readings at random. Someone hidden in the room, whose task is to make the readings correct, adjusts the temperature to match each reading as Temp consults it.\par
\noindent\textit{Belief} $p$: the room is at the temperature the thermometer shows.
 \textit{False lemma:} none. \textit{Causal connection:} no. \textit{Defeater:} none. \textit{Reflective access:} yes. \textit{Threshold} $\theta = 2$.\par
\begin{center}\begin{tabular}{@{}p{0.43\linewidth}rcccl@{}}
world & $d$ & close & $B$ & $p$ & tag \\ \midrule
actual world & 0 & yes & yes & yes & act \\
the thermometer shows another reading & 1 & yes & yes & yes & ok \\
the helper is absent & 3 & no & yes & no & tm \\
\end{tabular}\end{center}
\noindent\textit{Computed.} J yes, T yes, error distance $\delta = 3$, locally safe yes, safe basis-wide yes, Gettier case no, knowledge yes. Profile over $\theta = 0..4$: \texttt{SSSuu}. Class: \textbf{K[fx]}.
 \textit{Source verdict:} not knowledge (presented against a fourth condition). Agreement: no.
 \textit{Dispute (dep):} Pritchard's verdict is that Temp does not know. The belief is justified and safe; the truth is secured by another agent, as in MOR, where the source's verdict is knowledge.
\par\medskip
\needspace{12\baselineskip}
\subsection*{MOR \enspace Morris in Chicago}
\addcontentsline{toc}{subsection}{MOR Morris in Chicago}
\noindent\textit{Source.} \citet{Lackey2007}.\quad \textit{Route:} testimony.\quad \textit{Basis:} testimony of a passer-by.\par
\noindent\textit{Case.} Morris, just arrived in Chicago, asks the first passer-by for directions to the Sears Tower, and believes what he is told. The passer-by is a local and the directions are correct.\par
\noindent\textit{Belief} $p$: the tower is two blocks east.
 \textit{False lemma:} none. \textit{Causal connection:} yes. \textit{Defeater:} none. \textit{Reflective access:} yes. \textit{Threshold} $\theta = 2$.\par
\begin{center}\begin{tabular}{@{}p{0.43\linewidth}rcccl@{}}
world & $d$ & close & $B$ & $p$ & tag \\ \midrule
actual world & 0 & yes & yes & yes & act \\
he asks a different local & 1 & yes & yes & yes & ok \\
the passer-by is another visitor who guesses & 3 & no & yes & no & src \\
\end{tabular}\end{center}
\noindent\textit{Computed.} J yes, T yes, error distance $\delta = 3$, locally safe yes, safe basis-wide yes, Gettier case no, knowledge yes. Profile over $\theta = 0..4$: \texttt{SSSuu}. Class: \textbf{K}.
 \textit{Source verdict:} knowledge (presented against a fourth condition). Agreement: yes.
\par\medskip
\needspace{12\baselineskip}
\subsection*{LOT \enspace Lottery}
\addcontentsline{toc}{subsection}{LOT Lottery}
\noindent\textit{Source.} \citet{Hawthorne2004}.\quad \textit{Route:} statistics.\quad \textit{Basis:} the odds.\par
\noindent\textit{Case.} A holder of one ticket in a large and fair lottery believes, on the basis of the odds alone, that her ticket will lose. It does.\par
\noindent\textit{Belief} $p$: my ticket will lose.
 \textit{False lemma:} none. \textit{Causal connection:} no. \textit{Defeater:} none. \textit{Reflective access:} yes. \textit{Threshold} $\theta = 2$.\par
\begin{center}\begin{tabular}{@{}p{0.43\linewidth}rcccl@{}}
world & $d$ & close & $B$ & $p$ & tag \\ \midrule
actual world & 0 & yes & yes & yes & act \\
her ticket is drawn & 1 & yes & yes & no & draw \\
another losing ticket is drawn & 1 & yes & yes & yes & ok \\
\end{tabular}\end{center}
\noindent\textit{Computed.} J yes, T yes, error distance $\delta = 1$, locally safe no, safe basis-wide no, Gettier case yes, knowledge no. Profile over $\theta = 0..4$: \texttt{Suuuu}. Class: \textbf{G-lot/d}.
 \textit{Source verdict:} not knowledge (other). Agreement: yes.
\par\medskip
\needspace{12\baselineskip}
\subsection*{LTN \enspace Lottery, newspaper}
\addcontentsline{toc}{subsection}{LTN Lottery, newspaper}
\noindent\textit{Source.} \citet{Hawthorne2004}.\quad \textit{Route:} testimony.\quad \textit{Basis:} the published result.\par
\noindent\textit{Case.} The same holder believes that her ticket has lost on reading the published result in a reliable newspaper.\par
\noindent\textit{Belief} $p$: my ticket has lost.
 \textit{False lemma:} none. \textit{Causal connection:} yes. \textit{Defeater:} none. \textit{Reflective access:} yes. \textit{Threshold} $\theta = 2$.\par
\begin{center}\begin{tabular}{@{}p{0.43\linewidth}rcccl@{}}
world & $d$ & close & $B$ & $p$ & tag \\ \midrule
actual world & 0 & yes & yes & yes & act \\
her ticket is drawn and reported & 1 & yes & no & no & draw \\
another ticket is drawn and reported & 1 & yes & yes & yes & ok \\
the paper misprints the result & 3 & no & yes & no & src \\
\end{tabular}\end{center}
\noindent\textit{Computed.} J yes, T yes, error distance $\delta = 3$, locally safe yes, safe basis-wide yes, Gettier case no, knowledge yes. Profile over $\theta = 0..4$: \texttt{SSSuu}. Class: \textbf{K}.
 \textit{Source verdict:} knowledge (other). Agreement: yes.
\par\medskip
\needspace{12\baselineskip}
\subsection*{MAG \enspace Magoo's tree}
\addcontentsline{toc}{subsection}{MAG Magoo's tree}
\noindent\textit{Source.} \citet{Williamson2000}.\quad \textit{Route:} perception.\quad \textit{Basis:} visual estimate of height.\par
\noindent\textit{Case.} Mr Magoo looks at a distant tree and can judge its height only roughly. Its height is $i+1$ inches; he believes that it is not $i$ inches tall, though a tree of $i$ inches would look no different to him.\par
\noindent\textit{Belief} $p$: the tree is not $i$ inches tall.
 \textit{False lemma:} none. \textit{Causal connection:} yes. \textit{Defeater:} none. \textit{Reflective access:} yes. \textit{Threshold} $\theta = 2$.\par
\begin{center}\begin{tabular}{@{}p{0.43\linewidth}rcccl@{}}
world & $d$ & close & $B$ & $p$ & tag \\ \midrule
actual world & 0 & yes & yes & yes & act \\
the tree is $i$ inches tall & 1 & yes & yes & no & marg \\
the tree is $i+2$ inches tall & 1 & yes & yes & yes & marg \\
\end{tabular}\end{center}
\noindent\textit{Computed.} J yes, T yes, error distance $\delta = 1$, locally safe no, safe basis-wide no, Gettier case yes, knowledge no. Profile over $\theta = 0..4$: \texttt{Suuuu}. Class: \textbf{G-mrg/c}.
 \textit{Source verdict:} not knowledge (other). Agreement: yes.
\par\medskip
\needspace{12\baselineskip}
\subsection*{MGW \enspace Magoo, wide margin}
\addcontentsline{toc}{subsection}{MGW Magoo, wide margin}
\noindent\textit{Source.} \citet{Williamson2000}.\quad \textit{Route:} perception.\quad \textit{Basis:} visual estimate of height.\par
\noindent\textit{Case.} Mr Magoo, looking at the same tree, believes that it is not $i-50$ inches tall, a height that would look quite different to him.\par
\noindent\textit{Belief} $p$: the tree is not $i-50$ inches tall.
 \textit{False lemma:} none. \textit{Causal connection:} yes. \textit{Defeater:} none. \textit{Reflective access:} yes. \textit{Threshold} $\theta = 2$.\par
\begin{center}\begin{tabular}{@{}p{0.43\linewidth}rcccl@{}}
world & $d$ & close & $B$ & $p$ & tag \\ \midrule
actual world & 0 & yes & yes & yes & act \\
the tree is $i$ inches tall & 1 & yes & yes & yes & marg \\
the tree is $i-50$ inches tall & 4 & no & no & no & marg \\
\end{tabular}\end{center}
\noindent\textit{Computed.} J yes, T yes, error distance $\delta = \infty$, locally safe yes, safe basis-wide yes, Gettier case no, knowledge yes. Profile over $\theta = 0..4$: \texttt{SSSSS}. Class: \textbf{K}.
 \textit{Source verdict:} knowledge (other). Agreement: yes.
\par\medskip
\needspace{12\baselineskip}
\subsection*{ZBP \enspace Zebras}
\addcontentsline{toc}{subsection}{ZBP Zebras}
\noindent\textit{Source.} \citet{Dretske1970}.\quad \textit{Route:} perception.\quad \textit{Basis:} seeing zebra-like animals in the zebra pen.\par
\noindent\textit{Case.} A visitor to a zoo sees animals in a pen marked as the zebra enclosure, and believes that they are zebras.\par
\noindent\textit{Belief} $p$: the animals are zebras.
 \textit{Weaker} $q$: the animals are not mules cleverly disguised as zebras.
 \textit{False lemma:} none. \textit{Causal connection:} yes. \textit{Defeater:} none. \textit{Reflective access:} yes. \textit{Threshold} $\theta = 2$.\par
\begin{center}\begin{tabular}{@{}p{0.43\linewidth}rccclc@{}}
world & $d$ & close & $B$ & $p$ & tag & $q$ \\ \midrule
actual world & 0 & yes & yes & yes & act & yes \\
other zebras are in the pen & 1 & yes & yes & yes & ok & yes \\
the animals are disguised mules & 4 & no & yes & no & env & no \\
\end{tabular}\end{center}
\noindent\textit{Computed.} J yes, T yes, error distance $\delta = 4$, locally safe yes, safe basis-wide yes, Gettier case no, knowledge yes. Profile over $\theta = 0..4$: \texttt{SSSSu}. Class: \textbf{K}.
 \textit{Source verdict:} knowledge (other). Agreement: yes.
\par\medskip
\needspace{12\baselineskip}
\subsection*{ZBM \enspace Not disguised mules}
\addcontentsline{toc}{subsection}{ZBM Not disguised mules}
\noindent\textit{Source.} \citet{Dretske1970}.\quad \textit{Route:} deduction.\quad \textit{Basis:} deduction from the perceptual belief.\par
\noindent\textit{Case.} The same visitor infers from her belief that the animals are zebras that they are not mules cleverly disguised to look like zebras.\par
\noindent\textit{Belief} $p$: the animals are not cleverly disguised mules.
 \textit{False lemma:} none. \textit{Causal connection:} yes. \textit{Defeater:} none. \textit{Reflective access:} yes. \textit{Threshold} $\theta = 2$.\par
\begin{center}\begin{tabular}{@{}p{0.43\linewidth}rcccl@{}}
world & $d$ & close & $B$ & $p$ & tag \\ \midrule
actual world & 0 & yes & yes & yes & act \\
other zebras are in the pen & 1 & yes & yes & yes & ok \\
the animals are disguised mules & 4 & no & yes & no & env \\
\end{tabular}\end{center}
\noindent\textit{Computed.} J yes, T yes, error distance $\delta = 4$, locally safe yes, safe basis-wide yes, Gettier case no, knowledge yes. Profile over $\theta = 0..4$: \texttt{SSSSu}. Class: \textbf{K}.
 \textit{Source verdict:} not knowledge (other). Agreement: no.
 \textit{Dispute (dep):} Dretske denies knowledge, since the belief is insensitive; closure is defended by \citet{Williamson2000} and \citet{Hawthorne2004}, and by EPI-20.
\par\medskip
\needspace{12\baselineskip}
\subsection*{GRN \enspace The grandmother}
\addcontentsline{toc}{subsection}{GRN The grandmother}
\noindent\textit{Source.} \citet{Nozick1981}.\quad \textit{Route:} perception.\quad \textit{Basis:} seeing him.\par
\noindent\textit{Case.} A grandmother sees her grandson on a visit and believes that he is well. Had he been ill, the family would have kept it from her and told her that he was well.\par
\noindent\textit{Belief} $p$: my grandson is well.
 \textit{False lemma:} none. \textit{Causal connection:} yes. \textit{Defeater:} none. \textit{Reflective access:} yes. \textit{Threshold} $\theta = 2$.\par
\begin{center}\begin{tabular}{@{}p{0.43\linewidth}rcccl@{}}
world & $d$ & close & $B$ & $p$ & tag \\ \midrule
actual world & 0 & yes & yes & yes & act \\
he visits an hour later & 1 & yes & yes & yes & ok \\
he is ill and she is told he is well & 2 & yes & no & no & basis \\
\end{tabular}\end{center}
\noindent\textit{Computed.} J yes, T yes, error distance $\delta = \infty$, locally safe yes, safe basis-wide yes, Gettier case no, knowledge yes. Profile over $\theta = 0..4$: \texttt{SSSSS}. Class: \textbf{K[ev]}.
 \textit{Source verdict:} knowledge (presented against a fourth condition). Agreement: yes.
\par\medskip
\needspace{12\baselineskip}
\subsection*{HLF \enspace Halloween, fine basis}
\addcontentsline{toc}{subsection}{HLF Halloween, fine basis}
\noindent\textit{Source.} \citet{Comesana2005}.\quad \textit{Route:} testimony.\quad \textit{Basis:} directions from a truthful director.\par
\noindent\textit{Case.} A guest asks the way to a party of someone posted to give directions, who would, in circumstances that could easily have obtained, have given misleading ones. The guest receives correct directions and believes them. Here the basis is individuated to include the director's actual disposition to tell the truth.\par
\noindent\textit{Belief} $p$: the party is at that house.
 \textit{False lemma:} none. \textit{Causal connection:} yes. \textit{Defeater:} none. \textit{Reflective access:} yes. \textit{Threshold} $\theta = 2$.\par
\begin{center}\begin{tabular}{@{}p{0.43\linewidth}rcccl@{}}
world & $d$ & close & $B$ & $p$ & tag \\ \midrule
actual world & 0 & yes & yes & yes & act \\
the circumstances obtain; directions misleading & 1 & yes & no & no & src \\
he asks a moment later & 1 & yes & yes & yes & ok \\
\end{tabular}\end{center}
\noindent\textit{Computed.} J yes, T yes, error distance $\delta = \infty$, locally safe yes, safe basis-wide yes, Gettier case no, knowledge yes. Profile over $\theta = 0..4$: \texttt{SSSSS}. Class: \textbf{K}.
 \textit{Source verdict:} knowledge (presented against a fourth condition). Agreement: yes.
 \textit{Dispute (bas):} Basis individuated finely: the verdict agrees with Comesa\~na's.
\par\medskip
\needspace{12\baselineskip}
\subsection*{HLC \enspace Halloween, coarse basis}
\addcontentsline{toc}{subsection}{HLC Halloween, coarse basis}
\noindent\textit{Source.} \citet{Comesana2005}.\quad \textit{Route:} testimony.\quad \textit{Basis:} directions from the person posted.\par
\noindent\textit{Case.} The same case, with the basis individuated coarsely as directions from the person posted to give them, whatever his disposition.\par
\noindent\textit{Belief} $p$: the party is at that house.
 \textit{False lemma:} none. \textit{Causal connection:} yes. \textit{Defeater:} genuine. \textit{Reflective access:} yes. \textit{Threshold} $\theta = 2$.\par
\begin{center}\begin{tabular}{@{}p{0.43\linewidth}rcccl@{}}
world & $d$ & close & $B$ & $p$ & tag \\ \midrule
actual world & 0 & yes & yes & yes & act \\
the circumstances obtain; directions misleading & 1 & yes & yes & no & src \\
he asks a moment later & 1 & yes & yes & yes & ok \\
\end{tabular}\end{center}
\noindent\textit{Computed.} J yes, T yes, error distance $\delta = 1$, locally safe no, safe basis-wide no, Gettier case yes, knowledge no. Profile over $\theta = 0..4$: \texttt{Suuuu}. Class: \textbf{G-tst/c}.
 \textit{Source verdict:} knowledge (presented against a fourth condition). Agreement: no.
 \textit{Dispute (bas):} Basis individuated coarsely: a Gettier case, against Comesa\~na's verdict.
\par\medskip
\needspace{12\baselineskip}
\subsection*{ATF \enspace Atomic clock, fine basis}
\addcontentsline{toc}{subsection}{ATF Atomic clock, fine basis}
\noindent\textit{Source.} \citet{Bogardus2014}.\quad \textit{Route:} instrument.\quad \textit{Basis:} reading a functioning clock.\par
\noindent\textit{Case.} A subject reads the time from a highly accurate clock. A nearby radioactive source could easily have made it malfunction, but did not. Here the basis is reading a functioning clock.\par
\noindent\textit{Belief} $p$: it is the time the clock shows.
 \textit{False lemma:} none. \textit{Causal connection:} yes. \textit{Defeater:} none. \textit{Reflective access:} yes. \textit{Threshold} $\theta = 2$.\par
\begin{center}\begin{tabular}{@{}p{0.43\linewidth}rcccl@{}}
world & $d$ & close & $B$ & $p$ & tag \\ \midrule
actual world & 0 & yes & yes & yes & act \\
the clock malfunctions & 1 & yes & no & no & env \\
she reads it a minute later & 1 & yes & yes & yes & ok \\
\end{tabular}\end{center}
\noindent\textit{Computed.} J yes, T yes, error distance $\delta = \infty$, locally safe yes, safe basis-wide yes, Gettier case no, knowledge yes. Profile over $\theta = 0..4$: \texttt{SSSSS}. Class: \textbf{K}.
 \textit{Source verdict:} knowledge (presented against a fourth condition). Agreement: yes.
 \textit{Dispute (bas):} Basis individuated finely: the verdict agrees with Bogardus's.
\par\medskip
\needspace{12\baselineskip}
\subsection*{ATC \enspace Atomic clock, coarse basis}
\addcontentsline{toc}{subsection}{ATC Atomic clock, coarse basis}
\noindent\textit{Source.} \citet{Bogardus2014}.\quad \textit{Route:} instrument.\quad \textit{Basis:} reading that clock.\par
\noindent\textit{Case.} The same case, with the basis individuated as reading that clock, whether or not it functions.\par
\noindent\textit{Belief} $p$: it is the time the clock shows.
 \textit{False lemma:} none. \textit{Causal connection:} yes. \textit{Defeater:} genuine. \textit{Reflective access:} yes. \textit{Threshold} $\theta = 2$.\par
\begin{center}\begin{tabular}{@{}p{0.43\linewidth}rcccl@{}}
world & $d$ & close & $B$ & $p$ & tag \\ \midrule
actual world & 0 & yes & yes & yes & act \\
the clock malfunctions & 1 & yes & yes & no & env \\
she reads it a minute later & 1 & yes & yes & yes & ok \\
\end{tabular}\end{center}
\noindent\textit{Computed.} J yes, T yes, error distance $\delta = 1$, locally safe no, safe basis-wide no, Gettier case yes, knowledge no. Profile over $\theta = 0..4$: \texttt{Suuuu}. Class: \textbf{G-env/c}.
 \textit{Source verdict:} knowledge (presented against a fourth condition). Agreement: no.
 \textit{Dispute (bas):} Basis individuated coarsely: a Gettier case, against Bogardus's verdict.
\par\medskip
\needspace{12\baselineskip}
\subsection*{WAR \enspace Handouts}
\addcontentsline{toc}{subsection}{WAR Handouts}
\noindent\textit{Source.} \citet{Warfield2005}.\quad \textit{Route:} inference.\quad \textit{Basis:} inference from a careful count.\par
\noindent\textit{Case.} A speaker counts the audience at her talk, arrives at 53, and, having 100 handouts, concludes that she has enough. There are in fact 52 people present.\par
\noindent\textit{Belief} $p$: I have enough handouts.
 \textit{False lemma:} there are 53 people present. \textit{Causal connection:} yes. \textit{Defeater:} none. \textit{Reflective access:} yes. \textit{Threshold} $\theta = 2$.\par
\begin{center}\begin{tabular}{@{}p{0.43\linewidth}rcccl@{}}
world & $d$ & close & $B$ & $p$ & tag \\ \midrule
actual world & 0 & yes & yes & yes & act \\
she counts 54 & 1 & yes & yes & yes & marg \\
53 people are present & 1 & yes & yes & yes & ok \\
more than 100 people are present & 4 & no & no & no & tm \\
\end{tabular}\end{center}
\noindent\textit{Computed.} J yes, T yes, error distance $\delta = \infty$, locally safe yes, safe basis-wide yes, Gettier case no, knowledge yes. Profile over $\theta = 0..4$: \texttt{SSSSS}. Class: \textbf{K}.
 \textit{Source verdict:} knowledge (presented against a fourth condition). Agreement: yes.
 \textit{Dispute (int):} The verdict is generally accepted; whether the knowledge depends on the falsehood is disputed by \citet{Klein2008}.
\par\medskip
\needspace{12\baselineskip}
\subsection*{BKL \enspace Bank, low stakes}
\addcontentsline{toc}{subsection}{BKL Bank, low stakes}
\noindent\textit{Source.} \citet{DeRose1992}.\quad \textit{Route:} induction.\quad \textit{Basis:} memory of a recent Saturday.\par
\noindent\textit{Case.} On a Friday, deciding whether to deposit a check now or on Saturday, a person believes that the bank will be open on Saturday because she was there on a Saturday two weeks ago. Little depends on it, and no one has raised the possibility that the hours have changed.\par
\noindent\textit{Belief} $p$: the bank will be open on Saturday.
 \textit{False lemma:} none. \textit{Causal connection:} yes. \textit{Defeater:} none. \textit{Reflective access:} yes. \textit{Threshold} $\theta = 2$.\par
\begin{center}\begin{tabular}{@{}p{0.43\linewidth}rcccl@{}}
world & $d$ & close & $B$ & $p$ & tag \\ \midrule
actual world & 0 & yes & yes & yes & act \\
the bank has changed its hours & 3 & no & yes & no & tm \\
she was there three weeks ago & 1 & yes & yes & yes & ok \\
\end{tabular}\end{center}
\noindent\textit{Computed.} J yes, T yes, error distance $\delta = 3$, locally safe yes, safe basis-wide yes, Gettier case no, knowledge yes. Profile over $\theta = 0..4$: \texttt{SSSuu}. Class: \textbf{K}.
 \textit{Source verdict:} knowledge (other). Agreement: yes.
\par\medskip
\needspace{12\baselineskip}
\subsection*{BKH \enspace Bank, high stakes}
\addcontentsline{toc}{subsection}{BKH Bank, high stakes}
\noindent\textit{Source.} \citet{DeRose1992}.\quad \textit{Route:} induction.\quad \textit{Basis:} memory of a recent Saturday.\par
\noindent\textit{Case.} The same person and evidence; a great deal depends on the deposit, and the possibility that the bank has changed its hours has been raised. The attended possibility is close in this context.\par
\noindent\textit{Belief} $p$: the bank will be open on Saturday.
 \textit{False lemma:} none. \textit{Causal connection:} yes. \textit{Defeater:} none. \textit{Reflective access:} yes. \textit{Threshold} $\theta = 3$.\par
\begin{center}\begin{tabular}{@{}p{0.43\linewidth}rcccl@{}}
world & $d$ & close & $B$ & $p$ & tag \\ \midrule
actual world & 0 & yes & yes & yes & act \\
the bank has changed its hours & 3 & yes & yes & no & tm \\
she was there three weeks ago & 1 & yes & yes & yes & ok \\
\end{tabular}\end{center}
\noindent\textit{Computed.} J yes, T yes, error distance $\delta = 3$, locally safe no, safe basis-wide no, Gettier case yes, knowledge no. Profile over $\theta = 0..4$: \texttt{SSSuu}. Class: \textbf{G-nfl/c}.
 \textit{Source verdict:} not knowledge (other). Agreement: yes.
\par\medskip
\needspace{12\baselineskip}
\subsection*{BIV \enspace Not a brain in a vat}
\addcontentsline{toc}{subsection}{BIV Not a brain in a vat}
\noindent\textit{Source.} \citet{Sosa1999}.\quad \textit{Route:} perception.\quad \textit{Basis:} ordinary perception and what it implies.\par
\noindent\textit{Case.} A person sitting at her desk believes that she is not a handless brain in a vat being fed experiences as of sitting at a desk. Were she such a brain, she would believe the same.\par
\noindent\textit{Belief} $p$: I am not a handless brain in a vat.
 \textit{False lemma:} none. \textit{Causal connection:} yes. \textit{Defeater:} none. \textit{Reflective access:} yes. \textit{Threshold} $\theta = 2$.\par
\begin{center}\begin{tabular}{@{}p{0.43\linewidth}rcccl@{}}
world & $d$ & close & $B$ & $p$ & tag \\ \midrule
actual world & 0 & yes & yes & yes & act \\
she sits in another chair & 1 & yes & yes & yes & ok \\
she is a brain in a vat & 4 & no & yes & no & env \\
\end{tabular}\end{center}
\noindent\textit{Computed.} J yes, T yes, error distance $\delta = 4$, locally safe yes, safe basis-wide yes, Gettier case no, knowledge yes. Profile over $\theta = 0..4$: \texttt{SSSSu}. Class: \textbf{K}.
 \textit{Source verdict:} knowledge (presented against a fourth condition). Agreement: yes.
\par\medskip
\needspace{12\baselineskip}
\subsection*{JJM \enspace Jesse James}
\addcontentsline{toc}{subsection}{JJM Jesse James}
\noindent\textit{Source.} \citet{Nozick1981}.\quad \textit{Route:} perception.\quad \textit{Basis:} seeing the robber's face.\par
\noindent\textit{Case.} A bystander watches a masked man escape from a bank robbery. The mask slips, she sees the robber's face, recognizes Jesse James, and believes that Jesse James is robbing the bank. Had the mask not slipped, she would have formed no belief about who the robber was.\par
\noindent\textit{Belief} $p$: Jesse James is robbing the bank.
 \textit{False lemma:} none. \textit{Causal connection:} yes. \textit{Defeater:} none. \textit{Reflective access:} yes. \textit{Threshold} $\theta = 2$.\par
\begin{center}\begin{tabular}{@{}p{0.43\linewidth}rcccl@{}}
world & $d$ & close & $B$ & $p$ & tag \\ \midrule
actual world & 0 & yes & yes & yes & act \\
the mask does not slip & 1 & yes & no & yes & basis \\
the robber is a look-alike of Jesse James & 4 & no & yes & no & env \\
\end{tabular}\end{center}
\noindent\textit{Computed.} J yes, T yes, error distance $\delta = 4$, locally safe yes, safe basis-wide yes, Gettier case no, knowledge yes. Profile over $\theta = 0..4$: \texttt{SSSSu}. Class: \textbf{K[ev]}.
 \textit{Source verdict:} knowledge (presented against a fourth condition). Agreement: yes.
\par\medskip
\needspace{12\baselineskip}
\subsection*{THT \enspace The jury}
\addcontentsline{toc}{subsection}{THT The jury}
\noindent\textit{Source.} Plato, \emph{Theaetetus}, 201a--c.\quad \textit{Route:} testimony.\quad \textit{Basis:} the advocate's persuasion.\par
\noindent\textit{Case.} Jurors are brought by an advocate's skill to judge truly that a defendant committed a crime which only an eyewitness could know about. Their judgment is produced by persuasion, not by being taught the facts.\par
\noindent\textit{Belief} $p$: the defendant committed the crime.
 \textit{False lemma:} none. \textit{Causal connection:} no. \textit{Defeater:} none. \textit{Reflective access:} yes. \textit{Threshold} $\theta = 2$.\par
\begin{center}\begin{tabular}{@{}p{0.43\linewidth}rcccl@{}}
world & $d$ & close & $B$ & $p$ & tag \\ \midrule
actual world & 0 & yes & yes & yes & act \\
the defendant is innocent; the advocate argues as before & 1 & yes & yes & no & tm \\
the advocate argues more briefly & 1 & yes & yes & yes & ok \\
\end{tabular}\end{center}
\noindent\textit{Computed.} J no, T yes, error distance $\delta = 1$, locally safe no, safe basis-wide no, Gettier case no, knowledge no. Profile over $\theta = 0..4$: \texttt{Suuuu}. Class: \textbf{U}.
 \textit{Source verdict:} not knowledge (other). Agreement: yes.
\par\medskip
\needspace{12\baselineskip}
\subsection*{ICE \enspace Ice cubes}
\addcontentsline{toc}{subsection}{ICE Ice cubes}
\noindent\textit{Source.} \citet{Vogel1987}.\quad \textit{Route:} induction.\quad \textit{Basis:} induction from the heat and the time elapsed.\par
\noindent\textit{Case.} Hours ago a person left a glass of ice cubes outside on a hot summer day. Without looking, she believes that the ice has melted. Had the cubes not melted, because of some unusual intervention, she would still believe they had.\par
\noindent\textit{Belief} $p$: the ice cubes have melted.
 \textit{False lemma:} none. \textit{Causal connection:} yes. \textit{Defeater:} none. \textit{Reflective access:} yes. \textit{Threshold} $\theta = 2$.\par
\begin{center}\begin{tabular}{@{}p{0.43\linewidth}rcccl@{}}
world & $d$ & close & $B$ & $p$ & tag \\ \midrule
actual world & 0 & yes & yes & yes & act \\
she left the glass an hour earlier & 1 & yes & yes & yes & ok \\
someone has replaced them with fresh cubes & 4 & no & yes & no & tm \\
\end{tabular}\end{center}
\noindent\textit{Computed.} J yes, T yes, error distance $\delta = 4$, locally safe yes, safe basis-wide yes, Gettier case no, knowledge yes. Profile over $\theta = 0..4$: \texttt{SSSSu}. Class: \textbf{K}.
 \textit{Source verdict:} knowledge (presented against a fourth condition). Agreement: yes.
\par\medskip
\needspace{12\baselineskip}
\subsection*{SNC \enspace Faulty multiplication}
\addcontentsline{toc}{subsection}{SNC Faulty multiplication}
\noindent\textit{Source.} schematic.\quad \textit{Route:} calculation.\quad \textit{Basis:} the pupil's procedure.\par
\noindent\textit{Case.} A pupil computes $12 \times 13$ by a procedure that, through a systematic slip, gives wrong products for most nearby pairs but gives the right product, 156, for this one. She has every reason to trust the procedure.\par
\noindent\textit{Belief} $p$: $12 \times 13 = 156$.
 \textit{False lemma:} none. \textit{Causal connection:} no. \textit{Defeater:} none. \textit{Reflective access:} yes. \textit{Threshold} $\theta = 2$.\par
\begin{center}\begin{tabular}{@{}p{0.43\linewidth}rccclc@{}}
world & $d$ & close & $B$ & $p$ & tag & alt \\ \midrule
actual world & 0 & yes & yes & yes & act & -- \\
she computes $12 	imes 14$ & 1 & yes & no & yes & ok & no \\
she computes $11 	imes 13$ & 1 & yes & no & yes & ok & no \\
\end{tabular}\end{center}
\noindent\textit{Computed.} J yes, T yes, error distance $\delta = 1$, locally safe yes, safe basis-wide no, Gettier case yes, knowledge no. Profile over $\theta = 0..4$: \texttt{Suuuu}. Class: \textbf{G-nec/d}.
\par\medskip
\needspace{12\baselineskip}
\subsection*{SNS \enspace Sound multiplication}
\addcontentsline{toc}{subsection}{SNS Sound multiplication}
\noindent\textit{Source.} schematic.\quad \textit{Route:} calculation.\quad \textit{Basis:} the standard procedure.\par
\noindent\textit{Case.} A pupil computes $12 \times 13 = 156$ by the standard procedure, which she has mastered.\par
\noindent\textit{Belief} $p$: $12 \times 13 = 156$.
 \textit{False lemma:} none. \textit{Causal connection:} no. \textit{Defeater:} none. \textit{Reflective access:} yes. \textit{Threshold} $\theta = 2$.\par
\begin{center}\begin{tabular}{@{}p{0.43\linewidth}rccclc@{}}
world & $d$ & close & $B$ & $p$ & tag & alt \\ \midrule
actual world & 0 & yes & yes & yes & act & -- \\
she computes $12 	imes 14$ & 1 & yes & no & yes & ok & yes \\
she computes $11 	imes 13$ & 1 & yes & no & yes & ok & yes \\
\end{tabular}\end{center}
\noindent\textit{Computed.} J yes, T yes, error distance $\delta = \infty$, locally safe yes, safe basis-wide yes, Gettier case no, knowledge yes. Profile over $\theta = 0..4$: \texttt{SSSSS}. Class: \textbf{K}.
\par\medskip
\needspace{12\baselineskip}
\subsection*{SEV \enspace Witness at the window}
\addcontentsline{toc}{subsection}{SEV Witness at the window}
\noindent\textit{Source.} schematic.\quad \textit{Route:} perception.\quad \textit{Basis:} seeing the robbery.\par
\noindent\textit{Case.} A passer-by happens to look out of a window at the moment a robbery takes place across the street, and believes that the man in the green coat is the robber.\par
\noindent\textit{Belief} $p$: the man in the green coat is the robber.
 \textit{False lemma:} none. \textit{Causal connection:} yes. \textit{Defeater:} none. \textit{Reflective access:} yes. \textit{Threshold} $\theta = 2$.\par
\begin{center}\begin{tabular}{@{}p{0.43\linewidth}rcccl@{}}
world & $d$ & close & $B$ & $p$ & tag \\ \midrule
actual world & 0 & yes & yes & yes & act \\
she looks a moment later and sees nothing & 1 & yes & no & yes & basis \\
a look-alike in a green coat commits it & 4 & no & yes & no & env \\
\end{tabular}\end{center}
\noindent\textit{Computed.} J yes, T yes, error distance $\delta = 4$, locally safe yes, safe basis-wide yes, Gettier case no, knowledge yes. Profile over $\theta = 0..4$: \texttt{SSSSu}. Class: \textbf{K[ev]}.
\par\medskip
\needspace{12\baselineskip}
\subsection*{SCT \enspace Winner told}
\addcontentsline{toc}{subsection}{SCT Winner told}
\noindent\textit{Source.} schematic.\quad \textit{Route:} testimony.\quad \textit{Basis:} the official announcement.\par
\noindent\textit{Case.} A ticket holder learns from the official announcement that her ticket has won. It is a matter of chance that it won.\par
\noindent\textit{Belief} $p$: my ticket has won.
 \textit{False lemma:} none. \textit{Causal connection:} yes. \textit{Defeater:} none. \textit{Reflective access:} yes. \textit{Threshold} $\theta = 2$.\par
\begin{center}\begin{tabular}{@{}p{0.43\linewidth}rcccl@{}}
world & $d$ & close & $B$ & $p$ & tag \\ \midrule
actual world & 0 & yes & yes & yes & act \\
another ticket wins and is announced & 1 & yes & no & no & draw \\
the announcement is misprinted & 4 & no & yes & no & src \\
\end{tabular}\end{center}
\noindent\textit{Computed.} J yes, T yes, error distance $\delta = 4$, locally safe yes, safe basis-wide yes, Gettier case no, knowledge yes. Profile over $\theta = 0..4$: \texttt{SSSSu}. Class: \textbf{K[ct]}.
\par\medskip
\needspace{12\baselineskip}
\subsection*{SRF \enspace Unreflective birder}
\addcontentsline{toc}{subsection}{SRF Unreflective birder}
\noindent\textit{Source.} schematic.\quad \textit{Route:} perception.\quad \textit{Basis:} recognitional skill.\par
\noindent\textit{Case.} A walker identifies a bird as a wren by sight, using a skill whose reliability she has never checked and about which she has no beliefs; it looks to her as a wren looks, and it is one.\par
\noindent\textit{Belief} $p$: that bird is a wren.
 \textit{False lemma:} none. \textit{Causal connection:} yes. \textit{Defeater:} none. \textit{Reflective access:} no. \textit{Threshold} $\theta = 2$.\par
\begin{center}\begin{tabular}{@{}p{0.43\linewidth}rcccl@{}}
world & $d$ & close & $B$ & $p$ & tag \\ \midrule
actual world & 0 & yes & yes & yes & act \\
a different wren & 1 & yes & yes & yes & ok \\
a bird of a similar species & 3 & no & no & no & env \\
\end{tabular}\end{center}
\noindent\textit{Computed.} J yes, T yes, error distance $\delta = \infty$, locally safe yes, safe basis-wide yes, Gettier case no, knowledge yes. Profile over $\theta = 0..4$: \texttt{SSSSS}. Class: \textbf{K[rf]}.
\par\medskip
\needspace{12\baselineskip}
\subsection*{STL \enspace Mistaken liar}
\addcontentsline{toc}{subsection}{STL Mistaken liar}
\noindent\textit{Source.} schematic.\quad \textit{Route:} testimony.\quad \textit{Basis:} the liar's testimony.\par
\noindent\textit{Case.} A habitual liar, intending to deceive, tells a hearer what he believes to be false. Through his own mistake, it is true. The hearer has no reason to distrust him.\par
\noindent\textit{Belief} $p$: the train leaves at noon.
 \textit{False lemma:} none. \textit{Causal connection:} no. \textit{Defeater:} genuine. \textit{Reflective access:} yes. \textit{Threshold} $\theta = 2$.\par
\begin{center}\begin{tabular}{@{}p{0.43\linewidth}rcccl@{}}
world & $d$ & close & $B$ & $p$ & tag \\ \midrule
actual world & 0 & yes & yes & yes & act \\
the liar's belief is correct; he says noon & 1 & yes & yes & no & src \\
he lies about another train & 2 & yes & no & yes & basis \\
\end{tabular}\end{center}
\noindent\textit{Computed.} J yes, T yes, error distance $\delta = 1$, locally safe no, safe basis-wide no, Gettier case yes, knowledge no. Profile over $\theta = 0..4$: \texttt{Suuuu}. Class: \textbf{G-tst/d}.
\par\medskip


\section{Systematic variants}\label{app:b}

\paragraph{What is computed.} For each case of Appendix~A, the program applies each of the eleven operations of Section~\ref{s:cat} that is applicable: \texttt{nfl} where the basis relies on a false lemma, \texttt{fl} where the route is inferential and it does not; \texttt{sec} where the model has an error world tagged \texttt{tm}, \texttt{draw} or \texttt{time}, setting $B = 0$ at those error worlds and marking the truth-maker as connected; \texttt{nor} where it has an error world tagged \texttt{env} or \texttt{src}, removing those worlds; \texttt{hos} where the route is perception, instrument or testimony and there is no close error world of those tags, adding a look-alike at distance $1$; \texttt{rem} where the belief is true and unsafe, moving every error world to distance $4$; \texttt{fal} where the belief is true, making it false at the actual world; \texttt{unj} where it is justified, removing justification; \texttt{rel} and \texttt{lie} where it is true, making the belief a hearer's, on sincere testimony, or on testimony from a speaker who could easily have asserted it falsely (an added error world at distance $1$, tagged \texttt{src}); and \texttt{ded} where the case specifies a weaker proposition $q$, replacing $p$ by $q$ at every world. The table gives the variant's justification, truth, error distance, safety and knowledge, its class and the parent's, and the prediction checked. The predictions are those of Propositions~\ref{p:lemma} (\texttt{nfl}, \texttt{fl}: safety unchanged), \ref{p:env} (\texttt{sec}, \texttt{nor}: safe when no other error remains; \texttt{hos}: unsafe), \ref{p:margin} (\texttt{rem}: safe below $4$), \ref{p:testimony} (\texttt{rel}: safety transmitted; \texttt{lie}: unsafe) and \ref{p:deduction} (\texttt{ded}: safe whenever the parent is), together with the definitional predictions for \texttt{fal} (not a Gettier case) and \texttt{unj} (not knowledge).

\paragraph{Summary.} The program generates \stnvariants{} variants. Every prediction holds; the program would halt otherwise. Each variant is also checked against Theorem~\ref{t:partition} and Lemma~\ref{l:delta}.

{\scriptsize\setlength{\tabcolsep}{2.5pt}
\begin{longtable}{@{}lp{0.36\linewidth}cccccll@{}}
\toprule
Id & Operation & J & T & $\delta$ & Safe & K & Class (parent $\to$ variant) & Prediction \\
\midrule\endhead
\bottomrule\endfoot
GE1.nfl & false lemma removed; same evidence, true intermediates & \textsf{+} & \textsf{+} & 1 & \textsf{--} & \textsf{--} & G-fl/d $\to$ G-nfl/d & safety unchanged \checkmark \\
GE1.sec & truth-maker secured: the basis now bears on it & \textsf{+} & \textsf{+} & $\infty$ & \textsf{+} & \textsf{+} & G-fl/d $\to$ K & safe \checkmark \\
GE1.rem & every error world moved to distance 4 & \textsf{+} & \textsf{+} & 4 & \textsf{+} & \textsf{+} & G-fl/d $\to$ K & safe below 4 \checkmark \\
GE1.fal & the belief made false at the actual world & \textsf{+} & \textsf{--} & 0 & \textsf{--} & \textsf{--} & G-fl/d $\to$ F & not Gettier \checkmark \\
GE1.unj & justification removed & \textsf{--} & \textsf{+} & 1 & \textsf{--} & \textsf{--} & G-fl/d $\to$ U & not knowledge \checkmark \\
GE1.rel & a hearer believes on the subject's sincere testimony & \textsf{+} & \textsf{+} & 1 & \textsf{--} & \textsf{--} & G-fl/d $\to$ G-int/d & safety transmitted \checkmark \\
GE1.lie & relay from a speaker who could easily have lied & \textsf{+} & \textsf{+} & 1 & \textsf{--} & \textsf{--} & G-fl/d $\to$ G-tst/d & unsafe \checkmark \\
GE1.ded & competent deduction of the weaker $q$ & \textsf{+} & \textsf{+} & $\infty$ & \textsf{+} & \textsf{+} & G-fl/d $\to$ K & safe \checkmark \\
\midrule
GE2.nfl & false lemma removed; same evidence, true intermediates & \textsf{+} & \textsf{+} & 1 & \textsf{--} & \textsf{--} & G-fl/d $\to$ G-nfl/d & safety unchanged \checkmark \\
GE2.sec & truth-maker secured: the basis now bears on it & \textsf{+} & \textsf{+} & $\infty$ & \textsf{+} & \textsf{+} & G-fl/d $\to$ K & safe \checkmark \\
GE2.rem & every error world moved to distance 4 & \textsf{+} & \textsf{+} & 4 & \textsf{+} & \textsf{+} & G-fl/d $\to$ K & safe below 4 \checkmark \\
GE2.fal & the belief made false at the actual world & \textsf{+} & \textsf{--} & 0 & \textsf{--} & \textsf{--} & G-fl/d $\to$ F & not Gettier \checkmark \\
GE2.unj & justification removed & \textsf{--} & \textsf{+} & 1 & \textsf{--} & \textsf{--} & G-fl/d $\to$ U & not knowledge \checkmark \\
GE2.rel & a hearer believes on the subject's sincere testimony & \textsf{+} & \textsf{+} & 1 & \textsf{--} & \textsf{--} & G-fl/d $\to$ G-int/d & safety transmitted \checkmark \\
GE2.lie & relay from a speaker who could easily have lied & \textsf{+} & \textsf{+} & 1 & \textsf{--} & \textsf{--} & G-fl/d $\to$ G-tst/d & unsafe \checkmark \\
\midrule
NOG.nfl & false lemma removed; same evidence, true intermediates & \textsf{+} & \textsf{+} & 1 & \textsf{--} & \textsf{--} & G-fl/d $\to$ G-nfl/d & safety unchanged \checkmark \\
NOG.sec & truth-maker secured: the basis now bears on it & \textsf{+} & \textsf{+} & $\infty$ & \textsf{+} & \textsf{+} & G-fl/d $\to$ K & safe \checkmark \\
NOG.rem & every error world moved to distance 4 & \textsf{+} & \textsf{+} & 4 & \textsf{+} & \textsf{+} & G-fl/d $\to$ K & safe below 4 \checkmark \\
NOG.fal & the belief made false at the actual world & \textsf{+} & \textsf{--} & 0 & \textsf{--} & \textsf{--} & G-fl/d $\to$ F & not Gettier \checkmark \\
NOG.unj & justification removed & \textsf{--} & \textsf{+} & 1 & \textsf{--} & \textsf{--} & G-fl/d $\to$ U & not knowledge \checkmark \\
NOG.rel & a hearer believes on the subject's sincere testimony & \textsf{+} & \textsf{+} & 1 & \textsf{--} & \textsf{--} & G-fl/d $\to$ G-int/d & safety transmitted \checkmark \\
NOG.lie & relay from a speaker who could easily have lied & \textsf{+} & \textsf{+} & 1 & \textsf{--} & \textsf{--} & G-fl/d $\to$ G-tst/d & unsafe \checkmark \\
NOG.ded & competent deduction of the weaker $q$ & \textsf{+} & \textsf{+} & $\infty$ & \textsf{+} & \textsf{+} & G-fl/d $\to$ K & safe \checkmark \\
\midrule
FEL.fl & a false intermediate belief added to the route & \textsf{+} & \textsf{+} & 1 & \textsf{--} & \textsf{--} & G-nfl/d $\to$ G-fl/d & safety unchanged \checkmark \\
FEL.sec & truth-maker secured: the basis now bears on it & \textsf{+} & \textsf{+} & $\infty$ & \textsf{+} & \textsf{+} & G-nfl/d $\to$ K & safe \checkmark \\
FEL.rem & every error world moved to distance 4 & \textsf{+} & \textsf{+} & 4 & \textsf{+} & \textsf{+} & G-nfl/d $\to$ K & safe below 4 \checkmark \\
FEL.fal & the belief made false at the actual world & \textsf{+} & \textsf{--} & 0 & \textsf{--} & \textsf{--} & G-nfl/d $\to$ F & not Gettier \checkmark \\
FEL.unj & justification removed & \textsf{--} & \textsf{+} & 1 & \textsf{--} & \textsf{--} & G-nfl/d $\to$ U & not knowledge \checkmark \\
FEL.rel & a hearer believes on the subject's sincere testimony & \textsf{+} & \textsf{+} & 1 & \textsf{--} & \textsf{--} & G-nfl/d $\to$ G-int/d & safety transmitted \checkmark \\
FEL.lie & relay from a speaker who could easily have lied & \textsf{+} & \textsf{+} & 1 & \textsf{--} & \textsf{--} & G-nfl/d $\to$ G-tst/d & unsafe \checkmark \\
FEL.ded & competent deduction of the weaker $q$ & \textsf{+} & \textsf{+} & $\infty$ & \textsf{+} & \textsf{+} & G-nfl/d $\to$ K & safe \checkmark \\
\midrule
CLK.nfl & false lemma removed; same evidence, true intermediates & \textsf{+} & \textsf{+} & 1 & \textsf{--} & \textsf{--} & G-tmp/d (same) & safety unchanged \checkmark \\
CLK.sec & truth-maker secured: the basis now bears on it & \textsf{+} & \textsf{+} & $\infty$ & \textsf{+} & \textsf{+} & G-tmp/d $\to$ K & safe \checkmark \\
CLK.hos & hostile environment: a look-alike is placed nearby & \textsf{+} & \textsf{+} & 1 & \textsf{--} & \textsf{--} & G-tmp/d (same) & unsafe \checkmark \\
CLK.rem & every error world moved to distance 4 & \textsf{+} & \textsf{+} & 4 & \textsf{+} & \textsf{+} & G-tmp/d $\to$ K & safe below 4 \checkmark \\
CLK.fal & the belief made false at the actual world & \textsf{+} & \textsf{--} & 0 & \textsf{--} & \textsf{--} & G-tmp/d $\to$ F & not Gettier \checkmark \\
CLK.unj & justification removed & \textsf{--} & \textsf{+} & 1 & \textsf{--} & \textsf{--} & G-tmp/d $\to$ U & not knowledge \checkmark \\
CLK.rel & a hearer believes on the subject's sincere testimony & \textsf{+} & \textsf{+} & 1 & \textsf{--} & \textsf{--} & G-tmp/d (same) & safety transmitted \checkmark \\
CLK.lie & relay from a speaker who could easily have lied & \textsf{+} & \textsf{+} & 1 & \textsf{--} & \textsf{--} & G-tmp/d (same) & unsafe \checkmark \\
\midrule
BRN.nor & environment made normal: look-alikes and bad sources removed & \textsf{+} & \textsf{+} & $\infty$ & \textsf{+} & \textsf{+} & G-env/c $\to$ K & safe \checkmark \\
BRN.rem & every error world moved to distance 4 & \textsf{+} & \textsf{+} & 4 & \textsf{+} & \textsf{+} & G-env/c $\to$ K & safe below 4 \checkmark \\
BRN.fal & the belief made false at the actual world & \textsf{+} & \textsf{--} & 0 & \textsf{--} & \textsf{--} & G-env/c $\to$ F & not Gettier \checkmark \\
BRN.unj & justification removed & \textsf{--} & \textsf{+} & 1 & \textsf{--} & \textsf{--} & G-env/c $\to$ U & not knowledge \checkmark \\
BRN.rel & a hearer believes on the subject's sincere testimony & \textsf{+} & \textsf{+} & 1 & \textsf{--} & \textsf{--} & G-env/c (same) & safety transmitted \checkmark \\
BRN.lie & relay from a speaker who could easily have lied & \textsf{+} & \textsf{+} & 1 & \textsf{--} & \textsf{--} & G-env/c $\to$ G-tst/c & unsafe \checkmark \\
BRN.ded & competent deduction of the weaker $q$ & \textsf{+} & \textsf{+} & $\infty$ & \textsf{+} & \textsf{+} & G-env/c $\to$ K & safe \checkmark \\
\midrule
RED.nor & environment made normal: look-alikes and bad sources removed & \textsf{+} & \textsf{+} & $\infty$ & \textsf{+} & \textsf{+} & K (same) & safe \checkmark \\
RED.hos & hostile environment: a look-alike is placed nearby & \textsf{+} & \textsf{+} & 1 & \textsf{--} & \textsf{--} & K $\to$ G-env/c & unsafe \checkmark \\
RED.fal & the belief made false at the actual world & \textsf{+} & \textsf{--} & 0 & \textsf{--} & \textsf{--} & K $\to$ F & not Gettier \checkmark \\
RED.unj & justification removed & \textsf{--} & \textsf{+} & 3 & \textsf{+} & \textsf{--} & K $\to$ U & not knowledge \checkmark \\
RED.rel & a hearer believes on the subject's sincere testimony & \textsf{+} & \textsf{+} & 3 & \textsf{+} & \textsf{+} & K (same) & safety transmitted \checkmark \\
RED.lie & relay from a speaker who could easily have lied & \textsf{+} & \textsf{+} & 1 & \textsf{--} & \textsf{--} & K $\to$ G-tst/c & unsafe \checkmark \\
RED.ded & competent deduction of the weaker $q$ & \textsf{+} & \textsf{+} & 3 & \textsf{+} & \textsf{+} & K (same) & safe (closure) \checkmark \\
\midrule
RDP.nor & environment made normal: look-alikes and bad sources removed & \textsf{+} & \textsf{+} & $\infty$ & \textsf{+} & \textsf{+} & G-env/c $\to$ K & safe \checkmark \\
RDP.rem & every error world moved to distance 4 & \textsf{+} & \textsf{+} & 4 & \textsf{+} & \textsf{+} & G-env/c $\to$ K & safe below 4 \checkmark \\
RDP.fal & the belief made false at the actual world & \textsf{+} & \textsf{--} & 0 & \textsf{--} & \textsf{--} & G-env/c $\to$ F & not Gettier \checkmark \\
RDP.unj & justification removed & \textsf{--} & \textsf{+} & 1 & \textsf{--} & \textsf{--} & G-env/c $\to$ U & not knowledge \checkmark \\
RDP.rel & a hearer believes on the subject's sincere testimony & \textsf{+} & \textsf{+} & 1 & \textsf{--} & \textsf{--} & G-env/c (same) & safety transmitted \checkmark \\
RDP.lie & relay from a speaker who could easily have lied & \textsf{+} & \textsf{+} & 1 & \textsf{--} & \textsf{--} & G-env/c $\to$ G-tst/c & unsafe \checkmark \\
\midrule
SHP.nfl & false lemma removed; same evidence, true intermediates & \textsf{+} & \textsf{+} & 1 & \textsf{--} & \textsf{--} & G-int/d (same) & safety unchanged \checkmark \\
SHP.sec & truth-maker secured: the basis now bears on it & \textsf{+} & \textsf{+} & $\infty$ & \textsf{+} & \textsf{+} & G-int/d $\to$ K & safe \checkmark \\
SHP.hos & hostile environment: a look-alike is placed nearby & \textsf{+} & \textsf{+} & 1 & \textsf{--} & \textsf{--} & G-int/d $\to$ G-env/d & unsafe \checkmark \\
SHP.rem & every error world moved to distance 4 & \textsf{+} & \textsf{+} & 4 & \textsf{+} & \textsf{+} & G-int/d $\to$ K & safe below 4 \checkmark \\
SHP.fal & the belief made false at the actual world & \textsf{+} & \textsf{--} & 0 & \textsf{--} & \textsf{--} & G-int/d $\to$ F & not Gettier \checkmark \\
SHP.unj & justification removed & \textsf{--} & \textsf{+} & 1 & \textsf{--} & \textsf{--} & G-int/d $\to$ U & not knowledge \checkmark \\
SHP.rel & a hearer believes on the subject's sincere testimony & \textsf{+} & \textsf{+} & 1 & \textsf{--} & \textsf{--} & G-int/d (same) & safety transmitted \checkmark \\
SHP.lie & relay from a speaker who could easily have lied & \textsf{+} & \textsf{+} & 1 & \textsf{--} & \textsf{--} & G-int/d $\to$ G-tst/d & unsafe \checkmark \\
\midrule
ZDR.nfl & false lemma removed; same evidence, true intermediates & \textsf{+} & \textsf{+} & 1 & \textsf{--} & \textsf{--} & G-fl/d $\to$ G-nfl/d & safety unchanged \checkmark \\
ZDR.sec & truth-maker secured: the basis now bears on it & \textsf{+} & \textsf{+} & $\infty$ & \textsf{+} & \textsf{+} & G-fl/d $\to$ K & safe \checkmark \\
ZDR.rem & every error world moved to distance 4 & \textsf{+} & \textsf{+} & 4 & \textsf{+} & \textsf{+} & G-fl/d $\to$ K & safe below 4 \checkmark \\
ZDR.fal & the belief made false at the actual world & \textsf{+} & \textsf{--} & 0 & \textsf{--} & \textsf{--} & G-fl/d $\to$ F & not Gettier \checkmark \\
ZDR.unj & justification removed & \textsf{--} & \textsf{+} & 1 & \textsf{--} & \textsf{--} & G-fl/d $\to$ U & not knowledge \checkmark \\
ZDR.rel & a hearer believes on the subject's sincere testimony & \textsf{+} & \textsf{+} & 1 & \textsf{--} & \textsf{--} & G-fl/d $\to$ G-int/d & safety transmitted \checkmark \\
ZDR.lie & relay from a speaker who could easily have lied & \textsf{+} & \textsf{+} & 1 & \textsf{--} & \textsf{--} & G-fl/d $\to$ G-tst/d & unsafe \checkmark \\
\midrule
ZHB.nfl & false lemma removed; same evidence, true intermediates & \textsf{+} & \textsf{+} & 1 & \textsf{--} & \textsf{--} & G-int/d (same) & safety unchanged \checkmark \\
ZHB.sec & truth-maker secured: the basis now bears on it & \textsf{+} & \textsf{+} & $\infty$ & \textsf{+} & \textsf{+} & G-int/d $\to$ K & safe \checkmark \\
ZHB.hos & hostile environment: a look-alike is placed nearby & \textsf{+} & \textsf{+} & 1 & \textsf{--} & \textsf{--} & G-int/d $\to$ G-env/d & unsafe \checkmark \\
ZHB.rem & every error world moved to distance 4 & \textsf{+} & \textsf{+} & 4 & \textsf{+} & \textsf{+} & G-int/d $\to$ K & safe below 4 \checkmark \\
ZHB.fal & the belief made false at the actual world & \textsf{+} & \textsf{--} & 0 & \textsf{--} & \textsf{--} & G-int/d $\to$ F & not Gettier \checkmark \\
ZHB.unj & justification removed & \textsf{--} & \textsf{+} & 1 & \textsf{--} & \textsf{--} & G-int/d $\to$ U & not knowledge \checkmark \\
ZHB.rel & a hearer believes on the subject's sincere testimony & \textsf{+} & \textsf{+} & 1 & \textsf{--} & \textsf{--} & G-int/d (same) & safety transmitted \checkmark \\
ZHB.lie & relay from a speaker who could easily have lied & \textsf{+} & \textsf{+} & 1 & \textsf{--} & \textsf{--} & G-int/d $\to$ G-tst/d & unsafe \checkmark \\
\midrule
PYR.nfl & false lemma removed; same evidence, true intermediates & \textsf{+} & \textsf{+} & 1 & \textsf{--} & \textsf{--} & G-fl/d $\to$ G-nfl/d & safety unchanged \checkmark \\
PYR.sec & truth-maker secured: the basis now bears on it & \textsf{+} & \textsf{+} & $\infty$ & \textsf{+} & \textsf{+} & G-fl/d $\to$ K & safe \checkmark \\
PYR.rem & every error world moved to distance 4 & \textsf{+} & \textsf{+} & 4 & \textsf{+} & \textsf{+} & G-fl/d $\to$ K & safe below 4 \checkmark \\
PYR.fal & the belief made false at the actual world & \textsf{+} & \textsf{--} & 0 & \textsf{--} & \textsf{--} & G-fl/d $\to$ F & not Gettier \checkmark \\
PYR.unj & justification removed & \textsf{--} & \textsf{+} & 1 & \textsf{--} & \textsf{--} & G-fl/d $\to$ U & not knowledge \checkmark \\
PYR.rel & a hearer believes on the subject's sincere testimony & \textsf{+} & \textsf{+} & 1 & \textsf{--} & \textsf{--} & G-fl/d $\to$ G-int/d & safety transmitted \checkmark \\
PYR.lie & relay from a speaker who could easily have lied & \textsf{+} & \textsf{+} & 1 & \textsf{--} & \textsf{--} & G-fl/d $\to$ G-tst/d & unsafe \checkmark \\
\midrule
GRB.nor & environment made normal: look-alikes and bad sources removed & \textsf{+} & \textsf{+} & $\infty$ & \textsf{+} & \textsf{+} & K[md] (same) & safe \checkmark \\
GRB.hos & hostile environment: a look-alike is placed nearby & \textsf{+} & \textsf{+} & 1 & \textsf{--} & \textsf{--} & K[md] $\to$ G-env/c & unsafe \checkmark \\
GRB.fal & the belief made false at the actual world & \textsf{+} & \textsf{--} & 0 & \textsf{--} & \textsf{--} & K[md] $\to$ F & not Gettier \checkmark \\
GRB.unj & justification removed & \textsf{--} & \textsf{+} & 4 & \textsf{+} & \textsf{--} & K[md] $\to$ U & not knowledge \checkmark \\
GRB.rel & a hearer believes on the subject's sincere testimony & \textsf{+} & \textsf{+} & 4 & \textsf{+} & \textsf{+} & K[md] (same) & safety transmitted \checkmark \\
GRB.lie & relay from a speaker who could easily have lied & \textsf{+} & \textsf{+} & 1 & \textsf{--} & \textsf{--} & K[md] $\to$ G-tst/c & unsafe \checkmark \\
\midrule
HAR.hos & hostile environment: a look-alike is placed nearby & \textsf{+} & \textsf{+} & 1 & \textsf{--} & \textsf{--} & K[md] $\to$ G-env/c & unsafe \checkmark \\
HAR.fal & the belief made false at the actual world & \textsf{+} & \textsf{--} & 0 & \textsf{--} & \textsf{--} & K[md] $\to$ F & not Gettier \checkmark \\
HAR.unj & justification removed & \textsf{--} & \textsf{+} & $\infty$ & \textsf{+} & \textsf{--} & K[md] $\to$ U & not knowledge \checkmark \\
HAR.rel & a hearer believes on the subject's sincere testimony & \textsf{+} & \textsf{+} & $\infty$ & \textsf{+} & \textsf{+} & K[md] (same) & safety transmitted \checkmark \\
HAR.lie & relay from a speaker who could easily have lied & \textsf{+} & \textsf{+} & 1 & \textsf{--} & \textsf{--} & K[md] $\to$ G-tst/c & unsafe \checkmark \\
\midrule
CLV.fal & the belief made false at the actual world & \textsf{--} & \textsf{--} & 0 & \textsf{--} & \textsf{--} & U (same) & not Gettier \checkmark \\
CLV.rel & a hearer believes on the subject's sincere testimony & \textsf{--} & \textsf{+} & $\infty$ & \textsf{+} & \textsf{--} & U (same) & safety transmitted \checkmark \\
CLV.lie & relay from a speaker who could easily have lied & \textsf{--} & \textsf{+} & 1 & \textsf{--} & \textsf{--} & U (same) & unsafe \checkmark \\
\midrule
TTP.fal & the belief made false at the actual world & \textsf{--} & \textsf{--} & 0 & \textsf{--} & \textsf{--} & U (same) & not Gettier \checkmark \\
TTP.rel & a hearer believes on the subject's sincere testimony & \textsf{--} & \textsf{+} & $\infty$ & \textsf{+} & \textsf{--} & U (same) & safety transmitted \checkmark \\
TTP.lie & relay from a speaker who could easily have lied & \textsf{--} & \textsf{+} & 1 & \textsf{--} & \textsf{--} & U (same) & unsafe \checkmark \\
\midrule
CHU.fl & a false intermediate belief added to the route & \textsf{+} & \textsf{+} & 4 & \textsf{+} & \textsf{+} & K (same) & safety unchanged \checkmark \\
CHU.sec & truth-maker secured: the basis now bears on it & \textsf{+} & \textsf{+} & $\infty$ & \textsf{+} & \textsf{+} & K (same) & safe \checkmark \\
CHU.fal & the belief made false at the actual world & \textsf{+} & \textsf{--} & 0 & \textsf{--} & \textsf{--} & K $\to$ F & not Gettier \checkmark \\
CHU.unj & justification removed & \textsf{--} & \textsf{+} & 4 & \textsf{+} & \textsf{--} & K $\to$ U & not knowledge \checkmark \\
CHU.rel & a hearer believes on the subject's sincere testimony & \textsf{+} & \textsf{+} & 4 & \textsf{+} & \textsf{+} & K (same) & safety transmitted \checkmark \\
CHU.lie & relay from a speaker who could easily have lied & \textsf{+} & \textsf{+} & 1 & \textsf{--} & \textsf{--} & K $\to$ G-tst/c & unsafe \checkmark \\
\midrule
KAL.nor & environment made normal: look-alikes and bad sources removed & \textsf{+} & \textsf{+} & $\infty$ & \textsf{+} & \textsf{+} & G-env/c $\to$ K[rf] & safe \checkmark \\
KAL.rem & every error world moved to distance 4 & \textsf{+} & \textsf{+} & 4 & \textsf{+} & \textsf{+} & G-env/c $\to$ K[rf] & safe below 4 \checkmark \\
KAL.fal & the belief made false at the actual world & \textsf{+} & \textsf{--} & 0 & \textsf{--} & \textsf{--} & G-env/c $\to$ F & not Gettier \checkmark \\
KAL.unj & justification removed & \textsf{--} & \textsf{+} & 1 & \textsf{--} & \textsf{--} & G-env/c $\to$ U & not knowledge \checkmark \\
KAL.rel & a hearer believes on the subject's sincere testimony & \textsf{+} & \textsf{+} & 1 & \textsf{--} & \textsf{--} & G-env/c (same) & safety transmitted \checkmark \\
KAL.lie & relay from a speaker who could easily have lied & \textsf{+} & \textsf{+} & 1 & \textsf{--} & \textsf{--} & G-env/c $\to$ G-tst/c & unsafe \checkmark \\
\midrule
TMP.sec & truth-maker secured: the basis now bears on it & \textsf{+} & \textsf{+} & $\infty$ & \textsf{+} & \textsf{+} & K[fx] (same) & safe \checkmark \\
TMP.hos & hostile environment: a look-alike is placed nearby & \textsf{+} & \textsf{+} & 1 & \textsf{--} & \textsf{--} & K[fx] $\to$ G-env/d & unsafe \checkmark \\
TMP.fal & the belief made false at the actual world & \textsf{+} & \textsf{--} & 0 & \textsf{--} & \textsf{--} & K[fx] $\to$ F & not Gettier \checkmark \\
TMP.unj & justification removed & \textsf{--} & \textsf{+} & 3 & \textsf{+} & \textsf{--} & K[fx] $\to$ U & not knowledge \checkmark \\
TMP.rel & a hearer believes on the subject's sincere testimony & \textsf{+} & \textsf{+} & 3 & \textsf{+} & \textsf{+} & K[fx] (same) & safety transmitted \checkmark \\
TMP.lie & relay from a speaker who could easily have lied & \textsf{+} & \textsf{+} & 1 & \textsf{--} & \textsf{--} & K[fx] $\to$ G-tst/d & unsafe \checkmark \\
\midrule
MOR.nor & environment made normal: look-alikes and bad sources removed & \textsf{+} & \textsf{+} & $\infty$ & \textsf{+} & \textsf{+} & K (same) & safe \checkmark \\
MOR.hos & hostile environment: a look-alike is placed nearby & \textsf{+} & \textsf{+} & 1 & \textsf{--} & \textsf{--} & K $\to$ G-env/c & unsafe \checkmark \\
MOR.fal & the belief made false at the actual world & \textsf{+} & \textsf{--} & 0 & \textsf{--} & \textsf{--} & K $\to$ F & not Gettier \checkmark \\
MOR.unj & justification removed & \textsf{--} & \textsf{+} & 3 & \textsf{+} & \textsf{--} & K $\to$ U & not knowledge \checkmark \\
MOR.rel & a hearer believes on the subject's sincere testimony & \textsf{+} & \textsf{+} & 3 & \textsf{+} & \textsf{+} & K (same) & safety transmitted \checkmark \\
MOR.lie & relay from a speaker who could easily have lied & \textsf{+} & \textsf{+} & 1 & \textsf{--} & \textsf{--} & K $\to$ G-tst/c & unsafe \checkmark \\
\midrule
LOT.sec & truth-maker secured: the basis now bears on it & \textsf{+} & \textsf{+} & $\infty$ & \textsf{+} & \textsf{+} & G-lot/d $\to$ K & safe \checkmark \\
LOT.rem & every error world moved to distance 4 & \textsf{+} & \textsf{+} & 4 & \textsf{+} & \textsf{+} & G-lot/d $\to$ K & safe below 4 \checkmark \\
LOT.fal & the belief made false at the actual world & \textsf{+} & \textsf{--} & 0 & \textsf{--} & \textsf{--} & G-lot/d $\to$ F & not Gettier \checkmark \\
LOT.unj & justification removed & \textsf{--} & \textsf{+} & 1 & \textsf{--} & \textsf{--} & G-lot/d $\to$ U & not knowledge \checkmark \\
LOT.rel & a hearer believes on the subject's sincere testimony & \textsf{+} & \textsf{+} & 1 & \textsf{--} & \textsf{--} & G-lot/d (same) & safety transmitted \checkmark \\
LOT.lie & relay from a speaker who could easily have lied & \textsf{+} & \textsf{+} & 1 & \textsf{--} & \textsf{--} & G-lot/d (same) & unsafe \checkmark \\
\midrule
LTN.nor & environment made normal: look-alikes and bad sources removed & \textsf{+} & \textsf{+} & $\infty$ & \textsf{+} & \textsf{+} & K (same) & safe \checkmark \\
LTN.hos & hostile environment: a look-alike is placed nearby & \textsf{+} & \textsf{+} & 1 & \textsf{--} & \textsf{--} & K $\to$ G-env/c & unsafe \checkmark \\
LTN.fal & the belief made false at the actual world & \textsf{+} & \textsf{--} & 0 & \textsf{--} & \textsf{--} & K $\to$ F & not Gettier \checkmark \\
LTN.unj & justification removed & \textsf{--} & \textsf{+} & 3 & \textsf{+} & \textsf{--} & K $\to$ U & not knowledge \checkmark \\
LTN.rel & a hearer believes on the subject's sincere testimony & \textsf{+} & \textsf{+} & 3 & \textsf{+} & \textsf{+} & K (same) & safety transmitted \checkmark \\
LTN.lie & relay from a speaker who could easily have lied & \textsf{+} & \textsf{+} & 1 & \textsf{--} & \textsf{--} & K $\to$ G-tst/c & unsafe \checkmark \\
\midrule
MAG.hos & hostile environment: a look-alike is placed nearby & \textsf{+} & \textsf{+} & 1 & \textsf{--} & \textsf{--} & G-mrg/c (same) & unsafe \checkmark \\
MAG.rem & every error world moved to distance 4 & \textsf{+} & \textsf{+} & 4 & \textsf{+} & \textsf{+} & G-mrg/c $\to$ K & safe below 4 \checkmark \\
MAG.fal & the belief made false at the actual world & \textsf{+} & \textsf{--} & 0 & \textsf{--} & \textsf{--} & G-mrg/c $\to$ F & not Gettier \checkmark \\
MAG.unj & justification removed & \textsf{--} & \textsf{+} & 1 & \textsf{--} & \textsf{--} & G-mrg/c $\to$ U & not knowledge \checkmark \\
MAG.rel & a hearer believes on the subject's sincere testimony & \textsf{+} & \textsf{+} & 1 & \textsf{--} & \textsf{--} & G-mrg/c (same) & safety transmitted \checkmark \\
MAG.lie & relay from a speaker who could easily have lied & \textsf{+} & \textsf{+} & 1 & \textsf{--} & \textsf{--} & G-mrg/c (same) & unsafe \checkmark \\
\midrule
MGW.hos & hostile environment: a look-alike is placed nearby & \textsf{+} & \textsf{+} & 1 & \textsf{--} & \textsf{--} & K $\to$ G-env/c & unsafe \checkmark \\
MGW.fal & the belief made false at the actual world & \textsf{+} & \textsf{--} & 0 & \textsf{--} & \textsf{--} & K $\to$ F & not Gettier \checkmark \\
MGW.unj & justification removed & \textsf{--} & \textsf{+} & $\infty$ & \textsf{+} & \textsf{--} & K $\to$ U & not knowledge \checkmark \\
MGW.rel & a hearer believes on the subject's sincere testimony & \textsf{+} & \textsf{+} & $\infty$ & \textsf{+} & \textsf{+} & K (same) & safety transmitted \checkmark \\
MGW.lie & relay from a speaker who could easily have lied & \textsf{+} & \textsf{+} & 1 & \textsf{--} & \textsf{--} & K $\to$ G-tst/c & unsafe \checkmark \\
\midrule
ZBP.nor & environment made normal: look-alikes and bad sources removed & \textsf{+} & \textsf{+} & $\infty$ & \textsf{+} & \textsf{+} & K (same) & safe \checkmark \\
ZBP.hos & hostile environment: a look-alike is placed nearby & \textsf{+} & \textsf{+} & 1 & \textsf{--} & \textsf{--} & K $\to$ G-env/c & unsafe \checkmark \\
ZBP.fal & the belief made false at the actual world & \textsf{+} & \textsf{--} & 0 & \textsf{--} & \textsf{--} & K $\to$ F & not Gettier \checkmark \\
ZBP.unj & justification removed & \textsf{--} & \textsf{+} & 4 & \textsf{+} & \textsf{--} & K $\to$ U & not knowledge \checkmark \\
ZBP.rel & a hearer believes on the subject's sincere testimony & \textsf{+} & \textsf{+} & 4 & \textsf{+} & \textsf{+} & K (same) & safety transmitted \checkmark \\
ZBP.lie & relay from a speaker who could easily have lied & \textsf{+} & \textsf{+} & 1 & \textsf{--} & \textsf{--} & K $\to$ G-tst/c & unsafe \checkmark \\
ZBP.ded & competent deduction of the weaker $q$ & \textsf{+} & \textsf{+} & 4 & \textsf{+} & \textsf{+} & K (same) & safe (closure) \checkmark \\
\midrule
ZBM.fl & a false intermediate belief added to the route & \textsf{+} & \textsf{+} & 4 & \textsf{+} & \textsf{+} & K (same) & safety unchanged \checkmark \\
ZBM.nor & environment made normal: look-alikes and bad sources removed & \textsf{+} & \textsf{+} & $\infty$ & \textsf{+} & \textsf{+} & K (same) & safe \checkmark \\
ZBM.fal & the belief made false at the actual world & \textsf{+} & \textsf{--} & 0 & \textsf{--} & \textsf{--} & K $\to$ F & not Gettier \checkmark \\
ZBM.unj & justification removed & \textsf{--} & \textsf{+} & 4 & \textsf{+} & \textsf{--} & K $\to$ U & not knowledge \checkmark \\
ZBM.rel & a hearer believes on the subject's sincere testimony & \textsf{+} & \textsf{+} & 4 & \textsf{+} & \textsf{+} & K (same) & safety transmitted \checkmark \\
ZBM.lie & relay from a speaker who could easily have lied & \textsf{+} & \textsf{+} & 1 & \textsf{--} & \textsf{--} & K $\to$ G-tst/c & unsafe \checkmark \\
\midrule
GRN.hos & hostile environment: a look-alike is placed nearby & \textsf{+} & \textsf{+} & 1 & \textsf{--} & \textsf{--} & K[ev] $\to$ G-env/c & unsafe \checkmark \\
GRN.fal & the belief made false at the actual world & \textsf{+} & \textsf{--} & 0 & \textsf{--} & \textsf{--} & K[ev] $\to$ F & not Gettier \checkmark \\
GRN.unj & justification removed & \textsf{--} & \textsf{+} & $\infty$ & \textsf{+} & \textsf{--} & K[ev] $\to$ U & not knowledge \checkmark \\
GRN.rel & a hearer believes on the subject's sincere testimony & \textsf{+} & \textsf{+} & $\infty$ & \textsf{+} & \textsf{+} & K[ev] (same) & safety transmitted \checkmark \\
GRN.lie & relay from a speaker who could easily have lied & \textsf{+} & \textsf{+} & 1 & \textsf{--} & \textsf{--} & K[ev] $\to$ G-tst/c & unsafe \checkmark \\
\midrule
HLF.hos & hostile environment: a look-alike is placed nearby & \textsf{+} & \textsf{+} & 1 & \textsf{--} & \textsf{--} & K $\to$ G-env/c & unsafe \checkmark \\
HLF.fal & the belief made false at the actual world & \textsf{+} & \textsf{--} & 0 & \textsf{--} & \textsf{--} & K $\to$ F & not Gettier \checkmark \\
HLF.unj & justification removed & \textsf{--} & \textsf{+} & $\infty$ & \textsf{+} & \textsf{--} & K $\to$ U & not knowledge \checkmark \\
HLF.rel & a hearer believes on the subject's sincere testimony & \textsf{+} & \textsf{+} & $\infty$ & \textsf{+} & \textsf{+} & K (same) & safety transmitted \checkmark \\
HLF.lie & relay from a speaker who could easily have lied & \textsf{+} & \textsf{+} & 1 & \textsf{--} & \textsf{--} & K $\to$ G-tst/c & unsafe \checkmark \\
\midrule
HLC.nor & environment made normal: look-alikes and bad sources removed & \textsf{+} & \textsf{+} & $\infty$ & \textsf{+} & \textsf{+} & G-tst/c $\to$ K & safe \checkmark \\
HLC.rem & every error world moved to distance 4 & \textsf{+} & \textsf{+} & 4 & \textsf{+} & \textsf{+} & G-tst/c $\to$ K & safe below 4 \checkmark \\
HLC.fal & the belief made false at the actual world & \textsf{+} & \textsf{--} & 0 & \textsf{--} & \textsf{--} & G-tst/c $\to$ F & not Gettier \checkmark \\
HLC.unj & justification removed & \textsf{--} & \textsf{+} & 1 & \textsf{--} & \textsf{--} & G-tst/c $\to$ U & not knowledge \checkmark \\
HLC.rel & a hearer believes on the subject's sincere testimony & \textsf{+} & \textsf{+} & 1 & \textsf{--} & \textsf{--} & G-tst/c (same) & safety transmitted \checkmark \\
HLC.lie & relay from a speaker who could easily have lied & \textsf{+} & \textsf{+} & 1 & \textsf{--} & \textsf{--} & G-tst/c (same) & unsafe \checkmark \\
\midrule
ATF.hos & hostile environment: a look-alike is placed nearby & \textsf{+} & \textsf{+} & 1 & \textsf{--} & \textsf{--} & K $\to$ G-env/c & unsafe \checkmark \\
ATF.fal & the belief made false at the actual world & \textsf{+} & \textsf{--} & 0 & \textsf{--} & \textsf{--} & K $\to$ F & not Gettier \checkmark \\
ATF.unj & justification removed & \textsf{--} & \textsf{+} & $\infty$ & \textsf{+} & \textsf{--} & K $\to$ U & not knowledge \checkmark \\
ATF.rel & a hearer believes on the subject's sincere testimony & \textsf{+} & \textsf{+} & $\infty$ & \textsf{+} & \textsf{+} & K (same) & safety transmitted \checkmark \\
ATF.lie & relay from a speaker who could easily have lied & \textsf{+} & \textsf{+} & 1 & \textsf{--} & \textsf{--} & K $\to$ G-tst/c & unsafe \checkmark \\
\midrule
ATC.nor & environment made normal: look-alikes and bad sources removed & \textsf{+} & \textsf{+} & $\infty$ & \textsf{+} & \textsf{+} & G-env/c $\to$ K & safe \checkmark \\
ATC.rem & every error world moved to distance 4 & \textsf{+} & \textsf{+} & 4 & \textsf{+} & \textsf{+} & G-env/c $\to$ K & safe below 4 \checkmark \\
ATC.fal & the belief made false at the actual world & \textsf{+} & \textsf{--} & 0 & \textsf{--} & \textsf{--} & G-env/c $\to$ F & not Gettier \checkmark \\
ATC.unj & justification removed & \textsf{--} & \textsf{+} & 1 & \textsf{--} & \textsf{--} & G-env/c $\to$ U & not knowledge \checkmark \\
ATC.rel & a hearer believes on the subject's sincere testimony & \textsf{+} & \textsf{+} & 1 & \textsf{--} & \textsf{--} & G-env/c (same) & safety transmitted \checkmark \\
ATC.lie & relay from a speaker who could easily have lied & \textsf{+} & \textsf{+} & 1 & \textsf{--} & \textsf{--} & G-env/c $\to$ G-tst/c & unsafe \checkmark \\
\midrule
WAR.nfl & false lemma removed; same evidence, true intermediates & \textsf{+} & \textsf{+} & $\infty$ & \textsf{+} & \textsf{+} & K (same) & safety unchanged \checkmark \\
WAR.fal & the belief made false at the actual world & \textsf{+} & \textsf{--} & 0 & \textsf{--} & \textsf{--} & K $\to$ F & not Gettier \checkmark \\
WAR.unj & justification removed & \textsf{--} & \textsf{+} & $\infty$ & \textsf{+} & \textsf{--} & K $\to$ U & not knowledge \checkmark \\
WAR.rel & a hearer believes on the subject's sincere testimony & \textsf{+} & \textsf{+} & $\infty$ & \textsf{+} & \textsf{+} & K (same) & safety transmitted \checkmark \\
WAR.lie & relay from a speaker who could easily have lied & \textsf{+} & \textsf{+} & 1 & \textsf{--} & \textsf{--} & K $\to$ G-tst/c & unsafe \checkmark \\
\midrule
BKL.fl & a false intermediate belief added to the route & \textsf{+} & \textsf{+} & 3 & \textsf{+} & \textsf{+} & K (same) & safety unchanged \checkmark \\
BKL.sec & truth-maker secured: the basis now bears on it & \textsf{+} & \textsf{+} & $\infty$ & \textsf{+} & \textsf{+} & K (same) & safe \checkmark \\
BKL.fal & the belief made false at the actual world & \textsf{+} & \textsf{--} & 0 & \textsf{--} & \textsf{--} & K $\to$ F & not Gettier \checkmark \\
BKL.unj & justification removed & \textsf{--} & \textsf{+} & 3 & \textsf{+} & \textsf{--} & K $\to$ U & not knowledge \checkmark \\
BKL.rel & a hearer believes on the subject's sincere testimony & \textsf{+} & \textsf{+} & 3 & \textsf{+} & \textsf{+} & K (same) & safety transmitted \checkmark \\
BKL.lie & relay from a speaker who could easily have lied & \textsf{+} & \textsf{+} & 1 & \textsf{--} & \textsf{--} & K $\to$ G-tst/c & unsafe \checkmark \\
\midrule
BKH.fl & a false intermediate belief added to the route & \textsf{+} & \textsf{+} & 3 & \textsf{--} & \textsf{--} & G-nfl/c $\to$ G-fl/c & safety unchanged \checkmark \\
BKH.sec & truth-maker secured: the basis now bears on it & \textsf{+} & \textsf{+} & $\infty$ & \textsf{+} & \textsf{+} & G-nfl/c $\to$ K & safe \checkmark \\
BKH.rem & every error world moved to distance 4 & \textsf{+} & \textsf{+} & 4 & \textsf{+} & \textsf{+} & G-nfl/c $\to$ K & safe below 4 \checkmark \\
BKH.fal & the belief made false at the actual world & \textsf{+} & \textsf{--} & 0 & \textsf{--} & \textsf{--} & G-nfl/c $\to$ F & not Gettier \checkmark \\
BKH.unj & justification removed & \textsf{--} & \textsf{+} & 3 & \textsf{--} & \textsf{--} & G-nfl/c $\to$ U & not knowledge \checkmark \\
BKH.rel & a hearer believes on the subject's sincere testimony & \textsf{+} & \textsf{+} & 3 & \textsf{--} & \textsf{--} & G-nfl/c $\to$ G-int/c & safety transmitted \checkmark \\
BKH.lie & relay from a speaker who could easily have lied & \textsf{+} & \textsf{+} & 1 & \textsf{--} & \textsf{--} & G-nfl/c $\to$ G-tst/c & unsafe \checkmark \\
\midrule
BIV.nor & environment made normal: look-alikes and bad sources removed & \textsf{+} & \textsf{+} & $\infty$ & \textsf{+} & \textsf{+} & K (same) & safe \checkmark \\
BIV.hos & hostile environment: a look-alike is placed nearby & \textsf{+} & \textsf{+} & 1 & \textsf{--} & \textsf{--} & K $\to$ G-env/c & unsafe \checkmark \\
BIV.fal & the belief made false at the actual world & \textsf{+} & \textsf{--} & 0 & \textsf{--} & \textsf{--} & K $\to$ F & not Gettier \checkmark \\
BIV.unj & justification removed & \textsf{--} & \textsf{+} & 4 & \textsf{+} & \textsf{--} & K $\to$ U & not knowledge \checkmark \\
BIV.rel & a hearer believes on the subject's sincere testimony & \textsf{+} & \textsf{+} & 4 & \textsf{+} & \textsf{+} & K (same) & safety transmitted \checkmark \\
BIV.lie & relay from a speaker who could easily have lied & \textsf{+} & \textsf{+} & 1 & \textsf{--} & \textsf{--} & K $\to$ G-tst/c & unsafe \checkmark \\
\midrule
JJM.nor & environment made normal: look-alikes and bad sources removed & \textsf{+} & \textsf{+} & $\infty$ & \textsf{+} & \textsf{+} & K[ev] (same) & safe \checkmark \\
JJM.hos & hostile environment: a look-alike is placed nearby & \textsf{+} & \textsf{+} & 1 & \textsf{--} & \textsf{--} & K[ev] $\to$ G-env/c & unsafe \checkmark \\
JJM.fal & the belief made false at the actual world & \textsf{+} & \textsf{--} & 0 & \textsf{--} & \textsf{--} & K[ev] $\to$ F & not Gettier \checkmark \\
JJM.unj & justification removed & \textsf{--} & \textsf{+} & 4 & \textsf{+} & \textsf{--} & K[ev] $\to$ U & not knowledge \checkmark \\
JJM.rel & a hearer believes on the subject's sincere testimony & \textsf{+} & \textsf{+} & 4 & \textsf{+} & \textsf{+} & K[ev] (same) & safety transmitted \checkmark \\
JJM.lie & relay from a speaker who could easily have lied & \textsf{+} & \textsf{+} & 1 & \textsf{--} & \textsf{--} & K[ev] $\to$ G-tst/c & unsafe \checkmark \\
\midrule
THT.sec & truth-maker secured: the basis now bears on it & \textsf{--} & \textsf{+} & $\infty$ & \textsf{+} & \textsf{--} & U (same) & safe \checkmark \\
THT.hos & hostile environment: a look-alike is placed nearby & \textsf{--} & \textsf{+} & 1 & \textsf{--} & \textsf{--} & U (same) & unsafe \checkmark \\
THT.rem & every error world moved to distance 4 & \textsf{--} & \textsf{+} & 4 & \textsf{+} & \textsf{--} & U (same) & safe below 4 \checkmark \\
THT.fal & the belief made false at the actual world & \textsf{--} & \textsf{--} & 0 & \textsf{--} & \textsf{--} & U (same) & not Gettier \checkmark \\
THT.rel & a hearer believes on the subject's sincere testimony & \textsf{--} & \textsf{+} & 1 & \textsf{--} & \textsf{--} & U (same) & safety transmitted \checkmark \\
THT.lie & relay from a speaker who could easily have lied & \textsf{--} & \textsf{+} & 1 & \textsf{--} & \textsf{--} & U (same) & unsafe \checkmark \\
\midrule
ICE.fl & a false intermediate belief added to the route & \textsf{+} & \textsf{+} & 4 & \textsf{+} & \textsf{+} & K (same) & safety unchanged \checkmark \\
ICE.sec & truth-maker secured: the basis now bears on it & \textsf{+} & \textsf{+} & $\infty$ & \textsf{+} & \textsf{+} & K (same) & safe \checkmark \\
ICE.fal & the belief made false at the actual world & \textsf{+} & \textsf{--} & 0 & \textsf{--} & \textsf{--} & K $\to$ F & not Gettier \checkmark \\
ICE.unj & justification removed & \textsf{--} & \textsf{+} & 4 & \textsf{+} & \textsf{--} & K $\to$ U & not knowledge \checkmark \\
ICE.rel & a hearer believes on the subject's sincere testimony & \textsf{+} & \textsf{+} & 4 & \textsf{+} & \textsf{+} & K (same) & safety transmitted \checkmark \\
ICE.lie & relay from a speaker who could easily have lied & \textsf{+} & \textsf{+} & 1 & \textsf{--} & \textsf{--} & K $\to$ G-tst/c & unsafe \checkmark \\
\midrule
SNC.rem & every error world moved to distance 4 & \textsf{+} & \textsf{+} & 4 & \textsf{+} & \textsf{+} & G-nec/d $\to$ K & safe below 4 \checkmark \\
SNC.fal & the belief made false at the actual world & \textsf{+} & \textsf{--} & 0 & \textsf{--} & \textsf{--} & G-nec/d $\to$ F & not Gettier \checkmark \\
SNC.unj & justification removed & \textsf{--} & \textsf{+} & 1 & \textsf{--} & \textsf{--} & G-nec/d $\to$ U & not knowledge \checkmark \\
SNC.rel & a hearer believes on the subject's sincere testimony & \textsf{+} & \textsf{+} & 1 & \textsf{--} & \textsf{--} & G-nec/d (same) & safety transmitted \checkmark \\
SNC.lie & relay from a speaker who could easily have lied & \textsf{+} & \textsf{+} & 1 & \textsf{--} & \textsf{--} & G-nec/d $\to$ G-tst/d & unsafe \checkmark \\
\midrule
SNS.fal & the belief made false at the actual world & \textsf{+} & \textsf{--} & 0 & \textsf{--} & \textsf{--} & K $\to$ F & not Gettier \checkmark \\
SNS.unj & justification removed & \textsf{--} & \textsf{+} & $\infty$ & \textsf{+} & \textsf{--} & K $\to$ U & not knowledge \checkmark \\
SNS.rel & a hearer believes on the subject's sincere testimony & \textsf{+} & \textsf{+} & $\infty$ & \textsf{+} & \textsf{+} & K (same) & safety transmitted \checkmark \\
SNS.lie & relay from a speaker who could easily have lied & \textsf{+} & \textsf{+} & 1 & \textsf{--} & \textsf{--} & K $\to$ G-tst/d & unsafe \checkmark \\
\midrule
SEV.nor & environment made normal: look-alikes and bad sources removed & \textsf{+} & \textsf{+} & $\infty$ & \textsf{+} & \textsf{+} & K[ev] (same) & safe \checkmark \\
SEV.hos & hostile environment: a look-alike is placed nearby & \textsf{+} & \textsf{+} & 1 & \textsf{--} & \textsf{--} & K[ev] $\to$ G-env/c & unsafe \checkmark \\
SEV.fal & the belief made false at the actual world & \textsf{+} & \textsf{--} & 0 & \textsf{--} & \textsf{--} & K[ev] $\to$ F & not Gettier \checkmark \\
SEV.unj & justification removed & \textsf{--} & \textsf{+} & 4 & \textsf{+} & \textsf{--} & K[ev] $\to$ U & not knowledge \checkmark \\
SEV.rel & a hearer believes on the subject's sincere testimony & \textsf{+} & \textsf{+} & 4 & \textsf{+} & \textsf{+} & K[ev] (same) & safety transmitted \checkmark \\
SEV.lie & relay from a speaker who could easily have lied & \textsf{+} & \textsf{+} & 1 & \textsf{--} & \textsf{--} & K[ev] $\to$ G-tst/c & unsafe \checkmark \\
\midrule
SCT.nor & environment made normal: look-alikes and bad sources removed & \textsf{+} & \textsf{+} & $\infty$ & \textsf{+} & \textsf{+} & K[ct] (same) & safe \checkmark \\
SCT.hos & hostile environment: a look-alike is placed nearby & \textsf{+} & \textsf{+} & 1 & \textsf{--} & \textsf{--} & K[ct] $\to$ G-env/c & unsafe \checkmark \\
SCT.fal & the belief made false at the actual world & \textsf{+} & \textsf{--} & 0 & \textsf{--} & \textsf{--} & K[ct] $\to$ F & not Gettier \checkmark \\
SCT.unj & justification removed & \textsf{--} & \textsf{+} & 4 & \textsf{+} & \textsf{--} & K[ct] $\to$ U & not knowledge \checkmark \\
SCT.rel & a hearer believes on the subject's sincere testimony & \textsf{+} & \textsf{+} & 4 & \textsf{+} & \textsf{+} & K[ct] (same) & safety transmitted \checkmark \\
SCT.lie & relay from a speaker who could easily have lied & \textsf{+} & \textsf{+} & 1 & \textsf{--} & \textsf{--} & K[ct] $\to$ G-tst/c & unsafe \checkmark \\
\midrule
SRF.hos & hostile environment: a look-alike is placed nearby & \textsf{+} & \textsf{+} & 1 & \textsf{--} & \textsf{--} & K[rf] $\to$ G-env/c & unsafe \checkmark \\
SRF.fal & the belief made false at the actual world & \textsf{+} & \textsf{--} & 0 & \textsf{--} & \textsf{--} & K[rf] $\to$ F & not Gettier \checkmark \\
SRF.unj & justification removed & \textsf{--} & \textsf{+} & $\infty$ & \textsf{+} & \textsf{--} & K[rf] $\to$ U & not knowledge \checkmark \\
SRF.rel & a hearer believes on the subject's sincere testimony & \textsf{+} & \textsf{+} & $\infty$ & \textsf{+} & \textsf{+} & K[rf] (same) & safety transmitted \checkmark \\
SRF.lie & relay from a speaker who could easily have lied & \textsf{+} & \textsf{+} & 1 & \textsf{--} & \textsf{--} & K[rf] $\to$ G-tst/c & unsafe \checkmark \\
\midrule
STL.nor & environment made normal: look-alikes and bad sources removed & \textsf{+} & \textsf{+} & $\infty$ & \textsf{+} & \textsf{+} & G-tst/d $\to$ K[ev] & safe \checkmark \\
STL.rem & every error world moved to distance 4 & \textsf{+} & \textsf{+} & 4 & \textsf{+} & \textsf{+} & G-tst/d $\to$ K[ev] & safe below 4 \checkmark \\
STL.fal & the belief made false at the actual world & \textsf{+} & \textsf{--} & 0 & \textsf{--} & \textsf{--} & G-tst/d $\to$ F & not Gettier \checkmark \\
STL.unj & justification removed & \textsf{--} & \textsf{+} & 1 & \textsf{--} & \textsf{--} & G-tst/d $\to$ U & not knowledge \checkmark \\
STL.rel & a hearer believes on the subject's sincere testimony & \textsf{+} & \textsf{+} & 1 & \textsf{--} & \textsf{--} & G-tst/d (same) & safety transmitted \checkmark \\
STL.lie & relay from a speaker who could easily have lied & \textsf{+} & \textsf{+} & 1 & \textsf{--} & \textsf{--} & G-tst/d (same) & unsafe \checkmark \\
\end{longtable}}


\section{Profiles and rival conditions}\label{app:c}

\paragraph{What is computed.} For each of the \stnmodels{} models (cases and variants), the table gives the route, the threshold of the case, the error distance, and the profile: for $\theta = 0, 1, 2, 3, 4$ in turn, \texttt{S} if the belief is basis-wide safe at that threshold and \texttt{u} if not. The program checks that every profile has the form $\texttt{S}^k\texttt{u}^{5-k}$ with $k = \min(\delta, 5)$ (Lemma~\ref{l:delta} and Proposition~\ref{p:margin}(a)). It then gives the verdict on knowledge of justified true belief with each of four added conditions and with safety: NFL (the basis relies on no false lemma), CAU (the truth-maker is causally connected to the basis), SEN (in the nearest worlds in which $p$ is false, the subject does not believe it on the same basis; vacuously satisfied if $p$ is false at no world of the model), UND (no defeater exists), and SAF (justified and basis-wide safe, the analysis). The last columns give the source's verdict where there is one, and the class.

{\scriptsize\setlength{\tabcolsep}{3pt}
\begin{longtable}{@{}llcccccccccl@{}}
\toprule
Id & Route & $\theta$ & $\delta$ & Profile & NFL & CAU & SEN & UND & SAF & Lit. & Class \\
\midrule\endhead
\bottomrule\endfoot
GE1 & infer & 2 & 1 & \texttt{Suuuu} & \textsf{--} & \textsf{--} & \textsf{--} & \textsf{--} & \textsf{--} & N & G-fl/d \\
GE2 & infer & 2 & 1 & \texttt{Suuuu} & \textsf{--} & \textsf{--} & \textsf{--} & \textsf{--} & \textsf{--} & N & G-fl/d \\
NOG & infer & 2 & 1 & \texttt{Suuuu} & \textsf{--} & \textsf{--} & \textsf{--} & \textsf{--} & \textsf{--} & N & G-fl/d \\
FEL & infer & 2 & 1 & \texttt{Suuuu} & \textsf{+} & \textsf{--} & \textsf{--} & \textsf{--} & \textsf{--} & N & G-nfl/d \\
CLK & instr & 2 & 1 & \texttt{Suuuu} & \textsf{--} & \textsf{--} & \textsf{--} & \textsf{--} & \textsf{--} & N & G-tmp/d \\
BRN & perce & 2 & 1 & \texttt{Suuuu} & \textsf{+} & \textsf{+} & \textsf{--} & \textsf{--} & \textsf{--} & N & G-env/c \\
RED & perce & 2 & 3 & \texttt{SSSuu} & \textsf{+} & \textsf{+} & \textsf{+} & \textsf{+} & \textsf{+} & K & K \\
RDP & perce & 2 & 1 & \texttt{Suuuu} & \textsf{+} & \textsf{+} & \textsf{--} & \textsf{--} & \textsf{--} & N & G-env/c \\
SHP & perce & 2 & 1 & \texttt{Suuuu} & \textsf{--} & \textsf{--} & \textsf{--} & \textsf{--} & \textsf{--} & N & G-int/d \\
ZDR & infer & 2 & 1 & \texttt{Suuuu} & \textsf{--} & \textsf{--} & \textsf{--} & \textsf{--} & \textsf{--} & N & G-fl/d \\
ZHB & perce & 2 & 1 & \texttt{Suuuu} & \textsf{--} & \textsf{--} & \textsf{--} & \textsf{--} & \textsf{--} & N & G-int/d \\
PYR & induc & 2 & 1 & \texttt{Suuuu} & \textsf{--} & \textsf{--} & \textsf{--} & \textsf{--} & \textsf{--} & N & G-fl/d \\
GRB & perce & 2 & 4 & \texttt{SSSSu} & \textsf{+} & \textsf{+} & \textsf{--} & \textsf{--} & \textsf{+} & K & K[md] \\
HAR & testi & 2 & $\infty$ & \texttt{SSSSS} & \textsf{+} & \textsf{+} & \textsf{+} & \textsf{--} & \textsf{+} & N & K[md] \\
CLV & facul & 2 & $\infty$ & \texttt{SSSSS} & \textsf{--} & \textsf{--} & \textsf{--} & \textsf{--} & \textsf{--} & N & U \\
TTP & facul & 2 & $\infty$ & \texttt{SSSSS} & \textsf{--} & \textsf{--} & \textsf{--} & \textsf{--} & \textsf{--} & N & U \\
CHU & induc & 2 & 4 & \texttt{SSSSu} & \textsf{+} & \textsf{+} & \textsf{--} & \textsf{+} & \textsf{+} & K & K \\
KAL & perce & 2 & 1 & \texttt{Suuuu} & \textsf{+} & \textsf{+} & \textsf{--} & \textsf{--} & \textsf{--} & K & G-env/c \\
TMP & instr & 2 & 3 & \texttt{SSSuu} & \textsf{+} & \textsf{--} & \textsf{--} & \textsf{+} & \textsf{+} & N & K[fx] \\
MOR & testi & 2 & 3 & \texttt{SSSuu} & \textsf{+} & \textsf{+} & \textsf{--} & \textsf{+} & \textsf{+} & K & K \\
LOT & stati & 2 & 1 & \texttt{Suuuu} & \textsf{+} & \textsf{--} & \textsf{--} & \textsf{+} & \textsf{--} & N & G-lot/d \\
LTN & testi & 2 & 3 & \texttt{SSSuu} & \textsf{+} & \textsf{+} & \textsf{+} & \textsf{+} & \textsf{+} & K & K \\
MAG & perce & 2 & 1 & \texttt{Suuuu} & \textsf{+} & \textsf{+} & \textsf{--} & \textsf{+} & \textsf{--} & N & G-mrg/c \\
MGW & perce & 2 & $\infty$ & \texttt{SSSSS} & \textsf{+} & \textsf{+} & \textsf{+} & \textsf{+} & \textsf{+} & K & K \\
ZBP & perce & 2 & 4 & \texttt{SSSSu} & \textsf{+} & \textsf{+} & \textsf{--} & \textsf{+} & \textsf{+} & K & K \\
ZBM & deduc & 2 & 4 & \texttt{SSSSu} & \textsf{+} & \textsf{+} & \textsf{--} & \textsf{+} & \textsf{+} & N & K \\
GRN & perce & 2 & $\infty$ & \texttt{SSSSS} & \textsf{+} & \textsf{+} & \textsf{+} & \textsf{+} & \textsf{+} & K & K[ev] \\
HLF & testi & 2 & $\infty$ & \texttt{SSSSS} & \textsf{+} & \textsf{+} & \textsf{+} & \textsf{+} & \textsf{+} & K & K \\
HLC & testi & 2 & 1 & \texttt{Suuuu} & \textsf{+} & \textsf{+} & \textsf{--} & \textsf{--} & \textsf{--} & K & G-tst/c \\
ATF & instr & 2 & $\infty$ & \texttt{SSSSS} & \textsf{+} & \textsf{+} & \textsf{+} & \textsf{+} & \textsf{+} & K & K \\
ATC & instr & 2 & 1 & \texttt{Suuuu} & \textsf{+} & \textsf{+} & \textsf{--} & \textsf{--} & \textsf{--} & K & G-env/c \\
WAR & infer & 2 & $\infty$ & \texttt{SSSSS} & \textsf{--} & \textsf{+} & \textsf{+} & \textsf{+} & \textsf{+} & K & K \\
BKL & induc & 2 & 3 & \texttt{SSSuu} & \textsf{+} & \textsf{+} & \textsf{--} & \textsf{+} & \textsf{+} & K & K \\
BKH & induc & 3 & 3 & \texttt{SSSuu} & \textsf{+} & \textsf{+} & \textsf{--} & \textsf{+} & \textsf{--} & N & G-nfl/c \\
BIV & perce & 2 & 4 & \texttt{SSSSu} & \textsf{+} & \textsf{+} & \textsf{--} & \textsf{+} & \textsf{+} & K & K \\
JJM & perce & 2 & 4 & \texttt{SSSSu} & \textsf{+} & \textsf{+} & \textsf{--} & \textsf{+} & \textsf{+} & K & K[ev] \\
THT & testi & 2 & 1 & \texttt{Suuuu} & \textsf{--} & \textsf{--} & \textsf{--} & \textsf{--} & \textsf{--} & N & U \\
ICE & induc & 2 & 4 & \texttt{SSSSu} & \textsf{+} & \textsf{+} & \textsf{--} & \textsf{+} & \textsf{+} & K & K \\
SNC & calcu & 2 & 1 & \texttt{Suuuu} & \textsf{+} & \textsf{--} & \textsf{+} & \textsf{+} & \textsf{--} & -- & G-nec/d \\
SNS & calcu & 2 & $\infty$ & \texttt{SSSSS} & \textsf{+} & \textsf{--} & \textsf{+} & \textsf{+} & \textsf{+} & -- & K \\
SEV & perce & 2 & 4 & \texttt{SSSSu} & \textsf{+} & \textsf{+} & \textsf{--} & \textsf{+} & \textsf{+} & -- & K[ev] \\
SCT & testi & 2 & 4 & \texttt{SSSSu} & \textsf{+} & \textsf{+} & \textsf{+} & \textsf{+} & \textsf{+} & -- & K[ct] \\
SRF & perce & 2 & $\infty$ & \texttt{SSSSS} & \textsf{+} & \textsf{+} & \textsf{+} & \textsf{+} & \textsf{+} & -- & K[rf] \\
STL & testi & 2 & 1 & \texttt{Suuuu} & \textsf{+} & \textsf{--} & \textsf{--} & \textsf{--} & \textsf{--} & -- & G-tst/d \\
GE1.nfl & infer & 2 & 1 & \texttt{Suuuu} & \textsf{+} & \textsf{--} & \textsf{--} & \textsf{--} & \textsf{--} & -- & G-nfl/d \\
GE1.sec & infer & 2 & $\infty$ & \texttt{SSSSS} & \textsf{--} & \textsf{+} & \textsf{+} & \textsf{--} & \textsf{+} & -- & K \\
GE1.rem & infer & 2 & 4 & \texttt{SSSSu} & \textsf{--} & \textsf{--} & \textsf{--} & \textsf{--} & \textsf{+} & -- & K \\
GE1.fal & infer & 2 & 0 & \texttt{uuuuu} & \textsf{--} & \textsf{--} & \textsf{--} & \textsf{--} & \textsf{--} & -- & F \\
GE1.unj & infer & 2 & 1 & \texttt{Suuuu} & \textsf{--} & \textsf{--} & \textsf{--} & \textsf{--} & \textsf{--} & -- & U \\
GE1.rel & testi & 2 & 1 & \texttt{Suuuu} & \textsf{--} & \textsf{--} & \textsf{--} & \textsf{--} & \textsf{--} & -- & G-int/d \\
GE1.lie & testi & 2 & 1 & \texttt{Suuuu} & \textsf{--} & \textsf{--} & \textsf{--} & \textsf{--} & \textsf{--} & -- & G-tst/d \\
GE1.ded & deduc & 2 & $\infty$ & \texttt{SSSSS} & \textsf{--} & \textsf{--} & \textsf{+} & \textsf{--} & \textsf{+} & -- & K \\
GE2.nfl & infer & 2 & 1 & \texttt{Suuuu} & \textsf{+} & \textsf{--} & \textsf{--} & \textsf{--} & \textsf{--} & -- & G-nfl/d \\
GE2.sec & infer & 2 & $\infty$ & \texttt{SSSSS} & \textsf{--} & \textsf{+} & \textsf{+} & \textsf{--} & \textsf{+} & -- & K \\
GE2.rem & infer & 2 & 4 & \texttt{SSSSu} & \textsf{--} & \textsf{--} & \textsf{--} & \textsf{--} & \textsf{+} & -- & K \\
GE2.fal & infer & 2 & 0 & \texttt{uuuuu} & \textsf{--} & \textsf{--} & \textsf{--} & \textsf{--} & \textsf{--} & -- & F \\
GE2.unj & infer & 2 & 1 & \texttt{Suuuu} & \textsf{--} & \textsf{--} & \textsf{--} & \textsf{--} & \textsf{--} & -- & U \\
GE2.rel & testi & 2 & 1 & \texttt{Suuuu} & \textsf{--} & \textsf{--} & \textsf{--} & \textsf{--} & \textsf{--} & -- & G-int/d \\
GE2.lie & testi & 2 & 1 & \texttt{Suuuu} & \textsf{--} & \textsf{--} & \textsf{--} & \textsf{--} & \textsf{--} & -- & G-tst/d \\
NOG.nfl & infer & 2 & 1 & \texttt{Suuuu} & \textsf{+} & \textsf{--} & \textsf{--} & \textsf{--} & \textsf{--} & -- & G-nfl/d \\
NOG.sec & infer & 2 & $\infty$ & \texttt{SSSSS} & \textsf{--} & \textsf{+} & \textsf{+} & \textsf{--} & \textsf{+} & -- & K \\
NOG.rem & infer & 2 & 4 & \texttt{SSSSu} & \textsf{--} & \textsf{--} & \textsf{--} & \textsf{--} & \textsf{+} & -- & K \\
NOG.fal & infer & 2 & 0 & \texttt{uuuuu} & \textsf{--} & \textsf{--} & \textsf{--} & \textsf{--} & \textsf{--} & -- & F \\
NOG.unj & infer & 2 & 1 & \texttt{Suuuu} & \textsf{--} & \textsf{--} & \textsf{--} & \textsf{--} & \textsf{--} & -- & U \\
NOG.rel & testi & 2 & 1 & \texttt{Suuuu} & \textsf{--} & \textsf{--} & \textsf{--} & \textsf{--} & \textsf{--} & -- & G-int/d \\
NOG.lie & testi & 2 & 1 & \texttt{Suuuu} & \textsf{--} & \textsf{--} & \textsf{--} & \textsf{--} & \textsf{--} & -- & G-tst/d \\
NOG.ded & deduc & 2 & $\infty$ & \texttt{SSSSS} & \textsf{--} & \textsf{--} & \textsf{+} & \textsf{--} & \textsf{+} & -- & K \\
FEL.fl & infer & 2 & 1 & \texttt{Suuuu} & \textsf{--} & \textsf{--} & \textsf{--} & \textsf{--} & \textsf{--} & -- & G-fl/d \\
FEL.sec & infer & 2 & $\infty$ & \texttt{SSSSS} & \textsf{+} & \textsf{+} & \textsf{+} & \textsf{--} & \textsf{+} & -- & K \\
FEL.rem & infer & 2 & 4 & \texttt{SSSSu} & \textsf{+} & \textsf{--} & \textsf{--} & \textsf{--} & \textsf{+} & -- & K \\
FEL.fal & infer & 2 & 0 & \texttt{uuuuu} & \textsf{--} & \textsf{--} & \textsf{--} & \textsf{--} & \textsf{--} & -- & F \\
FEL.unj & infer & 2 & 1 & \texttt{Suuuu} & \textsf{--} & \textsf{--} & \textsf{--} & \textsf{--} & \textsf{--} & -- & U \\
FEL.rel & testi & 2 & 1 & \texttt{Suuuu} & \textsf{+} & \textsf{--} & \textsf{--} & \textsf{--} & \textsf{--} & -- & G-int/d \\
FEL.lie & testi & 2 & 1 & \texttt{Suuuu} & \textsf{+} & \textsf{--} & \textsf{--} & \textsf{--} & \textsf{--} & -- & G-tst/d \\
FEL.ded & deduc & 2 & $\infty$ & \texttt{SSSSS} & \textsf{+} & \textsf{--} & \textsf{+} & \textsf{--} & \textsf{+} & -- & K \\
CLK.nfl & instr & 2 & 1 & \texttt{Suuuu} & \textsf{+} & \textsf{--} & \textsf{--} & \textsf{--} & \textsf{--} & -- & G-tmp/d \\
CLK.sec & instr & 2 & $\infty$ & \texttt{SSSSS} & \textsf{--} & \textsf{+} & \textsf{+} & \textsf{--} & \textsf{+} & -- & K \\
CLK.hos & instr & 2 & 1 & \texttt{Suuuu} & \textsf{--} & \textsf{--} & \textsf{--} & \textsf{--} & \textsf{--} & -- & G-tmp/d \\
CLK.rem & instr & 2 & 4 & \texttt{SSSSu} & \textsf{--} & \textsf{--} & \textsf{--} & \textsf{--} & \textsf{+} & -- & K \\
CLK.fal & instr & 2 & 0 & \texttt{uuuuu} & \textsf{--} & \textsf{--} & \textsf{--} & \textsf{--} & \textsf{--} & -- & F \\
CLK.unj & instr & 2 & 1 & \texttt{Suuuu} & \textsf{--} & \textsf{--} & \textsf{--} & \textsf{--} & \textsf{--} & -- & U \\
CLK.rel & testi & 2 & 1 & \texttt{Suuuu} & \textsf{--} & \textsf{--} & \textsf{--} & \textsf{--} & \textsf{--} & -- & G-tmp/d \\
CLK.lie & testi & 2 & 1 & \texttt{Suuuu} & \textsf{--} & \textsf{--} & \textsf{--} & \textsf{--} & \textsf{--} & -- & G-tmp/d \\
BRN.nor & perce & 2 & $\infty$ & \texttt{SSSSS} & \textsf{+} & \textsf{+} & \textsf{+} & \textsf{--} & \textsf{+} & -- & K \\
BRN.rem & perce & 2 & 4 & \texttt{SSSSu} & \textsf{+} & \textsf{+} & \textsf{--} & \textsf{--} & \textsf{+} & -- & K \\
BRN.fal & perce & 2 & 0 & \texttt{uuuuu} & \textsf{--} & \textsf{--} & \textsf{--} & \textsf{--} & \textsf{--} & -- & F \\
BRN.unj & perce & 2 & 1 & \texttt{Suuuu} & \textsf{--} & \textsf{--} & \textsf{--} & \textsf{--} & \textsf{--} & -- & U \\
BRN.rel & testi & 2 & 1 & \texttt{Suuuu} & \textsf{+} & \textsf{+} & \textsf{--} & \textsf{--} & \textsf{--} & -- & G-env/c \\
BRN.lie & testi & 2 & 1 & \texttt{Suuuu} & \textsf{+} & \textsf{+} & \textsf{--} & \textsf{--} & \textsf{--} & -- & G-tst/c \\
BRN.ded & deduc & 2 & $\infty$ & \texttt{SSSSS} & \textsf{+} & \textsf{+} & \textsf{+} & \textsf{--} & \textsf{+} & -- & K \\
RED.nor & perce & 2 & $\infty$ & \texttt{SSSSS} & \textsf{+} & \textsf{+} & \textsf{+} & \textsf{+} & \textsf{+} & -- & K \\
RED.hos & perce & 2 & 1 & \texttt{Suuuu} & \textsf{+} & \textsf{+} & \textsf{--} & \textsf{+} & \textsf{--} & -- & G-env/c \\
RED.fal & perce & 2 & 0 & \texttt{uuuuu} & \textsf{--} & \textsf{--} & \textsf{--} & \textsf{--} & \textsf{--} & -- & F \\
RED.unj & perce & 2 & 3 & \texttt{SSSuu} & \textsf{--} & \textsf{--} & \textsf{--} & \textsf{--} & \textsf{--} & -- & U \\
RED.rel & testi & 2 & 3 & \texttt{SSSuu} & \textsf{+} & \textsf{+} & \textsf{+} & \textsf{+} & \textsf{+} & -- & K \\
RED.lie & testi & 2 & 1 & \texttt{Suuuu} & \textsf{+} & \textsf{+} & \textsf{--} & \textsf{--} & \textsf{--} & -- & G-tst/c \\
RED.ded & deduc & 2 & 3 & \texttt{SSSuu} & \textsf{+} & \textsf{+} & \textsf{+} & \textsf{+} & \textsf{+} & -- & K \\
RDP.nor & perce & 2 & $\infty$ & \texttt{SSSSS} & \textsf{+} & \textsf{+} & \textsf{+} & \textsf{--} & \textsf{+} & -- & K \\
RDP.rem & perce & 2 & 4 & \texttt{SSSSu} & \textsf{+} & \textsf{+} & \textsf{--} & \textsf{--} & \textsf{+} & -- & K \\
RDP.fal & perce & 2 & 0 & \texttt{uuuuu} & \textsf{--} & \textsf{--} & \textsf{--} & \textsf{--} & \textsf{--} & -- & F \\
RDP.unj & perce & 2 & 1 & \texttt{Suuuu} & \textsf{--} & \textsf{--} & \textsf{--} & \textsf{--} & \textsf{--} & -- & U \\
RDP.rel & testi & 2 & 1 & \texttt{Suuuu} & \textsf{+} & \textsf{+} & \textsf{--} & \textsf{--} & \textsf{--} & -- & G-env/c \\
RDP.lie & testi & 2 & 1 & \texttt{Suuuu} & \textsf{+} & \textsf{+} & \textsf{--} & \textsf{--} & \textsf{--} & -- & G-tst/c \\
SHP.nfl & perce & 2 & 1 & \texttt{Suuuu} & \textsf{+} & \textsf{--} & \textsf{--} & \textsf{--} & \textsf{--} & -- & G-int/d \\
SHP.sec & perce & 2 & $\infty$ & \texttt{SSSSS} & \textsf{--} & \textsf{+} & \textsf{+} & \textsf{--} & \textsf{+} & -- & K \\
SHP.hos & perce & 2 & 1 & \texttt{Suuuu} & \textsf{--} & \textsf{--} & \textsf{--} & \textsf{--} & \textsf{--} & -- & G-env/d \\
SHP.rem & perce & 2 & 4 & \texttt{SSSSu} & \textsf{--} & \textsf{--} & \textsf{--} & \textsf{--} & \textsf{+} & -- & K \\
SHP.fal & perce & 2 & 0 & \texttt{uuuuu} & \textsf{--} & \textsf{--} & \textsf{--} & \textsf{--} & \textsf{--} & -- & F \\
SHP.unj & perce & 2 & 1 & \texttt{Suuuu} & \textsf{--} & \textsf{--} & \textsf{--} & \textsf{--} & \textsf{--} & -- & U \\
SHP.rel & testi & 2 & 1 & \texttt{Suuuu} & \textsf{--} & \textsf{--} & \textsf{--} & \textsf{--} & \textsf{--} & -- & G-int/d \\
SHP.lie & testi & 2 & 1 & \texttt{Suuuu} & \textsf{--} & \textsf{--} & \textsf{--} & \textsf{--} & \textsf{--} & -- & G-tst/d \\
ZDR.nfl & infer & 2 & 1 & \texttt{Suuuu} & \textsf{+} & \textsf{--} & \textsf{--} & \textsf{--} & \textsf{--} & -- & G-nfl/d \\
ZDR.sec & infer & 2 & $\infty$ & \texttt{SSSSS} & \textsf{--} & \textsf{+} & \textsf{+} & \textsf{--} & \textsf{+} & -- & K \\
ZDR.rem & infer & 2 & 4 & \texttt{SSSSu} & \textsf{--} & \textsf{--} & \textsf{--} & \textsf{--} & \textsf{+} & -- & K \\
ZDR.fal & infer & 2 & 0 & \texttt{uuuuu} & \textsf{--} & \textsf{--} & \textsf{--} & \textsf{--} & \textsf{--} & -- & F \\
ZDR.unj & infer & 2 & 1 & \texttt{Suuuu} & \textsf{--} & \textsf{--} & \textsf{--} & \textsf{--} & \textsf{--} & -- & U \\
ZDR.rel & testi & 2 & 1 & \texttt{Suuuu} & \textsf{--} & \textsf{--} & \textsf{--} & \textsf{--} & \textsf{--} & -- & G-int/d \\
ZDR.lie & testi & 2 & 1 & \texttt{Suuuu} & \textsf{--} & \textsf{--} & \textsf{--} & \textsf{--} & \textsf{--} & -- & G-tst/d \\
ZHB.nfl & perce & 2 & 1 & \texttt{Suuuu} & \textsf{+} & \textsf{--} & \textsf{--} & \textsf{--} & \textsf{--} & -- & G-int/d \\
ZHB.sec & perce & 2 & $\infty$ & \texttt{SSSSS} & \textsf{--} & \textsf{+} & \textsf{+} & \textsf{--} & \textsf{+} & -- & K \\
ZHB.hos & perce & 2 & 1 & \texttt{Suuuu} & \textsf{--} & \textsf{--} & \textsf{--} & \textsf{--} & \textsf{--} & -- & G-env/d \\
ZHB.rem & perce & 2 & 4 & \texttt{SSSSu} & \textsf{--} & \textsf{--} & \textsf{--} & \textsf{--} & \textsf{+} & -- & K \\
ZHB.fal & perce & 2 & 0 & \texttt{uuuuu} & \textsf{--} & \textsf{--} & \textsf{--} & \textsf{--} & \textsf{--} & -- & F \\
ZHB.unj & perce & 2 & 1 & \texttt{Suuuu} & \textsf{--} & \textsf{--} & \textsf{--} & \textsf{--} & \textsf{--} & -- & U \\
ZHB.rel & testi & 2 & 1 & \texttt{Suuuu} & \textsf{--} & \textsf{--} & \textsf{--} & \textsf{--} & \textsf{--} & -- & G-int/d \\
ZHB.lie & testi & 2 & 1 & \texttt{Suuuu} & \textsf{--} & \textsf{--} & \textsf{--} & \textsf{--} & \textsf{--} & -- & G-tst/d \\
PYR.nfl & induc & 2 & 1 & \texttt{Suuuu} & \textsf{+} & \textsf{--} & \textsf{--} & \textsf{--} & \textsf{--} & -- & G-nfl/d \\
PYR.sec & induc & 2 & $\infty$ & \texttt{SSSSS} & \textsf{--} & \textsf{+} & \textsf{+} & \textsf{--} & \textsf{+} & -- & K \\
PYR.rem & induc & 2 & 4 & \texttt{SSSSu} & \textsf{--} & \textsf{--} & \textsf{--} & \textsf{--} & \textsf{+} & -- & K \\
PYR.fal & induc & 2 & 0 & \texttt{uuuuu} & \textsf{--} & \textsf{--} & \textsf{--} & \textsf{--} & \textsf{--} & -- & F \\
PYR.unj & induc & 2 & 1 & \texttt{Suuuu} & \textsf{--} & \textsf{--} & \textsf{--} & \textsf{--} & \textsf{--} & -- & U \\
PYR.rel & testi & 2 & 1 & \texttt{Suuuu} & \textsf{--} & \textsf{--} & \textsf{--} & \textsf{--} & \textsf{--} & -- & G-int/d \\
PYR.lie & testi & 2 & 1 & \texttt{Suuuu} & \textsf{--} & \textsf{--} & \textsf{--} & \textsf{--} & \textsf{--} & -- & G-tst/d \\
GRB.nor & perce & 2 & $\infty$ & \texttt{SSSSS} & \textsf{+} & \textsf{+} & \textsf{+} & \textsf{--} & \textsf{+} & -- & K[md] \\
GRB.hos & perce & 2 & 1 & \texttt{Suuuu} & \textsf{+} & \textsf{+} & \textsf{--} & \textsf{--} & \textsf{--} & -- & G-env/c \\
GRB.fal & perce & 2 & 0 & \texttt{uuuuu} & \textsf{--} & \textsf{--} & \textsf{--} & \textsf{--} & \textsf{--} & -- & F \\
GRB.unj & perce & 2 & 4 & \texttt{SSSSu} & \textsf{--} & \textsf{--} & \textsf{--} & \textsf{--} & \textsf{--} & -- & U \\
GRB.rel & testi & 2 & 4 & \texttt{SSSSu} & \textsf{+} & \textsf{+} & \textsf{--} & \textsf{--} & \textsf{+} & -- & K[md] \\
GRB.lie & testi & 2 & 1 & \texttt{Suuuu} & \textsf{+} & \textsf{+} & \textsf{--} & \textsf{--} & \textsf{--} & -- & G-tst/c \\
HAR.hos & testi & 2 & 1 & \texttt{Suuuu} & \textsf{+} & \textsf{+} & \textsf{--} & \textsf{--} & \textsf{--} & -- & G-env/c \\
HAR.fal & testi & 2 & 0 & \texttt{uuuuu} & \textsf{--} & \textsf{--} & \textsf{--} & \textsf{--} & \textsf{--} & -- & F \\
HAR.unj & testi & 2 & $\infty$ & \texttt{SSSSS} & \textsf{--} & \textsf{--} & \textsf{--} & \textsf{--} & \textsf{--} & -- & U \\
HAR.rel & testi & 2 & $\infty$ & \texttt{SSSSS} & \textsf{+} & \textsf{+} & \textsf{+} & \textsf{--} & \textsf{+} & -- & K[md] \\
HAR.lie & testi & 2 & 1 & \texttt{Suuuu} & \textsf{+} & \textsf{+} & \textsf{--} & \textsf{--} & \textsf{--} & -- & G-tst/c \\
CLV.fal & facul & 2 & 0 & \texttt{uuuuu} & \textsf{--} & \textsf{--} & \textsf{--} & \textsf{--} & \textsf{--} & -- & U \\
CLV.rel & testi & 2 & $\infty$ & \texttt{SSSSS} & \textsf{--} & \textsf{--} & \textsf{--} & \textsf{--} & \textsf{--} & -- & U \\
CLV.lie & testi & 2 & 1 & \texttt{Suuuu} & \textsf{--} & \textsf{--} & \textsf{--} & \textsf{--} & \textsf{--} & -- & U \\
TTP.fal & facul & 2 & 0 & \texttt{uuuuu} & \textsf{--} & \textsf{--} & \textsf{--} & \textsf{--} & \textsf{--} & -- & U \\
TTP.rel & testi & 2 & $\infty$ & \texttt{SSSSS} & \textsf{--} & \textsf{--} & \textsf{--} & \textsf{--} & \textsf{--} & -- & U \\
TTP.lie & testi & 2 & 1 & \texttt{Suuuu} & \textsf{--} & \textsf{--} & \textsf{--} & \textsf{--} & \textsf{--} & -- & U \\
CHU.fl & induc & 2 & 4 & \texttt{SSSSu} & \textsf{--} & \textsf{+} & \textsf{--} & \textsf{+} & \textsf{+} & -- & K \\
CHU.sec & induc & 2 & $\infty$ & \texttt{SSSSS} & \textsf{+} & \textsf{+} & \textsf{+} & \textsf{+} & \textsf{+} & -- & K \\
CHU.fal & induc & 2 & 0 & \texttt{uuuuu} & \textsf{--} & \textsf{--} & \textsf{--} & \textsf{--} & \textsf{--} & -- & F \\
CHU.unj & induc & 2 & 4 & \texttt{SSSSu} & \textsf{--} & \textsf{--} & \textsf{--} & \textsf{--} & \textsf{--} & -- & U \\
CHU.rel & testi & 2 & 4 & \texttt{SSSSu} & \textsf{+} & \textsf{+} & \textsf{--} & \textsf{+} & \textsf{+} & -- & K \\
CHU.lie & testi & 2 & 1 & \texttt{Suuuu} & \textsf{+} & \textsf{+} & \textsf{--} & \textsf{--} & \textsf{--} & -- & G-tst/c \\
KAL.nor & perce & 2 & $\infty$ & \texttt{SSSSS} & \textsf{+} & \textsf{+} & \textsf{+} & \textsf{--} & \textsf{+} & -- & K[rf] \\
KAL.rem & perce & 2 & 4 & \texttt{SSSSu} & \textsf{+} & \textsf{+} & \textsf{--} & \textsf{--} & \textsf{+} & -- & K[rf] \\
KAL.fal & perce & 2 & 0 & \texttt{uuuuu} & \textsf{--} & \textsf{--} & \textsf{--} & \textsf{--} & \textsf{--} & -- & F \\
KAL.unj & perce & 2 & 1 & \texttt{Suuuu} & \textsf{--} & \textsf{--} & \textsf{--} & \textsf{--} & \textsf{--} & -- & U \\
KAL.rel & testi & 2 & 1 & \texttt{Suuuu} & \textsf{+} & \textsf{+} & \textsf{--} & \textsf{--} & \textsf{--} & -- & G-env/c \\
KAL.lie & testi & 2 & 1 & \texttt{Suuuu} & \textsf{+} & \textsf{+} & \textsf{--} & \textsf{--} & \textsf{--} & -- & G-tst/c \\
TMP.sec & instr & 2 & $\infty$ & \texttt{SSSSS} & \textsf{+} & \textsf{+} & \textsf{+} & \textsf{+} & \textsf{+} & -- & K[fx] \\
TMP.hos & instr & 2 & 1 & \texttt{Suuuu} & \textsf{+} & \textsf{--} & \textsf{--} & \textsf{+} & \textsf{--} & -- & G-env/d \\
TMP.fal & instr & 2 & 0 & \texttt{uuuuu} & \textsf{--} & \textsf{--} & \textsf{--} & \textsf{--} & \textsf{--} & -- & F \\
TMP.unj & instr & 2 & 3 & \texttt{SSSuu} & \textsf{--} & \textsf{--} & \textsf{--} & \textsf{--} & \textsf{--} & -- & U \\
TMP.rel & testi & 2 & 3 & \texttt{SSSuu} & \textsf{+} & \textsf{--} & \textsf{--} & \textsf{+} & \textsf{+} & -- & K[fx] \\
TMP.lie & testi & 2 & 1 & \texttt{Suuuu} & \textsf{+} & \textsf{--} & \textsf{--} & \textsf{--} & \textsf{--} & -- & G-tst/d \\
MOR.nor & testi & 2 & $\infty$ & \texttt{SSSSS} & \textsf{+} & \textsf{+} & \textsf{+} & \textsf{+} & \textsf{+} & -- & K \\
MOR.hos & testi & 2 & 1 & \texttt{Suuuu} & \textsf{+} & \textsf{+} & \textsf{--} & \textsf{+} & \textsf{--} & -- & G-env/c \\
MOR.fal & testi & 2 & 0 & \texttt{uuuuu} & \textsf{--} & \textsf{--} & \textsf{--} & \textsf{--} & \textsf{--} & -- & F \\
MOR.unj & testi & 2 & 3 & \texttt{SSSuu} & \textsf{--} & \textsf{--} & \textsf{--} & \textsf{--} & \textsf{--} & -- & U \\
MOR.rel & testi & 2 & 3 & \texttt{SSSuu} & \textsf{+} & \textsf{+} & \textsf{--} & \textsf{+} & \textsf{+} & -- & K \\
MOR.lie & testi & 2 & 1 & \texttt{Suuuu} & \textsf{+} & \textsf{+} & \textsf{--} & \textsf{--} & \textsf{--} & -- & G-tst/c \\
LOT.sec & stati & 2 & $\infty$ & \texttt{SSSSS} & \textsf{+} & \textsf{+} & \textsf{+} & \textsf{+} & \textsf{+} & -- & K \\
LOT.rem & stati & 2 & 4 & \texttt{SSSSu} & \textsf{+} & \textsf{--} & \textsf{--} & \textsf{+} & \textsf{+} & -- & K \\
LOT.fal & stati & 2 & 0 & \texttt{uuuuu} & \textsf{--} & \textsf{--} & \textsf{--} & \textsf{--} & \textsf{--} & -- & F \\
LOT.unj & stati & 2 & 1 & \texttt{Suuuu} & \textsf{--} & \textsf{--} & \textsf{--} & \textsf{--} & \textsf{--} & -- & U \\
LOT.rel & testi & 2 & 1 & \texttt{Suuuu} & \textsf{+} & \textsf{--} & \textsf{--} & \textsf{+} & \textsf{--} & -- & G-lot/d \\
LOT.lie & testi & 2 & 1 & \texttt{Suuuu} & \textsf{+} & \textsf{--} & \textsf{--} & \textsf{--} & \textsf{--} & -- & G-lot/d \\
LTN.nor & testi & 2 & $\infty$ & \texttt{SSSSS} & \textsf{+} & \textsf{+} & \textsf{+} & \textsf{+} & \textsf{+} & -- & K \\
LTN.hos & testi & 2 & 1 & \texttt{Suuuu} & \textsf{+} & \textsf{+} & \textsf{--} & \textsf{+} & \textsf{--} & -- & G-env/c \\
LTN.fal & testi & 2 & 0 & \texttt{uuuuu} & \textsf{--} & \textsf{--} & \textsf{--} & \textsf{--} & \textsf{--} & -- & F \\
LTN.unj & testi & 2 & 3 & \texttt{SSSuu} & \textsf{--} & \textsf{--} & \textsf{--} & \textsf{--} & \textsf{--} & -- & U \\
LTN.rel & testi & 2 & 3 & \texttt{SSSuu} & \textsf{+} & \textsf{+} & \textsf{+} & \textsf{+} & \textsf{+} & -- & K \\
LTN.lie & testi & 2 & 1 & \texttt{Suuuu} & \textsf{+} & \textsf{+} & \textsf{--} & \textsf{--} & \textsf{--} & -- & G-tst/c \\
MAG.hos & perce & 2 & 1 & \texttt{Suuuu} & \textsf{+} & \textsf{+} & \textsf{--} & \textsf{+} & \textsf{--} & -- & G-mrg/c \\
MAG.rem & perce & 2 & 4 & \texttt{SSSSu} & \textsf{+} & \textsf{+} & \textsf{--} & \textsf{+} & \textsf{+} & -- & K \\
MAG.fal & perce & 2 & 0 & \texttt{uuuuu} & \textsf{--} & \textsf{--} & \textsf{--} & \textsf{--} & \textsf{--} & -- & F \\
MAG.unj & perce & 2 & 1 & \texttt{Suuuu} & \textsf{--} & \textsf{--} & \textsf{--} & \textsf{--} & \textsf{--} & -- & U \\
MAG.rel & testi & 2 & 1 & \texttt{Suuuu} & \textsf{+} & \textsf{+} & \textsf{--} & \textsf{+} & \textsf{--} & -- & G-mrg/c \\
MAG.lie & testi & 2 & 1 & \texttt{Suuuu} & \textsf{+} & \textsf{+} & \textsf{--} & \textsf{--} & \textsf{--} & -- & G-mrg/c \\
MGW.hos & perce & 2 & 1 & \texttt{Suuuu} & \textsf{+} & \textsf{+} & \textsf{--} & \textsf{+} & \textsf{--} & -- & G-env/c \\
MGW.fal & perce & 2 & 0 & \texttt{uuuuu} & \textsf{--} & \textsf{--} & \textsf{--} & \textsf{--} & \textsf{--} & -- & F \\
MGW.unj & perce & 2 & $\infty$ & \texttt{SSSSS} & \textsf{--} & \textsf{--} & \textsf{--} & \textsf{--} & \textsf{--} & -- & U \\
MGW.rel & testi & 2 & $\infty$ & \texttt{SSSSS} & \textsf{+} & \textsf{+} & \textsf{+} & \textsf{+} & \textsf{+} & -- & K \\
MGW.lie & testi & 2 & 1 & \texttt{Suuuu} & \textsf{+} & \textsf{+} & \textsf{--} & \textsf{--} & \textsf{--} & -- & G-tst/c \\
ZBP.nor & perce & 2 & $\infty$ & \texttt{SSSSS} & \textsf{+} & \textsf{+} & \textsf{+} & \textsf{+} & \textsf{+} & -- & K \\
ZBP.hos & perce & 2 & 1 & \texttt{Suuuu} & \textsf{+} & \textsf{+} & \textsf{--} & \textsf{+} & \textsf{--} & -- & G-env/c \\
ZBP.fal & perce & 2 & 0 & \texttt{uuuuu} & \textsf{--} & \textsf{--} & \textsf{--} & \textsf{--} & \textsf{--} & -- & F \\
ZBP.unj & perce & 2 & 4 & \texttt{SSSSu} & \textsf{--} & \textsf{--} & \textsf{--} & \textsf{--} & \textsf{--} & -- & U \\
ZBP.rel & testi & 2 & 4 & \texttt{SSSSu} & \textsf{+} & \textsf{+} & \textsf{--} & \textsf{+} & \textsf{+} & -- & K \\
ZBP.lie & testi & 2 & 1 & \texttt{Suuuu} & \textsf{+} & \textsf{+} & \textsf{--} & \textsf{--} & \textsf{--} & -- & G-tst/c \\
ZBP.ded & deduc & 2 & 4 & \texttt{SSSSu} & \textsf{+} & \textsf{+} & \textsf{--} & \textsf{+} & \textsf{+} & -- & K \\
ZBM.fl & deduc & 2 & 4 & \texttt{SSSSu} & \textsf{--} & \textsf{+} & \textsf{--} & \textsf{+} & \textsf{+} & -- & K \\
ZBM.nor & deduc & 2 & $\infty$ & \texttt{SSSSS} & \textsf{+} & \textsf{+} & \textsf{+} & \textsf{+} & \textsf{+} & -- & K \\
ZBM.fal & deduc & 2 & 0 & \texttt{uuuuu} & \textsf{--} & \textsf{--} & \textsf{--} & \textsf{--} & \textsf{--} & -- & F \\
ZBM.unj & deduc & 2 & 4 & \texttt{SSSSu} & \textsf{--} & \textsf{--} & \textsf{--} & \textsf{--} & \textsf{--} & -- & U \\
ZBM.rel & testi & 2 & 4 & \texttt{SSSSu} & \textsf{+} & \textsf{+} & \textsf{--} & \textsf{+} & \textsf{+} & -- & K \\
ZBM.lie & testi & 2 & 1 & \texttt{Suuuu} & \textsf{+} & \textsf{+} & \textsf{--} & \textsf{--} & \textsf{--} & -- & G-tst/c \\
GRN.hos & perce & 2 & 1 & \texttt{Suuuu} & \textsf{+} & \textsf{+} & \textsf{--} & \textsf{+} & \textsf{--} & -- & G-env/c \\
GRN.fal & perce & 2 & 0 & \texttt{uuuuu} & \textsf{--} & \textsf{--} & \textsf{--} & \textsf{--} & \textsf{--} & -- & F \\
GRN.unj & perce & 2 & $\infty$ & \texttt{SSSSS} & \textsf{--} & \textsf{--} & \textsf{--} & \textsf{--} & \textsf{--} & -- & U \\
GRN.rel & testi & 2 & $\infty$ & \texttt{SSSSS} & \textsf{+} & \textsf{+} & \textsf{+} & \textsf{+} & \textsf{+} & -- & K[ev] \\
GRN.lie & testi & 2 & 1 & \texttt{Suuuu} & \textsf{+} & \textsf{+} & \textsf{--} & \textsf{--} & \textsf{--} & -- & G-tst/c \\
HLF.hos & testi & 2 & 1 & \texttt{Suuuu} & \textsf{+} & \textsf{+} & \textsf{--} & \textsf{+} & \textsf{--} & -- & G-env/c \\
HLF.fal & testi & 2 & 0 & \texttt{uuuuu} & \textsf{--} & \textsf{--} & \textsf{--} & \textsf{--} & \textsf{--} & -- & F \\
HLF.unj & testi & 2 & $\infty$ & \texttt{SSSSS} & \textsf{--} & \textsf{--} & \textsf{--} & \textsf{--} & \textsf{--} & -- & U \\
HLF.rel & testi & 2 & $\infty$ & \texttt{SSSSS} & \textsf{+} & \textsf{+} & \textsf{+} & \textsf{+} & \textsf{+} & -- & K \\
HLF.lie & testi & 2 & 1 & \texttt{Suuuu} & \textsf{+} & \textsf{+} & \textsf{--} & \textsf{--} & \textsf{--} & -- & G-tst/c \\
HLC.nor & testi & 2 & $\infty$ & \texttt{SSSSS} & \textsf{+} & \textsf{+} & \textsf{+} & \textsf{--} & \textsf{+} & -- & K \\
HLC.rem & testi & 2 & 4 & \texttt{SSSSu} & \textsf{+} & \textsf{+} & \textsf{--} & \textsf{--} & \textsf{+} & -- & K \\
HLC.fal & testi & 2 & 0 & \texttt{uuuuu} & \textsf{--} & \textsf{--} & \textsf{--} & \textsf{--} & \textsf{--} & -- & F \\
HLC.unj & testi & 2 & 1 & \texttt{Suuuu} & \textsf{--} & \textsf{--} & \textsf{--} & \textsf{--} & \textsf{--} & -- & U \\
HLC.rel & testi & 2 & 1 & \texttt{Suuuu} & \textsf{+} & \textsf{+} & \textsf{--} & \textsf{--} & \textsf{--} & -- & G-tst/c \\
HLC.lie & testi & 2 & 1 & \texttt{Suuuu} & \textsf{+} & \textsf{+} & \textsf{--} & \textsf{--} & \textsf{--} & -- & G-tst/c \\
ATF.hos & instr & 2 & 1 & \texttt{Suuuu} & \textsf{+} & \textsf{+} & \textsf{--} & \textsf{+} & \textsf{--} & -- & G-env/c \\
ATF.fal & instr & 2 & 0 & \texttt{uuuuu} & \textsf{--} & \textsf{--} & \textsf{--} & \textsf{--} & \textsf{--} & -- & F \\
ATF.unj & instr & 2 & $\infty$ & \texttt{SSSSS} & \textsf{--} & \textsf{--} & \textsf{--} & \textsf{--} & \textsf{--} & -- & U \\
ATF.rel & testi & 2 & $\infty$ & \texttt{SSSSS} & \textsf{+} & \textsf{+} & \textsf{+} & \textsf{+} & \textsf{+} & -- & K \\
ATF.lie & testi & 2 & 1 & \texttt{Suuuu} & \textsf{+} & \textsf{+} & \textsf{--} & \textsf{--} & \textsf{--} & -- & G-tst/c \\
ATC.nor & instr & 2 & $\infty$ & \texttt{SSSSS} & \textsf{+} & \textsf{+} & \textsf{+} & \textsf{--} & \textsf{+} & -- & K \\
ATC.rem & instr & 2 & 4 & \texttt{SSSSu} & \textsf{+} & \textsf{+} & \textsf{--} & \textsf{--} & \textsf{+} & -- & K \\
ATC.fal & instr & 2 & 0 & \texttt{uuuuu} & \textsf{--} & \textsf{--} & \textsf{--} & \textsf{--} & \textsf{--} & -- & F \\
ATC.unj & instr & 2 & 1 & \texttt{Suuuu} & \textsf{--} & \textsf{--} & \textsf{--} & \textsf{--} & \textsf{--} & -- & U \\
ATC.rel & testi & 2 & 1 & \texttt{Suuuu} & \textsf{+} & \textsf{+} & \textsf{--} & \textsf{--} & \textsf{--} & -- & G-env/c \\
ATC.lie & testi & 2 & 1 & \texttt{Suuuu} & \textsf{+} & \textsf{+} & \textsf{--} & \textsf{--} & \textsf{--} & -- & G-tst/c \\
WAR.nfl & infer & 2 & $\infty$ & \texttt{SSSSS} & \textsf{+} & \textsf{+} & \textsf{+} & \textsf{+} & \textsf{+} & -- & K \\
WAR.fal & infer & 2 & 0 & \texttt{uuuuu} & \textsf{--} & \textsf{--} & \textsf{--} & \textsf{--} & \textsf{--} & -- & F \\
WAR.unj & infer & 2 & $\infty$ & \texttt{SSSSS} & \textsf{--} & \textsf{--} & \textsf{--} & \textsf{--} & \textsf{--} & -- & U \\
WAR.rel & testi & 2 & $\infty$ & \texttt{SSSSS} & \textsf{--} & \textsf{+} & \textsf{+} & \textsf{+} & \textsf{+} & -- & K \\
WAR.lie & testi & 2 & 1 & \texttt{Suuuu} & \textsf{--} & \textsf{+} & \textsf{--} & \textsf{--} & \textsf{--} & -- & G-tst/c \\
BKL.fl & induc & 2 & 3 & \texttt{SSSuu} & \textsf{--} & \textsf{+} & \textsf{--} & \textsf{+} & \textsf{+} & -- & K \\
BKL.sec & induc & 2 & $\infty$ & \texttt{SSSSS} & \textsf{+} & \textsf{+} & \textsf{+} & \textsf{+} & \textsf{+} & -- & K \\
BKL.fal & induc & 2 & 0 & \texttt{uuuuu} & \textsf{--} & \textsf{--} & \textsf{--} & \textsf{--} & \textsf{--} & -- & F \\
BKL.unj & induc & 2 & 3 & \texttt{SSSuu} & \textsf{--} & \textsf{--} & \textsf{--} & \textsf{--} & \textsf{--} & -- & U \\
BKL.rel & testi & 2 & 3 & \texttt{SSSuu} & \textsf{+} & \textsf{+} & \textsf{--} & \textsf{+} & \textsf{+} & -- & K \\
BKL.lie & testi & 2 & 1 & \texttt{Suuuu} & \textsf{+} & \textsf{+} & \textsf{--} & \textsf{--} & \textsf{--} & -- & G-tst/c \\
BKH.fl & induc & 3 & 3 & \texttt{SSSuu} & \textsf{--} & \textsf{+} & \textsf{--} & \textsf{+} & \textsf{--} & -- & G-fl/c \\
BKH.sec & induc & 3 & $\infty$ & \texttt{SSSSS} & \textsf{+} & \textsf{+} & \textsf{+} & \textsf{+} & \textsf{+} & -- & K \\
BKH.rem & induc & 3 & 4 & \texttt{SSSSu} & \textsf{+} & \textsf{+} & \textsf{--} & \textsf{+} & \textsf{+} & -- & K \\
BKH.fal & induc & 3 & 0 & \texttt{uuuuu} & \textsf{--} & \textsf{--} & \textsf{--} & \textsf{--} & \textsf{--} & -- & F \\
BKH.unj & induc & 3 & 3 & \texttt{SSSuu} & \textsf{--} & \textsf{--} & \textsf{--} & \textsf{--} & \textsf{--} & -- & U \\
BKH.rel & testi & 3 & 3 & \texttt{SSSuu} & \textsf{+} & \textsf{+} & \textsf{--} & \textsf{+} & \textsf{--} & -- & G-int/c \\
BKH.lie & testi & 3 & 1 & \texttt{Suuuu} & \textsf{+} & \textsf{+} & \textsf{--} & \textsf{--} & \textsf{--} & -- & G-tst/c \\
BIV.nor & perce & 2 & $\infty$ & \texttt{SSSSS} & \textsf{+} & \textsf{+} & \textsf{+} & \textsf{+} & \textsf{+} & -- & K \\
BIV.hos & perce & 2 & 1 & \texttt{Suuuu} & \textsf{+} & \textsf{+} & \textsf{--} & \textsf{+} & \textsf{--} & -- & G-env/c \\
BIV.fal & perce & 2 & 0 & \texttt{uuuuu} & \textsf{--} & \textsf{--} & \textsf{--} & \textsf{--} & \textsf{--} & -- & F \\
BIV.unj & perce & 2 & 4 & \texttt{SSSSu} & \textsf{--} & \textsf{--} & \textsf{--} & \textsf{--} & \textsf{--} & -- & U \\
BIV.rel & testi & 2 & 4 & \texttt{SSSSu} & \textsf{+} & \textsf{+} & \textsf{--} & \textsf{+} & \textsf{+} & -- & K \\
BIV.lie & testi & 2 & 1 & \texttt{Suuuu} & \textsf{+} & \textsf{+} & \textsf{--} & \textsf{--} & \textsf{--} & -- & G-tst/c \\
JJM.nor & perce & 2 & $\infty$ & \texttt{SSSSS} & \textsf{+} & \textsf{+} & \textsf{+} & \textsf{+} & \textsf{+} & -- & K[ev] \\
JJM.hos & perce & 2 & 1 & \texttt{Suuuu} & \textsf{+} & \textsf{+} & \textsf{--} & \textsf{+} & \textsf{--} & -- & G-env/c \\
JJM.fal & perce & 2 & 0 & \texttt{uuuuu} & \textsf{--} & \textsf{--} & \textsf{--} & \textsf{--} & \textsf{--} & -- & F \\
JJM.unj & perce & 2 & 4 & \texttt{SSSSu} & \textsf{--} & \textsf{--} & \textsf{--} & \textsf{--} & \textsf{--} & -- & U \\
JJM.rel & testi & 2 & 4 & \texttt{SSSSu} & \textsf{+} & \textsf{+} & \textsf{--} & \textsf{+} & \textsf{+} & -- & K[ev] \\
JJM.lie & testi & 2 & 1 & \texttt{Suuuu} & \textsf{+} & \textsf{+} & \textsf{--} & \textsf{--} & \textsf{--} & -- & G-tst/c \\
THT.sec & testi & 2 & $\infty$ & \texttt{SSSSS} & \textsf{--} & \textsf{--} & \textsf{--} & \textsf{--} & \textsf{--} & -- & U \\
THT.hos & testi & 2 & 1 & \texttt{Suuuu} & \textsf{--} & \textsf{--} & \textsf{--} & \textsf{--} & \textsf{--} & -- & U \\
THT.rem & testi & 2 & 4 & \texttt{SSSSu} & \textsf{--} & \textsf{--} & \textsf{--} & \textsf{--} & \textsf{--} & -- & U \\
THT.fal & testi & 2 & 0 & \texttt{uuuuu} & \textsf{--} & \textsf{--} & \textsf{--} & \textsf{--} & \textsf{--} & -- & U \\
THT.rel & testi & 2 & 1 & \texttt{Suuuu} & \textsf{--} & \textsf{--} & \textsf{--} & \textsf{--} & \textsf{--} & -- & U \\
THT.lie & testi & 2 & 1 & \texttt{Suuuu} & \textsf{--} & \textsf{--} & \textsf{--} & \textsf{--} & \textsf{--} & -- & U \\
ICE.fl & induc & 2 & 4 & \texttt{SSSSu} & \textsf{--} & \textsf{+} & \textsf{--} & \textsf{+} & \textsf{+} & -- & K \\
ICE.sec & induc & 2 & $\infty$ & \texttt{SSSSS} & \textsf{+} & \textsf{+} & \textsf{+} & \textsf{+} & \textsf{+} & -- & K \\
ICE.fal & induc & 2 & 0 & \texttt{uuuuu} & \textsf{--} & \textsf{--} & \textsf{--} & \textsf{--} & \textsf{--} & -- & F \\
ICE.unj & induc & 2 & 4 & \texttt{SSSSu} & \textsf{--} & \textsf{--} & \textsf{--} & \textsf{--} & \textsf{--} & -- & U \\
ICE.rel & testi & 2 & 4 & \texttt{SSSSu} & \textsf{+} & \textsf{+} & \textsf{--} & \textsf{+} & \textsf{+} & -- & K \\
ICE.lie & testi & 2 & 1 & \texttt{Suuuu} & \textsf{+} & \textsf{+} & \textsf{--} & \textsf{--} & \textsf{--} & -- & G-tst/c \\
SNC.rem & calcu & 2 & 4 & \texttt{SSSSu} & \textsf{+} & \textsf{--} & \textsf{+} & \textsf{+} & \textsf{+} & -- & K \\
SNC.fal & calcu & 2 & 0 & \texttt{uuuuu} & \textsf{--} & \textsf{--} & \textsf{--} & \textsf{--} & \textsf{--} & -- & F \\
SNC.unj & calcu & 2 & 1 & \texttt{Suuuu} & \textsf{--} & \textsf{--} & \textsf{--} & \textsf{--} & \textsf{--} & -- & U \\
SNC.rel & testi & 2 & 1 & \texttt{Suuuu} & \textsf{+} & \textsf{--} & \textsf{+} & \textsf{+} & \textsf{--} & -- & G-nec/d \\
SNC.lie & testi & 2 & 1 & \texttt{Suuuu} & \textsf{+} & \textsf{--} & \textsf{--} & \textsf{--} & \textsf{--} & -- & G-tst/d \\
SNS.fal & calcu & 2 & 0 & \texttt{uuuuu} & \textsf{--} & \textsf{--} & \textsf{--} & \textsf{--} & \textsf{--} & -- & F \\
SNS.unj & calcu & 2 & $\infty$ & \texttt{SSSSS} & \textsf{--} & \textsf{--} & \textsf{--} & \textsf{--} & \textsf{--} & -- & U \\
SNS.rel & testi & 2 & $\infty$ & \texttt{SSSSS} & \textsf{+} & \textsf{--} & \textsf{+} & \textsf{+} & \textsf{+} & -- & K \\
SNS.lie & testi & 2 & 1 & \texttt{Suuuu} & \textsf{+} & \textsf{--} & \textsf{--} & \textsf{--} & \textsf{--} & -- & G-tst/d \\
SEV.nor & perce & 2 & $\infty$ & \texttt{SSSSS} & \textsf{+} & \textsf{+} & \textsf{+} & \textsf{+} & \textsf{+} & -- & K[ev] \\
SEV.hos & perce & 2 & 1 & \texttt{Suuuu} & \textsf{+} & \textsf{+} & \textsf{--} & \textsf{+} & \textsf{--} & -- & G-env/c \\
SEV.fal & perce & 2 & 0 & \texttt{uuuuu} & \textsf{--} & \textsf{--} & \textsf{--} & \textsf{--} & \textsf{--} & -- & F \\
SEV.unj & perce & 2 & 4 & \texttt{SSSSu} & \textsf{--} & \textsf{--} & \textsf{--} & \textsf{--} & \textsf{--} & -- & U \\
SEV.rel & testi & 2 & 4 & \texttt{SSSSu} & \textsf{+} & \textsf{+} & \textsf{--} & \textsf{+} & \textsf{+} & -- & K[ev] \\
SEV.lie & testi & 2 & 1 & \texttt{Suuuu} & \textsf{+} & \textsf{+} & \textsf{--} & \textsf{--} & \textsf{--} & -- & G-tst/c \\
SCT.nor & testi & 2 & $\infty$ & \texttt{SSSSS} & \textsf{+} & \textsf{+} & \textsf{+} & \textsf{+} & \textsf{+} & -- & K[ct] \\
SCT.hos & testi & 2 & 1 & \texttt{Suuuu} & \textsf{+} & \textsf{+} & \textsf{--} & \textsf{+} & \textsf{--} & -- & G-env/c \\
SCT.fal & testi & 2 & 0 & \texttt{uuuuu} & \textsf{--} & \textsf{--} & \textsf{--} & \textsf{--} & \textsf{--} & -- & F \\
SCT.unj & testi & 2 & 4 & \texttt{SSSSu} & \textsf{--} & \textsf{--} & \textsf{--} & \textsf{--} & \textsf{--} & -- & U \\
SCT.rel & testi & 2 & 4 & \texttt{SSSSu} & \textsf{+} & \textsf{+} & \textsf{+} & \textsf{+} & \textsf{+} & -- & K[ct] \\
SCT.lie & testi & 2 & 1 & \texttt{Suuuu} & \textsf{+} & \textsf{+} & \textsf{--} & \textsf{--} & \textsf{--} & -- & G-tst/c \\
SRF.hos & perce & 2 & 1 & \texttt{Suuuu} & \textsf{+} & \textsf{+} & \textsf{--} & \textsf{+} & \textsf{--} & -- & G-env/c \\
SRF.fal & perce & 2 & 0 & \texttt{uuuuu} & \textsf{--} & \textsf{--} & \textsf{--} & \textsf{--} & \textsf{--} & -- & F \\
SRF.unj & perce & 2 & $\infty$ & \texttt{SSSSS} & \textsf{--} & \textsf{--} & \textsf{--} & \textsf{--} & \textsf{--} & -- & U \\
SRF.rel & testi & 2 & $\infty$ & \texttt{SSSSS} & \textsf{+} & \textsf{+} & \textsf{+} & \textsf{+} & \textsf{+} & -- & K[rf] \\
SRF.lie & testi & 2 & 1 & \texttt{Suuuu} & \textsf{+} & \textsf{+} & \textsf{--} & \textsf{--} & \textsf{--} & -- & G-tst/c \\
STL.nor & testi & 2 & $\infty$ & \texttt{SSSSS} & \textsf{+} & \textsf{--} & \textsf{+} & \textsf{--} & \textsf{+} & -- & K[ev] \\
STL.rem & testi & 2 & 4 & \texttt{SSSSu} & \textsf{+} & \textsf{--} & \textsf{--} & \textsf{--} & \textsf{+} & -- & K[ev] \\
STL.fal & testi & 2 & 0 & \texttt{uuuuu} & \textsf{--} & \textsf{--} & \textsf{--} & \textsf{--} & \textsf{--} & -- & F \\
STL.unj & testi & 2 & 1 & \texttt{Suuuu} & \textsf{--} & \textsf{--} & \textsf{--} & \textsf{--} & \textsf{--} & -- & U \\
STL.rel & testi & 2 & 1 & \texttt{Suuuu} & \textsf{+} & \textsf{--} & \textsf{--} & \textsf{--} & \textsf{--} & -- & G-tst/d \\
STL.lie & testi & 2 & 1 & \texttt{Suuuu} & \textsf{+} & \textsf{--} & \textsf{--} & \textsf{--} & \textsf{--} & -- & G-tst/d \\
\end{longtable}}


\paragraph{Classes across thresholds.} The second table gives, for each case of Appendix~A, its class at each threshold $\theta = 0, \ldots, 4$, with the threshold at which the case is assessed. At $\theta = 0$ only the actual world is close, and every justified true belief is classed as knowledge; this is the limiting context in which no error possibility is attended to. A Gettier case of the literature enters its class at $\theta = \delta$ and keeps it at every higher threshold, by Proposition~\ref{p:margin}(a). The rows of RED, LTN and the bank cases show knowledge that is lost only at thresholds at which a remote possibility, a red facade, a misprint, a change of hours, has been made close.

{\small\setlength{\tabcolsep}{4pt}
\begin{longtable}{@{}lccccccl@{}}
\toprule
Id & $\delta$ & $\theta=0$ & $1$ & $2$ & $3$ & $4$ & Case threshold \\
\midrule\endhead
\bottomrule\endfoot
GE1 & 1 & K & G-fl/d & G-fl/d & G-fl/d & G-fl/d & 2 \\
GE2 & 1 & K & G-fl/d & G-fl/d & G-fl/d & G-fl/d & 2 \\
NOG & 1 & K & G-fl/d & G-fl/d & G-fl/d & G-fl/d & 2 \\
FEL & 1 & K & G-nfl/d & G-nfl/d & G-nfl/d & G-nfl/d & 2 \\
CLK & 1 & K & G-tmp/d & G-tmp/d & G-tmp/d & G-tmp/d & 2 \\
BRN & 1 & K & G-env/c & G-env/c & G-env/c & G-env/c & 2 \\
RED & 3 & K & K & K & G-env/c & G-env/c & 2 \\
RDP & 1 & K & G-env/c & G-env/c & G-env/c & G-env/c & 2 \\
SHP & 1 & K & G-int/d & G-int/d & G-int/d & G-int/d & 2 \\
ZDR & 1 & K & G-fl/d & G-fl/d & G-fl/d & G-fl/d & 2 \\
ZHB & 1 & K & G-int/d & G-int/d & G-int/d & G-int/d & 2 \\
PYR & 1 & K & G-fl/d & G-fl/d & G-fl/d & G-fl/d & 2 \\
GRB & 4 & K[md] & K[md] & K[md] & K[md] & G-env/c & 2 \\
HAR & $\infty$ & K[md] & K[md] & K[md] & K[md] & K[md] & 2 \\
CLV & $\infty$ & U & U & U & U & U & 2 \\
TTP & $\infty$ & U & U & U & U & U & 2 \\
CHU & 4 & K & K & K & K & G-nfl/c & 2 \\
KAL & 1 & K[rf] & G-env/c & G-env/c & G-env/c & G-env/c & 2 \\
TMP & 3 & K[fx] & K[fx] & K[fx] & G-int/d & G-int/d & 2 \\
MOR & 3 & K & K & K & G-tst/c & G-tst/c & 2 \\
LOT & 1 & K & G-lot/d & G-lot/d & G-lot/d & G-lot/d & 2 \\
LTN & 3 & K & K & K & G-tst/c & G-tst/c & 2 \\
MAG & 1 & K & G-mrg/c & G-mrg/c & G-mrg/c & G-mrg/c & 2 \\
MGW & $\infty$ & K & K & K & K & K & 2 \\
ZBP & 4 & K & K & K & K & G-env/c & 2 \\
ZBM & 4 & K & K & K & K & G-env/c & 2 \\
GRN & $\infty$ & K & K & K[ev] & K[ev] & K[ev] & 2 \\
HLF & $\infty$ & K & K & K & K & K & 2 \\
HLC & 1 & K & G-tst/c & G-tst/c & G-tst/c & G-tst/c & 2 \\
ATF & $\infty$ & K & K & K & K & K & 2 \\
ATC & 1 & K & G-env/c & G-env/c & G-env/c & G-env/c & 2 \\
WAR & $\infty$ & K & K & K & K & K & 2 \\
BKL & 3 & K & K & K & G-nfl/c & G-nfl/c & 2 \\
BKH & 3 & K & K & K & G-nfl/c & G-nfl/c & 3 \\
BIV & 4 & K & K & K & K & G-env/c & 2 \\
JJM & 4 & K & K[ev] & K[ev] & K[ev] & G-env/c & 2 \\
THT & 1 & U & U & U & U & U & 2 \\
ICE & 4 & K & K & K & K & G-nfl/c & 2 \\
SNC & 1 & K & G-nec/d & G-nec/d & G-nec/d & G-nec/d & 2 \\
SNS & $\infty$ & K & K & K & K & K & 2 \\
SEV & 4 & K & K[ev] & K[ev] & K[ev] & G-env/c & 2 \\
SCT & 4 & K[ct] & K[ct] & K[ct] & K[ct] & G-tst/c & 2 \\
SRF & $\infty$ & K[rf] & K[rf] & K[rf] & K[rf] & K[rf] & 2 \\
STL & 1 & K & G-tst/d & G-tst/d & G-tst/d & G-tst/d & 2 \\
\end{longtable}}


\section{The typology: incidence and check}\label{app:d}

\paragraph{Classes.} Table~\ref{tab:classes} gives, for each class of Theorem~\ref{t:partition}, the number of literature and schematic cases and the number of all models that fall under it, with the first witnesses in the order of the catalogue. Every class is realized by at least one source case. G-nfl, G-int, G-tmp, G-lot, G-tst, G-mrg and G-nec are realized by one or two source cases each, the remainder of their members being variants; G-tst, the largest Gettier class among all models, owes its size to the relay variants \texttt{lie}, one for every case whose belief is true.

\begin{table}[h]\caption{Classes of the typology.}\label{tab:classes}
\begin{center}\small
\begin{tabular}{@{}lp{0.36\linewidth}rrp{0.3\linewidth}@{}}\toprule
Class & Description & Sources & All & Witnesses (first five) \\ \midrule
G-fl & inferential, through a false lemma; truth-maker varies & 5 & 7 & GE1, GE2, NOG, ZDR, PYR \\
G-nfl & inferential, no false lemma; truth-maker varies & 2 & 7 & FEL, BKH, GE1.nfl, GE2.nfl, NOG.nfl \\
G-int & non-inferential; truth-maker varies (intervening) & 2 & 13 & SHP, ZHB, GE1.rel, GE2.rel, NOG.rel \\
G-env & look-alike, trick or faulty device nearby & 4 & 26 & BRN, RDP, KAL, ATC, BRN.rel \\
G-tmp & moment of consultation varies (temporal) & 1 & 5 & CLK, CLK.nfl, CLK.hos, CLK.rel, CLK.lie \\
G-lot & outcome of a chance process varies (lottery) & 1 & 3 & LOT, LOT.rel, LOT.lie \\
G-tst & testimonial source varies or is insincere & 2 & 42 & HLC, STL, GE1.lie, GE2.lie, NOG.lie \\
G-mrg & parameter within the margin of discrimination & 1 & 4 & MAG, MAG.hos, MAG.rel, MAG.lie \\
G-nec & necessary content; basis-wide unsafety only & 1 & 2 & SNC, SNC.rel \\
K & knowledge (justified, basis-wide safe) & 22 & 105 & RED, GRB, HAR, CHU, TMP \\
U & unjustified & 3 & 56 & CLV, TTP, THT, GE1.unj, GE2.unj \\
F & justified false belief & 0 & 41 & GE1.fal, GE2.fal, NOG.fal, FEL.fal, CLK.fal \\
\bottomrule\end{tabular}\end{center}

\end{table}

\paragraph{Classes by route.} Table~\ref{tab:matrix} gives the number of Gettier models of each class by route of the basis. The inferential classes occur only with inferential routes, by definition; G-env occurs with perception, instruments and testimony, and the testimonial column is filled in every class reachable by relay (Proposition~\ref{p:testimony}).

\begin{table}[h]\caption{Gettier models by class and route.}\label{tab:matrix}
\begin{center}\small
\begin{tabular}{@{}lrrrrrrrrr@{}}\toprule
Class & infer. & induc. & perce. & instr. & testi. & stati. & deduc. & calcu. & facul. \\ \midrule
G-fl & 5 & 2 & . & . & . & . & . & . & . \\
G-nfl & 5 & 2 & . & . & . & . & . & . & . \\
G-int & . & . & 4 & . & 9 & . & . & . & . \\
G-env & . & . & 14 & 3 & 9 & . & . & . & . \\
G-tmp & . & . & . & 3 & 2 & . & . & . & . \\
G-lot & . & . & . & . & 2 & 1 & . & . & . \\
G-tst & . & . & . & . & 42 & . & . & . & . \\
G-mrg & . & . & 2 & . & 2 & . & . & . & . \\
G-nec & . & . & . & . & 1 & . & . & 1 & . \\
\bottomrule\end{tabular}\end{center}

\end{table}

\paragraph{Causal connection and false lemmas.} Table~\ref{tab:combos} gives, among the justified true models, the number with and without a causal connection, and with and without a false lemma, that are safe and unsafe. All eight cells are non-empty, which verifies Propositions~\ref{p:lemma} and \ref{p:causal}.

\begin{table}[h]\caption{Causal connection, false lemmas and safety among justified true models.}\label{tab:combos}
\begin{center}\small
\begin{tabular}{@{}llrp{0.48\linewidth}@{}}\toprule
Feature & Safe & Count & Witnesses (first six) \\ \midrule
causal yes & yes & 85 & RED, GRB, HAR, CHU, MOR, LTN \\
causal yes & no & 58 & BRN, RDP, KAL, MAG, HLC, ATC \\
causal no & yes & 20 & TMP, SNS, GE1.rem, GE1.ded, GE2.rem, NOG.rem \\
causal no & no & 51 & GE1, GE2, NOG, FEL, CLK, SHP \\
\midrule
false lemma yes & yes & 24 & WAR, GE1.sec, GE1.rem, GE1.ded, GE2.sec, GE2.rem \\
false lemma yes & no & 30 & GE1, GE2, NOG, CLK, SHP, ZDR \\
false lemma no & yes & 81 & RED, GRB, HAR, CHU, TMP, MOR \\
false lemma no & no & 79 & FEL, BRN, RDP, KAL, LOT, MAG \\
\bottomrule\end{tabular}\end{center}

\end{table}

\paragraph{Rival conditions.} Table~\ref{tab:rivals} gives, for each condition of Appendix~C, the number of the \stnundisputed{} undisputed literature cases on which it delivers the source's verdict, and the number of variants on which it agrees with safety.

\begin{table}[h]\caption{Rival conditions against the literature and against safety.}\label{tab:rivals}
\begin{center}\small
\begin{tabular}{@{}lrrr@{}}\toprule
Condition & Correct on undisputed sources & Rate & Agrees with safety on variants \\ \midrule
NFL & 23 / 29 & 0.79 & 176 / 267 \\
CAU & 26 / 29 & 0.90 & 198 / 267 \\
SEN & 21 / 29 & 0.72 & 232 / 267 \\
UND & 25 / 29 & 0.86 & 210 / 267 \\
SAF & 29 / 29 & 1.00 & -- \\
\bottomrule\end{tabular}\end{center}

\end{table}

\paragraph{Departures.} Table~\ref{tab:departures} lists the literature cases on which the computed verdict differs from the source's, with the kind of dispute and its note. Every departure carries a dispute record of kind \texttt{dep} or \texttt{bas}; the program asserts this.

\begin{table}[h]\caption{Departures from the sources' verdicts.}\label{tab:departures}
\begin{center}\small
\begin{tabular}{@{}lllp{0.5\linewidth}@{}}\toprule
Id & Source verdict & Computed & Kind and note \\ \midrule
HAR & not knowledge & K[md] & (dep) Harman's verdict is that Jill does not know. The case has the structure of GRB, where the source's verdict is knowledge: a misleading defeater the subject does not possess. \\
KAL & knowledge & G-env/c & (dep) Sosa grants animal knowledge (apt belief) and denies reflective knowledge; the belief is unsafe, so on the analysis it is not knowledge. \\
TMP & not knowledge & K[fx] & (dep) Pritchard's verdict is that Temp does not know. The belief is justified and safe; the truth is secured by another agent, as in MOR, where the source's verdict is knowledge. \\
ZBM & not knowledge & K & (dep) Dretske denies knowledge, since the belief is insensitive; closure is defended by \citet{Williamson2000} and \citet{Hawthorne2004}, and by EPI-20. \\
HLC & knowledge & G-tst/c & (bas) Basis individuated coarsely: a Gettier case, against Comesa\~na's verdict. \\
ATC & knowledge & G-env/c & (bas) Basis individuated coarsely: a Gettier case, against Bogardus's verdict. \\
\bottomrule\end{tabular}\end{center}

\end{table}

\section{Exhaustive check on small models}\label{app:e}

\paragraph{What is enumerated.} The propositions of Section~\ref{s:thm} are stated for all case models, and the catalogue tests them only on the models it contains. To test them independently of the encodings, the program enumerates every case model in a bounded space: at most two non-actual worlds; each world at a distance from $1$ to $4$, with one of the four combinations of belief and truth values, and, if it is an error world, one of the six error tags (otherwise the tag \texttt{ok}, or \texttt{basis} when the subject lacks the basis and $p$ is true); justification and actual truth each $0$ or $1$; route inference (with or without a false lemma), perception or testimony; and causal connection present or absent. There are \stcnconfigs{} configurations of worlds (ordered, so that each unordered configuration with two distinct worlds occurs twice) and \stcnmodels{} models. The value alt is not varied, so the class G-nec does not occur in the census; Proposition~\ref{p:nec} is checked on the catalogue.

\paragraph{What is checked.} For every model and every threshold $\theta = 0, \ldots, 4$, the program checks Lemma~\ref{l:delta}, Theorem~\ref{t:partition} (the class is unique, and it is a Gettier class exactly when the belief is justified, true and unsafe), Proposition~\ref{p:veritic}, the monotonicity of Proposition~\ref{p:margin}(a), and the characterization of EPI-01-a (Theorem~3.3 and Corollary~3.4) in the form \enquote{a justified true belief is knowledge if and only if it is not a Gettier case}; for inferential models it checks Proposition~\ref{p:lemma} by toggling the lemma, and for perceptual and testimonial models Proposition~\ref{p:env}(c) by adding a look-alike. In all, \stcnchecks{} model checks are performed and every one succeeds. The enumeration is exhaustive within the bound because it ranges over the product of the finite sets just listed. The table gives the distribution of the classes at the default threshold.

\begin{center}\small
\begin{tabular}{@{}lr@{}}\toprule
Class at $\theta = 2$ & Models \\ \midrule
G-fl & 224 \\
G-nfl & 224 \\
G-int & 448 \\
G-env & 928 \\
G-tmp & 1056 \\
G-lot & 1024 \\
G-tst & 960 \\
G-mrg & 992 \\
G-nec & 0 \\
K & 4808 \\
U & 21328 \\
F & 10664 \\
\midrule Total & 42656 \\
\bottomrule\end{tabular}\end{center}


\begin{thebibliography}{99}\small
\bibitem[BonJour(1980)]{BonJour1980} BonJour, L. (1980). Externalist theories of empirical knowledge. \emph{Midwest Studies in Philosophy}, 5, 53--73.
\bibitem[Bogardus(2014)]{Bogardus2014} Bogardus, T. (2014). Knowledge under threat. \emph{Philosophy and Phenomenological Research}, 88(2), 289--313.
\bibitem[Chisholm(1966)]{Chisholm1966} Chisholm, R.~M. (1966). \emph{Theory of Knowledge}. Prentice-Hall.
\bibitem[Clark(1963)]{Clark1963} Clark, M. (1963). Knowledge and grounds: A comment on Mr. Gettier's paper. \emph{Analysis}, 24(2), 46--48.
\bibitem[Cola{\c c}o et~al.(2014)]{Colaco2014} Cola{\c c}o, D., Buckwalter, W., Stich, S. and Machery, E. (2014). Epistemic intuitions in fake-barn thought experiments. \emph{Episteme}, 11(2), 199--212.
\bibitem[Comesa{\~n}a(2005)]{Comesana2005} Comesa{\~n}a, J. (2005). Unsafe knowledge. \emph{Synthese}, 146(3), 395--404.
\bibitem[DeRose(1992)]{DeRose1992} DeRose, K. (1992). Contextualism and knowledge attributions. \emph{Philosophy and Phenomenological Research}, 52(4), 913--929.
\bibitem[Dretske(1970)]{Dretske1970} Dretske, F. (1970). Epistemic operators. \emph{The Journal of Philosophy}, 67(24), 1007--1023.
\bibitem[Feldman(1974)]{Feldman1974} Feldman, R. (1974). An alleged defect in Gettier counter-examples. \emph{Australasian Journal of Philosophy}, 52(1), 68--69.
\bibitem[Gendler and Hawthorne(2005)]{GendlerHawthorne2005} Gendler, T.~S. and Hawthorne, J. (2005). The real guide to fake barns: A catalogue of gifts for your epistemic enemies. \emph{Philosophical Studies}, 124(3), 331--352.
\bibitem[Gettier(1963)]{Gettier1963} Gettier, E.~L. (1963). Is justified true belief knowledge? \emph{Analysis}, 23(6), 121--123.
\bibitem[Goldman(1967)]{Goldman1967} Goldman, A.~I. (1967). A causal theory of knowing. \emph{The Journal of Philosophy}, 64(12), 357--372.
\bibitem[Goldman(1976)]{Goldman1976} Goldman, A.~I. (1976). Discrimination and perceptual knowledge. \emph{The Journal of Philosophy}, 73(20), 771--791.
\bibitem[Harman(1973)]{Harman1973} Harman, G. (1973). \emph{Thought}. Princeton University Press.
\bibitem[Hawthorne(2004)]{Hawthorne2004} Hawthorne, J. (2004). \emph{Knowledge and Lotteries}. Oxford University Press.
\bibitem[Klein(2008)]{Klein2008} Klein, P. (2008). Useful false beliefs. In Q. Smith (Ed.), \emph{Epistemology: New Essays}. Oxford University Press.
\bibitem[Kripke(2011)]{Kripke2011} Kripke, S.~A. (2011). Nozick on knowledge. In \emph{Philosophical Troubles: Collected Papers, Volume 1} (pp.~162--224). Oxford University Press.
\bibitem[Lackey(2007)]{Lackey2007} Lackey, J. (2007). Why we don't deserve credit for everything we know. \emph{Synthese}, 158(3), 345--361.
\bibitem[Lehrer(1965)]{Lehrer1965} Lehrer, K. (1965). Knowledge, truth and evidence. \emph{Analysis}, 25(5), 168--175.
\bibitem[Lehrer(1990)]{Lehrer1990} Lehrer, K. (1990). \emph{Theory of Knowledge}. Westview Press.
\bibitem[Lehrer and Paxson(1969)]{LehrerPaxson1969} Lehrer, K. and Paxson, T. (1969). Knowledge: Undefeated justified true belief. \emph{The Journal of Philosophy}, 66(8), 225--237.
\bibitem[Machery et~al.(2017)]{Machery2017} Machery, E., Stich, S., Rose, D., Chatterjee, A., Karasawa, K., Struchiner, N., Sirker, S., Usui, N. and Hashimoto, T. (2017). Gettier across cultures. \emph{No\^us}, 51(3), 645--664.
\bibitem[Nozick(1981)]{Nozick1981} Nozick, R. (1981). \emph{Philosophical Explanations}. Harvard University Press.
\bibitem[Plato(1997)]{PlatoTht} Plato (1997). \emph{Theaetetus}. Trans. M.~J. Levett, rev. M. Burnyeat. In J.~M. Cooper (Ed.), \emph{Plato: Complete Works}. Hackett.
\bibitem[Pritchard(2005)]{Pritchard2005} Pritchard, D. (2005). \emph{Epistemic Luck}. Oxford University Press.
\bibitem[Pritchard(2012)]{Pritchard2012} Pritchard, D. (2012). Anti-luck virtue epistemology. \emph{The Journal of Philosophy}, 109(3), 247--279.
\bibitem[Russell(1948)]{Russell1948} Russell, B. (1948). \emph{Human Knowledge: Its Scope and Limits}. George Allen \& Unwin.
\bibitem[Shope(1983)]{Shope1983} Shope, R.~K. (1983). \emph{The Analysis of Knowing: A Decade of Research}. Princeton University Press.
\bibitem[Skyrms(1967)]{Skyrms1967} Skyrms, B. (1967). The explication of \enquote{X knows that p}. \emph{The Journal of Philosophy}, 64(12), 373--389.
\bibitem[Sosa(1999)]{Sosa1999} Sosa, E. (1999). How to defeat opposition to Moore. \emph{Philosophical Perspectives}, 13, 141--153.
\bibitem[Sosa(2007)]{Sosa2007} Sosa, E. (2007). \emph{A Virtue Epistemology: Apt Belief and Reflective Knowledge, Volume I}. Oxford University Press.
\bibitem[Starmans and Friedman(2012)]{StarmansFriedman2012} Starmans, C. and Friedman, O. (2012). The folk conception of knowledge. \emph{Cognition}, 124(3), 272--283.
\bibitem[Unger(1968)]{Unger1968} Unger, P. (1968). An analysis of factual knowledge. \emph{The Journal of Philosophy}, 65(6), 157--170.
\bibitem[Vogel(1987)]{Vogel1987} Vogel, J. (1987). Tracking, closure, and inductive knowledge. In S. Luper-Foy (Ed.), \emph{The Possibility of Knowledge: Nozick and His Critics}. Rowman \& Littlefield.
\bibitem[Warfield(2005)]{Warfield2005} Warfield, T.~A. (2005). Knowledge from falsehood. \emph{Philosophical Perspectives}, 19, 405--416.
\bibitem[Weinberg et~al.(2001)]{WeinbergNicholsStich2001} Weinberg, J.~M., Nichols, S. and Stich, S. (2001). Normativity and epistemic intuitions. \emph{Philosophical Topics}, 29(1--2), 429--460.
\bibitem[Williamson(2000)]{Williamson2000} Williamson, T. (2000). \emph{Knowledge and Its Limits}. Oxford University Press.
\bibitem[Zagzebski(1994)]{Zagzebski1994} Zagzebski, L. (1994). The inescapability of Gettier problems. \emph{The Philosophical Quarterly}, 44(174), 65--73.
\end{thebibliography}

\end{document}
